\chapter[The Virasoro VOA: Verma modules, Kac determinant, screenings]{The Virasoro vertex operator algebra: Verma modules, the Kac determinant, screenings}\label{ch:L06}
\begingroup\setlength{\emergencystretch}{3em}

\providecommand{\lsixvac}{\mathbf{1}}
\providecommand{\lsixVir}{\mathfrak{Vir}}
\providecommand{\lsixU}{U(\mathfrak{Vir})}
\providecommand{\lsixV}{V(c,0)}
\providecommand{\lsixL}{L(c,0)}
\providecommand{\lsixM}{M(c,h)}
\providecommand{\lsixK}{K}
\providecommand{\lsixhrs}{h_{r,s}}
\providecommand{\lsixct}{c(t)}
\providecommand{\lsixpp}{\mathrm{p}}
\providecommand{\lsixpt}{\mathrm{p}_{2}}
\providecommand{\lsixG}{G}
\providecommand{\lsixform}[2]{\langle #1,#2\rangle}

This chapter documents the release of the LeanVOA library named in the front matter. Every Lean quotation below is checked by \texttt{scripts/textbook/sync.py}
to occur verbatim (up to whitespace and comments) in the source files of that release. The chapter was
first written against an earlier state of the library; checks that were run at that earlier state and not re-run for this edition are marked as such. The chapter covers the 117 files
\leanpath{VOA/\allowbreak Examples/\allowbreak Virasoro*.lean} that \texttt{Notes/Textbook/assignment.txt} assigns to it
(the Virasoro files about rigidity, Benoit--Saint-Aubin vectors, fusion and the boundary models belong
to Chapter~\ref{ch:L07} and Chapter~\ref{ch:L19}). A small number of definitions that these files
use live in files of other chapters --- most importantly the general Verma module $M(c,h)$, which is
defined in \leanpath{VOA/\allowbreak Examples/\allowbreak InverseDSVirasoroVerma.lean} (a file of
Chapter~\ref{ch:L19}); we quote them here, with their file, because the statements of this chapter
cannot be read without them.

\paragraph{What the chapter establishes.} In one paragraph, and with the honesty that the rest of the
chapter will justify in detail:
\begin{enumerate}
\item For \emph{every} complex number $c$ the library constructs the universal Virasoro vertex
  operator algebra $\lsixV$ on the universal vacuum module $U(\lsixVir)/J(c)$
  (\lean{VirModule.virasoroVOA}), proves that its creation monomials indexed by partitions into
  parts $\ge2$ form a basis (so $\dim\lsixV_n$ is the number of such partitions), and constructs its
  simple quotient $\lsixL$ (\lean{VirSimple.}\allowbreak\lean{virasoroSimple}), proved simple for every $c$
  (\lean{VirSimple.}\allowbreak\lean{virasoroSimple_}\allowbreak\lean{isSimple}). Unconditional; \lean{VOA.pbw} (the PBW theorem of Chapter~\ref{ch:L02}; formerly a project axiom, now a theorem) is
  \emph{not} used.
\item For every $(c,h)\in\C^2$ it constructs the Verma module $M(c,h)$, its PBW basis, its
  contravariant (Shapovalov) form, the Gram matrices of that form at every level, and the irreducible
  quotient $L(c,h)$ with $\dim L(c,h)_n = \operatorname{rank}G_n(c,h)$. Unconditional.
\item \textbf{The Kac determinant formula} for Virasoro Verma modules, at every level and every
  $(c,h)$, is proved --- but it is \emph{not} stated as a single library declaration. The library
  states a conditional version (\lean{vermaWeightDeterminant_}\allowbreak\lean{eq_}\allowbreak\lean{kac_}\allowbreak\lean{of_}\allowbreak\lean{singular}, which assumes the
  existence of singular vectors), an extension theorem (\lean{verma_}\allowbreak\lean{kac_}\allowbreak\lean{all_}\allowbreak\lean{of_}\allowbreak\lean{nonreal_}\allowbreak\lean{singular}), and,
  in a separate file, an unconditional construction of the singular vectors from screening operators
  (\lean{screening_}\allowbreak\lean{nonreal_}\allowbreak\lean{kac_}\allowbreak\lean{singular}). The composition is two lines; we performed it in Lean against the
  library (the check is not part of this release), and it
  compiled with only the three standard axioms (Section~\ref{l06:sec:kacunconditional}). Several
  docstrings in the Kac files still say that the singular-vector premise ``remains open''; they were
  written before the screening files landed and are outdated (Warning~\ref{l06:warn:docstrings}).
\item For the \emph{vacuum} module the library proves the explicit Gram determinants at levels
  $4,5,6$ (the Kac determinant of $\lsixV$ there), the general degree and leading coefficient of the
  vacuum Gram determinant, and the fact that it has a \emph{simple} zero at every minimal-model central
  charge $c_{p,q}$ at the first singular level $(p-1)(q-1)$; it does \emph{not} state the full
  vacuum determinant formula at all levels.
\item It computes graded dimensions of $\lsixL$ at Lee--Yang ($c = -22/5$: $\dim = 1,1,2$ at levels
  $4,5,6$), Ising ($c = 1/2$: $\dim L_6 = 3$) and $M(2,7)$ ($c = -68/7$: $\dim L_6 = 3$). It does
  \emph{not} prove the full character formulas (Rogers--Ramanujan, the Ising vacuum character) beyond
  these levels.
\end{enumerate}

\paragraph{How to read the chapter.} Section~\ref{l06:sec:physics} gives the physical picture: the
stress tensor, the central charge, Verma modules as conformal families, null vectors, the Kac table,
the minimal models and the Coulomb gas. Sections~\ref{l06:sec:modes}--\ref{l06:sec:simple} build the
universal and the simple Virasoro vertex operator algebras. Section~\ref{l06:sec:gram} treats the
invariant form on the vacuum module and its Gram matrices, and Section~\ref{l06:sec:dims} the graded
dimensions of the minimal models that the library computes. Sections~\ref{l06:sec:verma}
and~\ref{l06:sec:shap} treat general Verma modules and their contravariant forms;
Section~\ref{l06:sec:kac} the Kac determinant; Section~\ref{l06:sec:screening} the screening
construction of singular vectors; Section~\ref{l06:sec:vacuumkac} the passage from Verma modules back
to the vacuum module, which is the input of the graded (first-order) rigidity theorem of Chapter~\ref{ch:L07}.
Section~\ref{l06:sec:others} lists the remaining files. The chapter closes with the warnings and dead
ends, the verification record, the referee's checklist and the references.

The Lie-algebraic objects (the Witt algebra, the Gel'fand--Fuks cocycle, the Virasoro algebra
$\lsixVir$, its universal enveloping algebra and the universal vacuum module) are defined in
Chapter~\ref{ch:L02}, Sections ``The Witt algebra and the Virasoro algebra'' and ``The universal
Virasoro vacuum module''. The notions of vertex algebra, vertex operator algebra, field, locality,
mode, ideal, simplicity, quotient and invariant bilinear form are those of Chapters~\ref{ch:L01}--\ref{ch:L04} and
Chapters~\ref{ch:L03}--\ref{ch:L04}. We recall each one at the point where it is used.

%======================================================================
\section{The physical picture}\label{l06:sec:physics}
%======================================================================

\subsection{The stress tensor and the central charge}

A two-dimensional conformal field theory has a conserved, traceless stress--energy tensor. In complex
coordinates its traceless part splits into a holomorphic field $T(z)$ and an antiholomorphic field
$\bar T(\bar z)$; the chiral algebra studied in this book is the algebra generated by holomorphic
fields such as $T(z)$. Expanding in modes,
\begin{equation}\label{l06:eq:Tmodes}
 T(z) = \sum_{n\in\Z}L_n\,z^{-n-2},
\end{equation}
the coefficients $L_n$ are the generators of infinitesimal conformal transformations
$z\mapsto z+\epsilon z^{n+1}$ acting on the space of states (radial quantisation: states live on
circles $|z| = r$ and $L_0$ generates dilatations, i.e.\ it is the Hamiltonian). The operator product
expansion of $T$ with itself is
\begin{equation}\label{l06:eq:TT}
 T(z)T(w) \sim \frac{c/2}{(z-w)^4} + \frac{2T(w)}{(z-w)^2} + \frac{\partial T(w)}{z-w},
\end{equation}
and in modes it is the Virasoro algebra
\begin{equation}\label{l06:eq:vir}
 [L_m,L_n] = (m-n)L_{m+n} + \frac{c}{12}(m^3-m)\,\delta_{m+n,0}.
\end{equation}
The number $c$ is the \emph{central charge}. Classically (for the Witt algebra of vector fields
$-z^{n+1}\partial_z$) it is absent; quantum mechanically it appears as an anomaly: the stress tensor
does not transform as a tensor under finite conformal maps but picks up $\frac{c}{12}$ times the
Schwarzian derivative; on a curved surface the trace of the stress tensor is $\propto c R$ (the
conformal or trace anomaly); and on a cylinder of circumference $L$ the ground-state energy is shifted
by $-\pi c/(6L)$ (the Casimir energy). The central charge ``counts degrees of freedom'': a free boson
has $c = 1$, a free Majorana fermion $c = 1/2$.

\subsection{Primary fields, Verma modules and descendants}

A \emph{primary} field $\phi(z)$ of conformal weight $h$ has OPE
$T(z)\phi(w)\sim h\phi(w)/(z-w)^2 + \partial\phi(w)/(z-w)$; the corresponding state
$|h\rangle = \phi(0)|0\rangle$ satisfies $L_0|h\rangle = h|h\rangle$ and $L_n|h\rangle = 0$ for $n>0$.
Such a state is called a \emph{highest-weight vector}. Acting with the creation modes $L_{-n}$, $n>0$,
produces the \emph{descendants} $L_{-n_1}\cdots L_{-n_k}|h\rangle$, of weight $h+n_1+\dots+n_k$; together
they form the \emph{conformal family} of $\phi$. The \emph{Verma module} $M(c,h)$ is the conformal family
in which no relation is imposed among descendants other than those forced by \eqref{l06:eq:vir}: by the
Poincar\'e--Birkhoff--Witt theorem a basis is given by the ordered monomials
$L_{-n_1}\cdots L_{-n_k}|h\rangle$ with $n_1\ge\dots\ge n_k\ge1$, so the number of states at level $N$
(weight $h+N$) is the number $\lsixpp(N)$ of partitions of $N$. The generating function of these
numbers is the \emph{character}
\begin{equation*}
 \chi_{M(c,h)}(q) = \operatorname{tr}_{M(c,h)}q^{L_0-c/24} = q^{h-c/24}\prod_{n\ge1}\frac{1}{1-q^n}.
\end{equation*}
The \emph{vacuum} conformal family is special: the vacuum is invariant under the global conformal group
$\mathrm{SL}_2$ generated by $L_{-1},L_0,L_1$, so $L_{-1}|0\rangle = 0$ as well. The universal vacuum
module $\lsixV$ is therefore spanned by monomials $L_{-n_1}\cdots L_{-n_k}|0\rangle$ with $n_i\ge2$, and its
character is $q^{-c/24}\prod_{n\ge2}(1-q^n)^{-1}$; the state $L_{-2}|0\rangle$ is the stress tensor
itself ($T(z)\leftrightarrow L_{-2}|0\rangle$ by the state--field correspondence).

\subsection{The Shapovalov form and null vectors}

Unitarity (or, more generally, the existence of a two-point function) is encoded in a bilinear form
with $\langle h|h\rangle = 1$ and $L_n^\dagger = L_{-n}$. Its matrix at level $N$, in the PBW basis, is the
\emph{Gram} or \emph{Shapovalov} matrix $G_N(c,h)$; for example
\begin{equation*}
 G_1 = (2h),\qquad
 G_2 = \begin{pmatrix} 4h(2h+1) & 6h\\ 6h & 4h+c/2\end{pmatrix}
 \quad\text{in the basis }(L_{-1}^2|h\rangle, L_{-2}|h\rangle).
\end{equation*}
A \emph{null vector} is a nonzero state orthogonal to every state. The basic example is a
\emph{singular vector}: a descendant $|\chi\rangle$ at level $N>0$ that is itself annihilated by all
$L_n$, $n>0$; it is automatically orthogonal to everything, because pairing it with a descendant
$L_{-n_1}\cdots|h\rangle$ moves the $L_{-n_i}$ across as $L_{n_i}$, which kill $|\chi\rangle$. Null
vectors \emph{decouple}: all their correlation functions vanish, and in a consistent theory they are
set to zero. Setting $|\chi\rangle = 0$ inside a correlation function turns into a linear differential
equation for the correlator, because $L_{-n}$ acts on a primary field at $w$ inside correlators as a
differential operator in the other insertion points --- this is the mechanism by which Belavin,
Polyakov and Zamolodchikov solved the minimal models.

\subsection{The Kac determinant and the Kac table}

Null vectors exist at level $N$ exactly when $\det G_N(c,h) = 0$. Kac's theorem computes this
determinant. Parametrise
\begin{equation}\label{l06:eq:kacparam}
 c = 13-6t-\frac6t,\qquad
 h_{r,s}(t) = \frac{(rt-s)^2-(t-1)^2}{4t}\qquad(r,s\ge1).
\end{equation}
Then, up to a nonzero constant,
\begin{equation}\label{l06:eq:kac}
 \det G_N(c,h) \;\propto\; \prod_{\substack{r,s\ge1\\ rs\le N}}\bigl(h-h_{r,s}(t)\bigr)^{\lsixpp(N-rs)}.
\end{equation}
The factor $h-h_{r,s}$ reflects a singular vector at level $rs$ in $M(c,h_{r,s})$; its exponent
$\lsixpp(N-rs)$ is the number of its descendants at level $N$. When $t = p/q$ with $p,q\ge2$ coprime,
\begin{equation*}
 c = c_{p,q} = 1-\frac{6(p-q)^2}{pq},
\end{equation*}
and the weights $h_{r,s}$, $1\le r\le q-1$, $1\le s\le p-1$ (with the identification
$h_{r,s} = h_{q-r,p-s}$) form the \emph{Kac table}: the $(p-1)(q-1)/2$ primary fields of the
\emph{minimal model} $M(p,q)$. Each $h_{r,s}$ in the table carries two independent singular vectors
and the irreducible module $L(c_{p,q},h_{r,s})$ is much smaller than the Verma module. Examples:
the Ising model $M(4,3)$, $c = 1/2$, with $h\in\{0,\frac1{16},\frac12\}$; the Lee--Yang model $M(5,2)$,
$c = -22/5$, with $h\in\{0,-\frac15\}$; the trivial model $M(3,2)$, $c = 0$.

For the vacuum $h = 0$: the label $(1,1)$ always gives $h_{1,1} = 0$, the singular vector being
$L_{-1}|0\rangle$, which the vacuum module already removes. At $c = c_{p,q}$ the next zero-weight label
is $(q-1,p-1)$ (for $p>q$), giving a vacuum null vector at level $(p-1)(q-1)$: level $4$ for
Lee--Yang, $6$ for Ising, $2$ for $c = 0$ (where $L_{-2}|0\rangle$ itself is null: $\langle T,T\rangle = c/2 = 0$).

\subsection{Characters of minimal models}

The irreducible vacuum module $L(c_{p,q},0)$ has character (Rocha-Caridi)
\begin{equation*}
 \chi_{0}(q) = \frac{q^{-c/24}}{\prod_{n\ge1}(1-q^n)}\sum_{k\in\Z}\bigl(q^{a(k)}-q^{b(k)}\bigr)
\end{equation*}
for suitable quadratic exponents $a(k),b(k)$. For Lee--Yang ($c=-22/5$, so $-c/24 = 11/60$) it is the Rogers--Ramanujan series $H(q)$,
$q^{-11/60}\chi_0 = \sum_{n\ge0}q^{n(n+1)}/((1-q)\cdots(1-q^n)) = 1+q^2+q^3+q^4+q^5+2q^6+\cdots$;
for Ising it is $1+q^2+q^3+2q^4+2q^5+3q^6+\cdots$. The library computes some of these coefficients
(Section~\ref{l06:sec:dims}).

\subsection{The Coulomb gas and screening charges}

Feigin and Fuchs, and in physics Dotsenko and Fateev, realised all these modules with a single free
boson $\varphi$ with a \emph{background charge}: modify the free-boson stress tensor
$T = -\frac12(\partial\varphi)^2$ to $T = -\frac12(\partial\varphi)^2 + i\sqrt2\,\alpha_0\partial^2\varphi$; then
$c = 1-24\alpha_0^2$ (in one common normalisation), and the vertex operators $e^{i\alpha\varphi}$ are
primary of weight $\alpha^2-2\alpha\alpha_0$ (again normalisation-dependent). Two values $\alpha_\pm$ give
fields $V_\pm = e^{i\alpha_\pm\varphi}$ of weight exactly one; their contour integrals
$Q_\pm = \oint V_\pm(z)\,dz$ commute with the Virasoro algebra (the integrand changes by a total
derivative). These are the \emph{screening charges}. Applying $r$ of them,
\begin{equation*}
 Q^r|p\rangle = \oint\!\cdots\!\oint dz_1\cdots dz_r\,\prod_{i<j}(z_i-z_j)^{2t}\prod_i z_i^{(\cdots)}
 \exp\Bigl(\beta\sum_{n>0}\frac{a_{-n}}{n}\sum_i z_i^n\Bigr)|p\rangle ,
\end{equation*}
produces a state in a Fock module that is annihilated by all $L_n$, $n>0$ --- a singular vector ---
at level $rs$ and highest weight $h_{r,s}$. When $t = \beta^2$ is a positive integer the factor
$\prod(z_i-z_j)^{2t}$ is a polynomial, the contour integral is the extraction of a single coefficient,
and everything is finite algebra. This is exactly the route the library takes
(Section~\ref{l06:sec:screening}): it constructs the screening states for $t = m^2$, proves they are
nonzero and singular, deduces that $\det G_{rs}$ vanishes on the Kac curve at $t = 1,4,9,\dots$, and then
uses that the determinant is a Laurent polynomial in $t$ to conclude it vanishes for every $t$.

\subsection{The Sugawara construction}

When a theory has currents $J^a(z)$ forming an affine Lie algebra at level $k$, its stress tensor can be
built from them: $T(z) = \frac{1}{2(k+h^\vee)}\sum_a{:}J^aJ^a{:}(z)$, the \emph{Sugawara construction},
with $c = k\dim\mathfrak g/(k+h^\vee)$. For one free boson ($\mathfrak g = \C$, $h^\vee = 0$) this is
$T = \frac{1}{2k}{:}JJ{:}$ with $c = 1$. The library's file \leanpath{VirasoroSugawara.lean} treats
exactly this abelian case (Section~\ref{l06:sec:sugawara}); the non-abelian affine $\mathfrak{sl}_2$
Sugawara vector is in Chapter~\ref{ch:L08}.

%======================================================================
\section{The Virasoro modes on the universal vacuum module}\label{l06:sec:modes}
%======================================================================

\subsection{The vacuum module, recalled}

\begin{definition}[Universal vacuum module; Chapter~\ref{ch:L02}]\label{l06:def:Vc}
Let $\lsixVir$ be the Virasoro Lie algebra with basis $L_n$ ($n\in\Z$), $C$ and bracket
$[L_m,L_n] = (m-n)L_{m+n}+\frac{m^3-m}{12}\delta_{m+n,0}C$, $C$ central, and let $\lsixU$ be its universal
enveloping algebra. For $c\in\C$ let $J(c)\subseteq\lsixU$ be the left ideal generated by $L_n$
($n\ge-1$) and $C-c$. The universal vacuum module is $\lsixV := \lsixU/J(c)$, with vacuum
$\lsixvac := [1]$ and $\lsixVir$ acting by left multiplication. Write $L_n$ also for the operator by
which $L_n$ acts.
\end{definition}

File \leanpath{VOA/Foundations/VirasoroModule.lean}.
\begin{leancode}
abbrev UVir : Type := UniversalEnvelopingAlgebra ℂ Virasoro

def vacGens (c : ℂ) : Set UVir :=
  {a | (∃ n : ℤ, -1 ≤ n ∧ a = ι ℂ (L n)) ∨ a = ι ℂ C - algebraMap ℂ UVir c}

noncomputable def vacIdeal (c : ℂ) : Submodule UVir UVir :=
  Submodule.span _ (vacGens c)

abbrev VirUniv (c : ℂ) : Type := UVir ⧸ vacIdeal c

noncomputable def virVacRho (c : ℂ) : Virasoro →ₗ⁅ℂ⁆ Module.End ℂ (VirUniv c) :=
  (Algebra.lsmul ℂ ℂ (VirUniv c)).toLieHom.comp (ι ℂ)

noncomputable def vacVec (c : ℂ) : VirUniv c := Submodule.Quotient.mk 1
\end{leancode}
File \leanpath{VOA/Examples/VirasoroField.lean}.
\begin{leancode}
noncomputable def Lop (n : ℤ) : End ℂ (VirUniv c) := virVacRho c (L n)

theorem Lop_bracket (m n : ℤ) :
    Lop c m * Lop c n - Lop c n * Lop c m
      = ((m : ℂ) - n) • Lop c (m + n)
        + (if m + n = 0 then ((m : ℂ) ^ 3 - m) / 12 * c else 0) • (1 : End ℂ (VirUniv c)) := by

theorem Lop_vacuum {n : ℤ} (hn : -1 ≤ n) : Lop c n (vacVec c) = 0 :=
\end{leancode}

\begin{dictionary}
\textbf{Decoding.} \lean{UniversalEnvelopingAlgebra ℂ Virasoro} is Mathlib's $\lsixU$, and \lean{ι ℂ}
the canonical Lie map $\lsixVir\to\lsixU$. A \lean{Submodule UVir UVir} is a left ideal ($\lsixU$
acting on itself by left multiplication), \lean{Submodule.span} is the left ideal generated by a set,
and \lean{⧸} is the quotient module. \lean{algebraMap ℂ UVir c} is the scalar $c\cdot1$. So
\lean{VirUniv c} is $\lsixV$ exactly. \lean{Algebra.lsmul} is left multiplication of $\lsixU$ on the
quotient, and composing with \lean{ι ℂ} gives the Lie action; \lean{Lop c n} is the operator $L_n$.
\lean{Lop_bracket} is \eqref{l06:eq:vir} as an identity of operators, with $C$ replaced by $c$ (this
uses that $C$ acts by $c$ on the whole module, \lean{vac_acts_central}, Chapter~\ref{ch:L02}); the term
\lean{if m + n = 0 then … else 0} is the Kronecker delta, and \lean{(1 : End ℂ (VirUniv c))} is the
identity operator. \lean{Lop_vacuum} is $L_n\lsixvac = 0$ for $n\ge-1$.

\textbf{The name.} The docstrings of \leanpath{VirasoroField.lean} and of some neighbouring files call
this module ``$L(c,0)$''. That is a historical misnomer: \lean{VirUniv c} is the \emph{universal}
module $\lsixV$, which differs from the simple $\lsixL$ at the degenerate central charges. The
library recorded the renaming (Chapter~\ref{ch:L02}); we write $\lsixV$ throughout.

\textbf{Vacuity.} If $\lsixV$ were the zero module every statement here would be trivially true. That
it is not is \lean{vacVec_ne_zero} (proved for every $c$, without \lean{VOA.pbw},
in \leanpath{VOA/\allowbreak Examples/\allowbreak VirasoroNonvanishing.lean}; Chapter~\ref{ch:L02} gives the statement and the
linear-dilaton proof, and we recall it in Theorem~\ref{l06:thm:nonvanishing}).
\end{dictionary}

\begin{physics}
$\lsixV$ is the conformal family of the identity operator, before any null state has been removed:
every state that can be created from the vacuum by the stress tensor. The relation $L_{-1}\lsixvac = 0$
is translation invariance of the vacuum, $L_0\lsixvac = 0$ says it has energy zero, and $L_n\lsixvac = 0$
for $n\ge1$ says $T(z)\lsixvac$ is regular at $z = 0$.
\end{physics}

\begin{theorem}[The vacuum is nonzero, for every $c$]\label{l06:thm:nonvanishing}
For every $c\in\C$, $\lsixvac\ne0$ in $\lsixV$.
\end{theorem}

File \leanpath{VOA/Examples/VirasoroNonvanishing.lean}.
\begin{leancode}
theorem exists_lam (c : ℂ) : ∃ lam : ℂ, 1 - 12 * lam ^ 2 = c := by

theorem vacVec_ne_zero (c : ℂ) : VirModule.vacVec c ≠ 0 := by
\end{leancode}

\begin{dictionary}
The statement is exactly ``$\lsixvac\ne0$''. Its only content is that the left ideal $J(c)$ does not
contain $1$. There are no hypotheses. Without it, $\lsixV = 0$ would be a model of every statement of
this chapter about $\lsixV$, so every later non-vacuity argument rests on this theorem.
\end{dictionary}

\begin{physics}
The vacuum exists at every central charge: the abstract ``conformal family of the identity'' is a
nonzero space. The witness is physical: a free boson with a background charge (a linear dilaton) has
a translation-invariant vacuum and central charge $1-12\lambda^2$, which takes every complex value.
\end{physics}

\begin{proof}[Proof on paper]
On the free-boson Fock space $\C[x_1,x_2,\dots]$ with oscillators $a_n$, the operators
$L^\lambda_n = \frac12\sum_m{:}a_{n-m}a_m{:}+\lambda(n+1)a_n$ satisfy \eqref{l06:eq:vir} with
$c = 1-12\lambda^2$ (\lean{Lld_virasoro}, \leanpath{VirasoroFock.lean}), and kill the Fock vacuum $1$ for
$n\ge-1$ (\lean{Lld_vacuum}: for $n=-1$ the background term carries the factor $n+1 = 0$). Choose
$\lambda$ with $1-12\lambda^2 = c$ (\lean{exists_lam}; $\C$ is algebraically closed). By the universal
property, $a\mapsto a\cdot1$ is a $\lsixU$-linear map $\lsixU\to\mathrm{Fock}$ killing the generators of
$J(c)$, hence all of $J(c)$; it therefore factors through $\lsixV$ and sends $\lsixvac$ to $1\ne0$.
\end{proof}

\begin{status}
\lean{VOA.vacVec_ne_zero}: \texttt{\#print axioms} gives \lean{[propext, Classical.}\allowbreak\lean{choice, Quot.}\allowbreak\lean{sound]}
(run for this chapter, Section~\ref{l06:sec:verification}). Unconditional, \lean{VOA.pbw}-free.
\end{status}

\subsection{The conformal field and its self-locality}

\begin{definition}[Conformal vector and conformal field]\label{l06:def:omega}
$\omega := L_{-2}\lsixvac\in\lsixV$, and $T(z) = Y(\omega,z) := \sum_{n\in\Z}L_n z^{-n-2}$, a formal
distribution with coefficients in $\End\lsixV$.
\end{definition}

File \leanpath{VOA/Examples/VirasoroField.lean}.
\begin{leancode}
noncomputable def omegaVir : VirUniv c := Lop c (-2) (vacVec c)

noncomputable def YomegaVir : FormalDist (End ℂ (VirUniv c)) := fun t => Lop c (-t - 2)

theorem mode_YomegaVir (m : ℤ) : mode (YomegaVir c) m = Lop c (m - 1) := by
\end{leancode}

\begin{dictionary}
\textbf{Index conventions --- read carefully.} A \lean{FormalDist A} is a function $\Z\to A$, the
value at $t$ being the coefficient of $z^t$ (Chapter~\ref{ch:L01}, Definition~\ref{l01:def:formaldist}).
So \lean{YomegaVir c t = Lop c (-t-2)} says the coefficient of $z^t$ is $L_{-t-2}$, i.e.\
$\sum_tL_{-t-2}z^t = \sum_nL_nz^{-n-2}$, which is \eqref{l06:eq:Tmodes}. The \emph{vertex-algebra mode}
\lean{mode f n} is the coefficient of $z^{-n-1}$, i.e.\ \lean{mode f n = f (-n-1)}
(Definition~\ref{l01:def:mode}). Hence \lean{mode (YomegaVir c) m} $= L_{m-1}$, i.e.\
$\omega_{(m)} = L_{m-1}$, the standard convention $L_n = \omega_{(n+1)}$.
\end{dictionary}

\begin{physics}
$\omega$ is the state of the stress tensor, and $Y(\omega,z) = T(z)$. The weight of $\omega$ is $2$:
the stress tensor has dimension $(2,0)$.
\end{physics}

\begin{theorem}[The $TT$ OPE and self-locality]\label{l06:thm:selflocal}
As formal distributions in two variables,
\begin{equation*}
 [T(z),T(w)] = \partial T(w)\,\delta(z-w) + 2T(w)\,\partial_w\delta(z-w) + \frac{c}{12}\,\partial_w^3\delta(z-w),
\end{equation*}
and consequently $(z-w)^4[T(z),T(w)] = 0$: $T$ is local with itself. Moreover $T$ is a field: for
every $v\in\lsixV$, $L_nv = 0$ for all sufficiently large $n$.
\end{theorem}

File \leanpath{VOA/Examples/VirasoroField.lean}.
\begin{leancode}
theorem bracketDist_YomegaVir_deltaDeriv :
    bracketDist (YomegaVir c) (YomegaVir c)
      = deltaDeriv 0 (YomegaVir_deriv c) + deltaDeriv 1 ((2 : ℂ) • YomegaVir c)
        + deltaDeriv 3 (YomegaVir_unit c) := by

theorem YomegaVir_self_local : Local (YomegaVir c) (YomegaVir c) := by
\end{leancode}
File \leanpath{VOA/Examples/VirasoroIsField.lean}.
\begin{leancode}
theorem YomegaVir_isField : IsField (YomegaVir c) := by
\end{leancode}

\begin{dictionary}
\lean{bracketDist a b} is the commutator $[a(z),b(w)]$ as a two-variable distribution;
\lean{deltaDeriv j f} is $f(w)\cdot\frac{1}{j!}\partial_w^j\delta(z-w)$ in the library's normalisation
(Chapter~\ref{ch:L01}); \lean{YomegaVir_deriv c} is $\partial T$ (its coefficient of $w^u$ is
$(u+1)L_{-u-3}$) and \lean{YomegaVir_unit c} is the constant distribution $\frac{c}{12}\cdot1$ sitting at
$w^0$. With $\frac1{3!}\partial^3$ in \lean{deltaDeriv 3}, the last term is $\frac{c}{12}\cdot\frac{1}{6}
\cdot 6\,\partial^3_w\delta = $ the $\frac{c/2}{(z-w)^4}$ pole of \eqref{l06:eq:TT} (the residue $c/2$ of the
fourth-order pole corresponds to $\frac{c}{2}\cdot\frac{1}{3!}\partial^3\delta$). \lean{Local a b} is
$\exists N,\ (z-w)^N[a(z),b(w)] = 0$ (Definition~\ref{l01:def:bracket}); the proof takes $N = 4$.
\lean{IsField f} is: for every vector $x$ there is $B$ with \lean{f j x = 0} for $j<-B$
(Definition~\ref{l01:def:field}), i.e.\ $L_nx = 0$ for $n$ large. These are the standard notions, and the
statements are exactly the theorem.
\end{dictionary}

\begin{physics}
This is the $TT$ OPE \eqref{l06:eq:TT} written as a commutator of fields: the $(z-w)^{-4}$ term is the
central charge (the two-point function $\langle T(z)T(w)\rangle = \frac{c/2}{(z-w)^4}$), the
$(z-w)^{-2}$ term says $T$ has weight $2$, the $(z-w)^{-1}$ term says $T$ generates translations.
Locality ($T(z)T(w) = T(w)T(z)$ up to singularities at $z = w$) is the chiral version of the statement
that the stress tensor is a local observable.
\end{physics}

\begin{proof}[Proof on paper]
The coefficient of $z^sw^t$ of $[T(z),T(w)]$ is $[L_{-s-2},L_{-t-2}]$, which by \eqref{l06:eq:vir}
equals $(t-s)L_{-s-t-4}$ plus the central term at $s+t+4 = 0$. Expanding the three
$\delta$-derivative terms coefficientwise produces exactly the same numbers
(\lean{bracketDist_YomegaVir}); multiplication by $(z-w)^{j+1}$ kills $\partial^j\delta(z-w)$, so
$(z-w)^4$ kills all three. For the field property: $\lsixV$ is spanned by $a\cdot\lsixvac$ with
$a\in\lsixU$; $L_m\lsixvac = 0$ for $m\ge-1$, and commuting $L_m$ past a generator $\iota(x)$ of
$\lsixU$ gives $[L_m,x]$, which for $m$ large is again a mode of large index (plus a central term that
vanishes for $m$ large); induction on the length of $a$ gives a bound (\lean{Lop_ann}).
\end{proof}

\begin{status}
\lean{VOA.}\allowbreak\lean{VirModule.}\allowbreak\lean{YomegaVir_}\allowbreak\lean{self_}\allowbreak\lean{local}, \lean{VOA.}\allowbreak\lean{VirModule.}\allowbreak\lean{YomegaVir_}\allowbreak\lean{isField}:
\lean{[propext, Classical.}\allowbreak\lean{choice, Quot.}\allowbreak\lean{sound]}. Unconditional. No free-field realisation is used: the
locality comes from the abstract relation \eqref{l06:eq:vir} alone.
\end{status}

%======================================================================
\section{The grading and the partition basis of the vacuum module}\label{l06:sec:basis}
%======================================================================

\subsection{Creation words and their energy}

\begin{definition}[Creation monomials and the energy grading]\label{l06:def:grading}
For a list $w = [i_1,\dots,i_k]$ of natural numbers put
\begin{equation*}
 L_w\lsixvac := L_{-i_1-2}L_{-i_2-2}\cdots L_{-i_k-2}\lsixvac,\qquad
 E(w) := \sum_j(i_j+2),
\end{equation*}
and let $\lsixV_n$ be the span of all $L_w\lsixvac$ with $E(w) = n$ (and $\lsixV_n = 0$ for $n<0$).
\end{definition}

File \leanpath{VOA/Examples/VirasoroGrading.lean}.
\begin{leancode}
def energyVir : List ℕ → ℕ
  | [] => 0
  | (i :: rest) => (i + 2) + energyVir rest

noncomputable def gradingVir (n : ℤ) : Submodule ℂ (VirUniv c) :=
  Submodule.span ℂ {x | ∃ ws : List ℕ, (energyVir ws : ℤ) = n ∧ negWordVir c ws = x}
\end{leancode}
File \leanpath{VOA/Examples/VirasoroSpanning.lean}.
\begin{leancode}
noncomputable def negWordVir : List ℕ → VirUniv c
  | [] => vacVec c
  | n :: ws => Lop c (-(n : ℤ) - 2) (negWordVir ws)
\end{leancode}

\begin{dictionary}
\textbf{The shift by two.} The list entry $i$ stands for the mode $L_{-i-2}$, so the entries are
$\ge0$ exactly when the modes are $L_{-2},L_{-3},\dots$ --- the modes that act nontrivially on the
vacuum. This encoding makes it impossible to write $L_{-1}$ or a positive mode in a creation word.
\lean{gradingVir c n} is the $\C$-span (\lean{Submodule.span ℂ}) of the monomials of energy $n$; for
$n<0$ the set is empty and the span is $0$. Note that the definition does \emph{not} say the
monomials are independent, nor that $\lsixV$ is the sum of these spaces; both are theorems below.
\end{dictionary}

\begin{physics}
$L_{-i_1-2}\cdots L_{-i_k-2}\lsixvac$ is the state of the normal-ordered product of derivatives of the
stress tensor, ${:}\partial^{i_1}T\cdots\partial^{i_k}T{:}$ (up to factorials). Its energy $E(w)$ is its
scaling dimension.
\end{physics}

\begin{theorem}[$\lsixV$ is graded by $L_0$ with finite-dimensional pieces]\label{l06:thm:grading}
\begin{enumerate}
\item $L_0(L_w\lsixvac) = E(w)\,L_w\lsixvac$ for every word $w$.
\item The creation monomials span $\lsixV$.
\item $\lsixV = \bigoplus_{n\in\Z}\lsixV_n$ (an internal direct sum), $\lsixV_n = 0$ for $n<0$,
  $\lsixV_0 = \C\lsixvac$, and each $\lsixV_n$ is finite-dimensional.
\end{enumerate}
\end{theorem}

File \leanpath{VOA/Examples/VirasoroGrading.lean}.
\begin{leancode}
theorem L0_negWord (ws : List ℕ) :
    Lop c 0 (negWordVir c ws) = (energyVir ws : ℂ) • negWordVir c ws := by

theorem gradingVir_neg (n : ℤ) (hn : n < 0) : gradingVir c n = ⊥ := by

theorem gradingVir_zero : gradingVir c 0 = ℂ ∙ vacVec c := by

theorem gradingVir_isInternal : DirectSum.IsInternal (gradingVir c) := by

theorem gradingVir_finiteDimensional (n : ℤ) : FiniteDimensional ℂ (gradingVir c n) := by
\end{leancode}
File \leanpath{VOA/Examples/VirasoroSpanning.lean}.
\begin{leancode}
noncomputable def negSpanVir : Submodule ℂ (VirUniv c) :=
  Submodule.span ℂ (Set.range (negWordVir c))

theorem negSpanVir_eq_top : negSpanVir c = ⊤ := by
\end{leancode}

\begin{dictionary}
\lean{⊥} is the zero subspace, \lean{⊤} the whole space, \lean{ℂ ∙ v} the line spanned by $v$.
\lean{DirectSum.}\allowbreak\lean{IsInternal (gradingVir c)} is Mathlib's ``the canonical map
$\bigoplus_n\lsixV_n\to\lsixV$ is bijective'', i.e.\ $\lsixV$ is the internal direct sum of the
$\lsixV_n$. \lean{FiniteDimensional ℂ (gradingVir c n)} is literal finite-dimensionality. The theorem
is exactly the textbook statement; nothing about the \emph{dimensions} is claimed here (that is
Theorem~\ref{l06:thm:pbw}).
\end{dictionary}

\begin{physics}
$L_0$ is the Hamiltonian in radial quantisation; its eigenvalue on a descendant is the scaling dimension
$h+N$. The finite-dimensionality of each $\lsixV_n$ is what makes the partition function
$\operatorname{tr}q^{L_0-c/24}$ a well-defined formal power series.
\end{physics}

\begin{proof}[Proof on paper]
(1) Commute $L_0$ to the right: $[L_0,L_{-n}] = nL_{-n}$ with no central term (as $-n\ne0$), and
$L_0\lsixvac = 0$. (2) $\lsixV$ is generated by $\lsixvac$ over $\lsixU$ (cyclicity), and $\lsixU$ is
spanned by monomials in the $L_n$ and $C$; the span of the creation monomials contains $\lsixvac$ and is
stable under every $L_n$ (for $n\le-2$ by definition; for $n\ge-1$ by commuting $L_n$ to the right,
which produces shorter creation monomials and finally hits $L_n\lsixvac = 0$) and under $C = c$
(\lean{negSpan_rho_stable}); hence it is all of $\lsixV$. (3) By (1), $\lsixV_n$ lies in the
$n$-eigenspace of $L_0$; eigenspaces for distinct eigenvalues are independent, and by (2) the
$\lsixV_n$ span. A word of energy $n$ has length at most $n/2$ and entries at most $n$, so there are
finitely many of them (\lean{finite_energyEq}), giving finite dimension. $\lsixV_0$ is spanned by the
empty word, i.e.\ by $\lsixvac$.
\end{proof}

\begin{status}
\lean{VOA.}\allowbreak\lean{VirModule.}\allowbreak\lean{gradingVir_}\allowbreak\lean{isInternal}: \lean{[propext, Classical.}\allowbreak\lean{choice, Quot.}\allowbreak\lean{sound]}; the other
items of the theorem are used by \lean{virasoroVOA}, whose axiom list is the same. Unconditional and
\lean{VOA.pbw}-free: the internal direct sum comes from $L_0$-eigenspaces, not from a PBW basis (as for the
affine analogue in Chapter~\ref{ch:L02}, Section~\ref{l02:sec:grading}).
\end{status}

\subsection{The partition basis}

\begin{definition}[Sorted words]\label{l06:def:sorted}
A word $w$ is \emph{sorted} if it is weakly decreasing. A sorted word of energy $n$ is the same thing
as a partition of $n$ into parts $\ge2$ (part $i_j+2$). Write $\lsixpt(n)$ for the number of such
partitions.
\end{definition}

File \leanpath{VOA/Examples/VirasoroPartitionBasis.lean}.
\begin{leancode}
def SortedWord (w : List ℕ) : Prop := w.Pairwise (· ≥ ·)

def SortedOfEnergy (n : ℕ) : Type := {w : List ℕ // SortedWord w ∧ energyVir w = n}
\end{leancode}

\begin{dictionary}
\lean{w.Pairwise (· ≥ ·)} says every earlier entry is $\ge$ every later entry: weakly decreasing.
\lean{SortedOfEnergy n} is the subtype of such lists with energy $n$, a finite type. The same predicate
\lean{SortedWord} is reused for Verma modules (Section~\ref{l06:sec:verma}), where entry $i$ means
$L_{-i-1}$ instead; the reader must keep track of which shift is in force. The two conventions are
stated in the docstrings and we flag them each time.
\end{dictionary}

\begin{physics}
Partitions into parts $\ge2$ count the states of the identity family: $\prod_{n\ge2}(1-q^n)^{-1} =
1+q^2+q^3+2q^4+2q^5+4q^6+\cdots$. The absence of parts equal to $1$ is the absence of $L_{-1}\lsixvac$,
i.e.\ translation invariance of the vacuum.
\end{physics}

\begin{theorem}[PBW basis of $\lsixV$, at every $c$]\label{l06:thm:pbw}
For every $c\in\C$ and $n\in\N$, the monomials $L_w\lsixvac$ with $w$ sorted of energy $n$ form a
basis of $\lsixV_n$. Hence $\dim\lsixV_n = \lsixpt(n)$, and the sorted monomials of all energies form a
basis of $\lsixV$. In particular $\dim\lsixV_n = 1,0,1,1,2,2,4,4,7$ for $n = 0,\dots,8$ and
$\dim\lsixV_{16} = 55$.
\end{theorem}

File \leanpath{VOA/Examples/VirasoroPartitionBasis.lean}.
\begin{leancode}
noncomputable def gradedBasis (c : ℂ) (n : ℕ) : Module.Basis (SortedOfEnergy n) ℂ
    (gradingVir c (n : ℤ)) :=
  Module.Basis.mk (gradedWord_linearIndependent c n) (gradedWord_span c n)

theorem finrank_gradingVir (c : ℂ) (n : ℕ) :
    Module.finrank ℂ (gradingVir c (n : ℤ)) = Fintype.card (SortedOfEnergy n) :=

theorem sortedWord_linearIndependent (c : ℂ) :
    LinearIndependent ℂ (fun w : {w : List ℕ // SortedWord w} => negWordVir c w.1) := by
\end{leancode}
File \leanpath{VOA/Examples/VirasoroTopSymbol.lean}.
\begin{leancode}
def sortedWords (n : ℕ) : Finset (List ℕ) := sortedAux n n n

theorem mem_sortedWords (n : ℕ) (w : List ℕ) :
    w ∈ sortedWords n ↔ (SortedWord w ∧ energyVir w = n) := by

theorem finrank_gradingVir_eq_card (c : ℂ) (n : ℕ) :
    Module.finrank ℂ (gradingVir c (n : ℤ)) = (sortedWords n).card := by

theorem card_sortedWords_le_eight :
    (sortedWords 0).card = 1 ∧ (sortedWords 1).card = 0 ∧ (sortedWords 2).card = 1 ∧
      (sortedWords 3).card = 1 ∧ (sortedWords 4).card = 2 ∧ (sortedWords 5).card = 2 ∧
      (sortedWords 6).card = 4 ∧ (sortedWords 7).card = 4 ∧ (sortedWords 8).card = 7 :=

theorem finrank_gradingVir_sixteen (c : ℂ) : Module.finrank ℂ (gradingVir c 16) = 55 := by
\end{leancode}

\begin{dictionary}
\textbf{Decoding.} \lean{Module.Basis ι ℂ M} is a basis of $M$ indexed by $\iota$;
\lean{Module.Basis.mk} builds one from a linearly independent spanning family. So \lean{gradedBasis c n}
is a basis of $\lsixV_n$ indexed by partitions of $n$ into parts $\ge2$, whose vectors are (by
\lean{gradedWord_val}) the monomials $L_w\lsixvac$. \lean{Module.finrank} is the dimension;
\lean{Fintype.card} the number of elements of a finite type. \lean{sortedWords n} is a computable finite
set of lists, proved equal (\lean{mem_sortedWords}) to the set of sorted words of energy $n$, so that
the dimension at a given level is a kernel computation (\lean{card_}\allowbreak\lean{sortedWords_}\allowbreak\lean{le_}\allowbreak\lean{eight}).

\textbf{Uniformity in $c$.} The index type \lean{SortedOfEnergy n} does not depend on $c$; this is what
lets the Gram matrices of Section~\ref{l06:sec:gram} be matrices with $c$-independent rows and columns,
and their determinants polynomials in $c$.

\textbf{Could it be wrong?} A basis needs independence, which fails in many superficially similar
modules (e.g.\ in the simple quotient $\lsixL$). It holds here at \emph{every} $c$, including all the
minimal-model charges; that is correct, because $\lsixV$ is the universal module (free over the
subalgebra generated by $L_{-2},L_{-3},\dots$); the null vectors at degenerate $c$ are nonzero vectors
of $\lsixV$, not linear relations among the monomials.
\end{dictionary}

\begin{physics}
The universal vacuum character is $q^{-c/24}\prod_{n\ge2}(1-q^n)^{-1}$ for every $c$. All the physics of
special central charges is in the passage to the irreducible quotient, where null vectors are set to
zero.
\end{physics}

\begin{proof}[Proof on paper]
\emph{Spanning by sorted words.} Transposing two adjacent creation modes costs a merged mode:
$L_{-i-2}L_{-j-2} = L_{-j-2}L_{-i-2} + (j-i)L_{-(i+j+2)-2}$ (\lean{negWordVir_swap}; no central term since
$-i-2-j-2\ne0$). The merged word is shorter and has the same energy, so induction on energy and length
sorts every word (\lean{Lop_sorted_insert}).

\emph{Independence: the top symbol.} Realise $\lsixV$ in the free-boson Fock space
$\C[X_0,X_1,\dots]$ by the linear dilaton of Theorem~\ref{l06:thm:nonvanishing}. The creation mode
$L^\lambda_{-n}$ ($n\ge2$) splits, by polynomial degree, into multiplication by the quadratic
$P_n = \frac12\sum_{a+b=n-2}X_aX_b$ (degree $+2$), a degree-preserving part, and $\lambda(1-n)X_{n-1}$
(degree $+1$). Applying $k$ modes to the vacuum $1$, the degree-$2k$ component of the image of
$L_w\lsixvac$ is therefore $\prod_jP_{i_j+2}$ --- with \emph{no} $\lambda$ in it
(\lean{topSymbol_fockWordVir}). These products, over sorted words, are linearly independent: substitute
$X_p\mapsto t^{B^p}$ for a large base $B$; then $\prod_jP_{i_j+2}$ becomes a one-variable polynomial
whose degree encodes the partition in base $B$, so distinct partitions have distinct degrees. Given a
relation $\sum_wa_wL_w\lsixvac = 0$, map it to Fock, take the top-degree component of the longest words
present, and conclude their coefficients vanish; induct downward on length
(\lean{negWordVir_coef_zero}). Since $\lambda$ never enters the top symbol, the argument is uniform in $c$.
\end{proof}

\begin{status}
\lean{VOA.}\allowbreak\lean{VirModule.}\allowbreak\lean{gradedBasis}, \lean{VOA.}\allowbreak\lean{VirModule.}\allowbreak\lean{finrank_}\allowbreak\lean{gradingVir},
\lean{VOA.}\allowbreak\lean{VirModule.}\allowbreak\lean{finrank_}\allowbreak\lean{gradingVir_}\allowbreak\lean{sixteen}: \lean{[propext, Classical.}\allowbreak\lean{choice, Quot.}\allowbreak\lean{sound]}.
Unconditional, \lean{VOA.pbw}-free.
\end{status}

\begin{theorem}[Feigin--Fuchs embedding of the universal vacuum module]\label{l06:thm:ffembed}
For every $c$ and every $\lambda$ with $1-12\lambda^2 = c$, the map $\Phi_\lambda:\lsixV\to\mathrm{Fock}$,
$[a]\mapsto a\cdot1$ (linear-dilaton action), is injective; its kernel on $\lsixU$ is exactly $J(c)$.
\end{theorem}

File \leanpath{VOA/Examples/VirasoroTopSymbol.lean}.
\begin{leancode}
theorem evAlg_vacuum_eq_zero_iff {c lam : ℂ} (hlam : 1 - 12 * lam ^ 2 = c) (a : UVir) :
    evAlg lam a vacuum = 0 ↔ a ∈ vacIdeal c := by

theorem fockMap_injective {c lam : ℂ} (hlam : 1 - 12 * lam ^ 2 = c) :
    Function.Injective (fockMap hlam) := by

theorem exists_injective_fockMap (c : ℂ) :
    ∃ Φ : VirUniv c →ₗ[ℂ] Fock, Function.Injective Φ ∧ Φ (vacVec c) = 1 := by
\end{leancode}

\begin{dictionary}
\lean{evAlg lam} is the algebra map $\lsixU\to\End(\mathrm{Fock})$ extending the linear-dilaton
representation; \lean{vacuum} is the Fock vacuum $1$; \lean{fockMap hlam} is the induced linear map
$\lsixV\to\mathrm{Fock}$. The statements hold at \emph{every} $c$, the degenerate minimal-model charges
included. A reader may expect the opposite (a free-field realisation ``loses'' null vectors); the
docstring explains why it does not: injectivity is a statement about the universal module, whose null
vectors are nonzero vectors that map to nonzero (singular) vectors of the Fock module. The non-injective
maps in the Felder complex are the screening maps \emph{between} Fock modules, which this theorem does
not consider.
\end{dictionary}

\begin{physics}
Every state of the identity family is realised by a free boson with background charge (``Coulomb gas''),
faithfully, at every central charge. Null states of $\lsixV$ are realised by genuinely nonzero Fock
states; they become zero only in the quotient that defines the minimal model.
\end{physics}

\begin{proof}[Proof on paper]
Expand $[a]$ in the partition basis; its image is the same linear combination of the Fock images of the
sorted monomials, and these are independent by the top-symbol argument of Theorem~\ref{l06:thm:pbw}.
\end{proof}

\begin{status}
\lean{VOA.}\allowbreak\lean{VirModule.}\allowbreak\lean{fockMap_}\allowbreak\lean{injective}: \lean{[propext, Classical.}\allowbreak\lean{choice, Quot.}\allowbreak\lean{sound]}. Unconditional.
\end{status}

%======================================================================
\section{The universal Virasoro vertex operator algebra}\label{l06:sec:voa}
%======================================================================

\subsection{Word fields}

To make $\lsixV$ a vertex operator algebra one must assign to every state a field. The plan is the
reconstruction theorem (Chapter~\ref{ch:L03}, Theorem~\ref{l03:thm:reconstruct}):
give fields for a spanning family of states, check that they are mutually local, create their states
from the vacuum, and are translation covariant; then a unique vertex algebra structure exists.

\begin{definition}[Word fields]\label{l06:def:wordfields}
For $i\in\N$ let $T_i(z) := \frac{1}{i!}\partial_z^iT(z)$, the field of $L_{-i-2}\lsixvac$. For a word
$w = [i_1,\dots,i_k]$ define $Y_w(z) := {:}T_{i_1}(z){:}T_{i_2}(z)\cdots{:}T_{i_k}(z)\mathbf{1}{:}\cdots{:}$, the
iterated normal-ordered product (with $Y_{[\,]} = \mathrm{id}$).
\end{definition}

File \leanpath{VOA/Examples/VirasoroWordFields.lean}.
\begin{leancode}
noncomputable def genFieldVir (i : ℕ) : FormalDist (End ℂ (VirUniv c)) :=
  (i.factorial : ℂ)⁻¹ • (VOA.deriv^[i] (YomegaVir c))

noncomputable def idFieldVir : FormalDist (End ℂ (VirUniv c)) :=
  fun j => if j = 0 then LinearMap.id else 0

noncomputable def YwordVirF :
    List ℕ → { f : FormalDist (End ℂ (VirUniv c)) // IsField f }
  | [] => ⟨idFieldVir c, idFieldVir_isField c⟩
  | (i :: rest) =>
      let p := YwordVirF rest
      ⟨nprod (genFieldVir_isField c i) p.2, IsField.nprod (genFieldVir_isField c i) p.2⟩

noncomputable def YwordVir (ws : List ℕ) : FormalDist (End ℂ (VirUniv c)) := (YwordVirF c ws).1

theorem YwordVir_local : ∀ (ws₁ ws₂ : List ℕ), Local (YwordVir c ws₁) (YwordVir c ws₂) := by

theorem YwordVir_vacuum_base (ws : List ℕ) :
    (YwordVir c ws) 0 (vacVec c) = negWordVir c ws := by
\end{leancode}

\begin{dictionary}
\lean{VOA.deriv^[i]} is the $i$-fold formal derivative $\partial_z^i$; \lean{nprod} is the normal-ordered
product ${:}a(z)b(z){:}$ of two fields (Chapter~\ref{ch:L03}), which needs proofs that both factors are
fields --- hence the recursive packaging \lean{YwordVirF} of a distribution with its field
certificate. \lean{YwordVir_local} is mutual locality of all word fields (via Dong's lemma,
Theorem~\ref{l03:thm:dong}); \lean{YwordVir_vacuum_base} says the $z^0$-coefficient of $Y_w(z)\lsixvac$ is
$L_w\lsixvac$, i.e.\ $Y_w$ creates the creation monomial (there is also \lean{YwordVir_vacuum_neg}: no
negative powers of $z$ appear in $Y_w(z)\lsixvac$).
\end{dictionary}

\begin{physics}
$Y_w$ is the field ${:}\partial^{i_1}T\cdots\partial^{i_k}T{:}$ (divided by factorials) --- all composite
operators that can be built from $T$. Their mutual locality is the statement that composites of local
fields are local.
\end{physics}

\begin{theorem}[Reconstruction of the vertex operators]\label{l06:thm:virYA}
There is a linear map $Y:\lsixV\to\End(\lsixV)[[z^{\pm1}]]$ with $Y(L_w\lsixvac,z) = Y_w(z)$ for every
word $w$, every $Y(u,z)$ a field, all mutually local, $Y(u,z)\lsixvac = u+O(z)$, and
$[L_{-1},Y(u,z)] = \partial_zY(u,z)$.
\end{theorem}

File \leanpath{VOA/Examples/VirasoroReconstruct.lean}.
\begin{leancode}
theorem virYA_exists :
    ∃ YA : VirUniv c →ₗ[ℂ] FormalDist (End ℂ (VirUniv c)),
      (∀ ws : List ℕ, YA (negWordVir c ws) = YwordVir c ws) ∧
      (∀ u : VirUniv c, IsField (YA u)) ∧
      (∀ u v : VirUniv c, Local (YA u) (YA v)) ∧
      (∀ u : VirUniv c, (YA u) 0 (vacVec c) = u) ∧
      (∀ u : VirUniv c, ∀ q : ℤ, q < 0 → (YA u) q (vacVec c) = 0) ∧
      (∀ u : VirUniv c, DCovariant (Ltrans c) (YA u)) :=
\end{leancode}
File \leanpath{VOA/Examples/VirasoroConformal.lean}.
\begin{leancode}
noncomputable def virYA : VirUniv c →ₗ[ℂ] FormalDist (End ℂ (VirUniv c)) := (virYA_exists c).choose

theorem virYA_omega : virYA c (omegaVir c) = YomegaVir c := by
\end{leancode}

\begin{dictionary}
\textbf{Decoding.} The six conjuncts are, in order: $Y$ agrees with the word fields on creation
monomials; every $Y(u)$ is a field; mutual locality; the $z^0$ coefficient of $Y(u,z)\lsixvac$ is $u$;
negative powers of $z$ in $Y(u,z)\lsixvac$ vanish; \lean{DCovariant (Ltrans c) f} is
$[L_{-1},f(z)] = \partial_zf(z)$ (\lean{Ltrans c = Lop c (-1)}, Definition~\ref{l03:def:dcov}).
\lean{virYA c} is \emph{a} map with these properties, chosen by \lean{Classical.choice} (\lean{.choose});
by Goddard uniqueness (Theorem~\ref{l03:thm:goddard}) there is only one, so the choice is harmless.
\lean{virYA_omega} says the chosen map sends $\omega$ to $T(z)$.

\textbf{Hypotheses.} None beyond $c\in\C$. The term proof (not shown in full) passes the word fields,
\lean{negSpanVir_eq_top}, \lean{YwordVir_local}, the creation tower and translation covariance of the
word fields to the general theorem \lean{Reconstruction.}\allowbreak\lean{reconstruct} of Chapter~\ref{ch:L03}.
\end{dictionary}

\begin{physics}
The state--field correspondence for the identity family: every state of $\lsixV$ is created from the
vacuum by a unique local field built from $T$. The theorem says that the $TT$ OPE \eqref{l06:eq:TT}
alone determines the whole chiral algebra.
\end{physics}

\begin{proof}[Proof on paper]
This is the reconstruction theorem: given a spanning family of states with mutually local fields
that create them and are translation covariant, define $Y$ on the span; well-definedness (two words with
the same state give the same field) follows from Goddard's uniqueness: a field $F$ that is local with
every $Y_w$ and satisfies $F(z)\lsixvac = 0$ is zero (\lean{field_}\allowbreak\lean{eq_}\allowbreak\lean{zero_}\allowbreak\lean{of_}\allowbreak\lean{vacuum}), because
$F(z)Y_w(x)\lsixvac$ and $Y_w(x)F(z)\lsixvac = 0$ agree after multiplying by a power of $(z-x)$, and setting
$x = 0$ gives $F(z)L_w\lsixvac = 0$ for all $w$, which span. Locality, the creation property and
covariance of the extended fields are then inherited from the generators by linearity and Dong's lemma.
\end{proof}

\begin{status}
\lean{VOA.}\allowbreak\lean{VirModule.}\allowbreak\lean{virYA_}\allowbreak\lean{exists}: \lean{[propext, Classical.}\allowbreak\lean{choice, Quot.}\allowbreak\lean{sound]}. Unconditional.
\end{status}

\subsection{The vertex operator algebra}

\begin{definition}[The universal Virasoro VOA $\lsixV$]\label{l06:def:virVOA}
For $c\in\C$, the vertex operator algebra with state space $\lsixV$, grading $\lsixV_n$, vertex
operators $Y$ of Theorem~\ref{l06:thm:virYA}, vacuum $\lsixvac$, conformal vector $\omega$ and central
charge $c$.
\end{definition}

File \leanpath{VOA/Examples/VirasoroInstance.lean}.
\begin{leancode}
noncomputable def virasoroVOA : VertexOperatorAlgebra (VirUniv c) where
  grading := gradingVir c
  isInternal := gradingVir_isInternal c
  finrank := gradingVir_finiteDimensional c
  Y := virYA c
  vacuum := vacVec c
  vacuum_mem := vacVec_mem_gradingVir_zero c
  grading_mode := fun u hu v hv n => virYA_grading_mode c u hu v hv n
  vacuum_Y := virYA_vacuum c
  creation := virYA_creation c
  locality := virYA_local c
  regularity := virYA_regularity c
  conformalVector := omegaVir c
  conformal_mem := omegaVir_mem_gradingVir_two c
  centralCharge := c
  virasoro := virYA_virasoro c
  L0_grading := fun {_} v hv => virYA_L0_grading c v hv
  translation := virYA_omega_translation c
  vacuum_grade := gradingVir_zero c
  positivity := gradingVir_neg c

theorem virasoroVOA_centralCharge : (virasoroVOA c).centralCharge = c := rfl

theorem omegaVir_ne_zero : omegaVir c ≠ 0 := LeeYang.L2vac_ne_zero c
\end{leancode}
File \leanpath{VOA/Examples/VirasoroConformal.lean}.
\begin{leancode}
theorem virYA_virasoro (m n : ℤ) :
    ⁅mode (virYA c (omegaVir c)) (m + 1), mode (virYA c (omegaVir c)) (n + 1)⁆
      = ((m : ℂ) - n) • mode (virYA c (omegaVir c)) (m + n + 1)
        + (if m + n = 0 then (((m : ℂ) ^ 3 - m) / 12 * c) • (1 : End ℂ (VirUniv c)) else 0) := by
\end{leancode}

\begin{dictionary}
\textbf{Every field of the structure} \lean{VertexOperatorAlgebra} (Definition~\ref{l03:def:VOA}) is
discharged by a named theorem. Reading them against the definition: \lean{grading}/\lean{isInternal}/
\lean{finrank} are Theorem~\ref{l06:thm:grading} (the field named \lean{finrank} is the
finite-dimensionality of each $\lsixV_n$); \lean{grading_mode} is $u_{(n)}\lsixV_\ell\subseteq
\lsixV_{k+\ell-n-1}$ for $u\in\lsixV_k$ (\lean{virYA_grading_mode}, \leanpath{VirasoroGradedField.lean});
\lean{vacuum_Y} is $Y(\lsixvac,z) = \mathrm{id}$; \lean{creation} is $u_{(-1)}\lsixvac = u$ and $u_{(n)}\lsixvac = 0$
for $n\ge0$; \lean{locality} and \lean{regularity} ($u_{(n)}v = 0$ for $n\gg0$) are from
Theorem~\ref{l06:thm:virYA}; \lean{virasoro} is \lean{virYA_virasoro}: the modes
$L_n = \omega_{(n+1)}$ satisfy \eqref{l06:eq:vir} with central charge \lean{c} (the bracket \lean{⁅·,·⁆}
on $\End$ is the commutator); \lean{L0_grading} is $L_0 = n$ on $\lsixV_n$; \lean{translation} is
$Y(L_{-1}u,z) = \partial_zY(u,z)$; \lean{vacuum_grade} is $\lsixV_0 = \C\lsixvac$; \lean{positivity} is
$\lsixV_n = 0$ for $n<0$.

\textbf{Universal, not simple.} The docstring of \lean{virasoroVOA} reads ``\emph{the general
central-charge Virasoro vertex operator algebra} $\mathrm{Vir}(c) = L(c,0)$, the universal Virasoro vacuum
VOA''. The equation with $L(c,0)$ is a misnomer: this is the \emph{universal} $\lsixV$
(Warning~\ref{l06:warn:docstrings}); the simple one is Definition~\ref{l06:def:simple}.

\textbf{Non-degeneracy.} \lean{omegaVir_ne_zero}: $\omega\ne0$, for every $c$. (An earlier version assumed
$c\ne0$, because its proof used $L_2\omega = \frac c2\lsixvac$; for $c = 0$, $\omega\ne0$ too --- it is a basis
vector by Theorem~\ref{l06:thm:pbw}. This book flagged the hypothesis as redundant (Warning~\ref{l06:warn:omega}), and the library has
since removed it: the current statement, quoted above, is hypothesis-free and is proved via the Lee--Yang
computation \lean{LeeYang.L2vac_ne_zero}.)
\end{dictionary}

\begin{physics}
The chiral algebra of a CFT of central charge $c$ generated by the stress tensor alone, with no relation
imposed. Every CFT with central charge $c$ contains a quotient of it (the vacuum sector of its Virasoro
symmetry).
\end{physics}

\begin{proof}[Proof on paper]
All axioms but three are restatements of Theorems~\ref{l06:thm:grading} and~\ref{l06:thm:virYA}.
\emph{Virasoro relation}: since $Y(\omega) = T$, the modes $\omega_{(n+1)}$ are the operators $L_n$, which
satisfy \eqref{l06:eq:vir} by Theorem~\ref{l06:thm:selflocal}'s source \lean{Lop_bracket}.
\emph{Translation}: both $Y(L_{-1}v,z)$ and $\partial_zY(v,z)$ are $L_{-1}$-covariant, local with all word
fields, and agree on the vacuum to first order; Goddard uniqueness identifies them.
\emph{Grading of modes}: each word field is a ``graded field of weight $E(w)$'' (its $z^j$ coefficient
shifts the weight by $E(w)+j$), a property preserved by derivatives and normal-ordered products
(\lean{GradedFieldVir.nprod}).
\end{proof}

\begin{status}
\lean{VOA.}\allowbreak\lean{VirModule.}\allowbreak\lean{virasoroVOA}, \lean{VOA.}\allowbreak\lean{VirModule.}\allowbreak\lean{omegaVir_}\allowbreak\lean{ne_}\allowbreak\lean{zero}:
\lean{[propext, Classical.}\allowbreak\lean{choice, Quot.}\allowbreak\lean{sound]}. Unconditional, for every $c\in\C$, \lean{VOA.pbw}-free.
\end{status}

\subsection{The free boson at \texorpdfstring{$c = 1$}{c = 1}: the abelian Sugawara construction}\label{l06:sec:sugawara}

Before the general construction existed, the library realised a Virasoro VOA at $c = 1$ inside the
Heisenberg VOA of Chapter~\ref{ch:L05}. It remains in the library.

\begin{theorem}[Sugawara at $c = 1$]\label{l06:thm:sugawara}
For every $k\ne0$, in the rank-one Heisenberg VOA at level $k$ with conformal vector
$\omega = \frac1{2k}a_{-1}^2\lsixvac$, the modes $L_n = \omega_{(n+1)}$ satisfy \eqref{l06:eq:vir} with
$c = 1$. At $k = 1$ this VOA is recorded as \lean{virasoroVOA1}.
\end{theorem}

File \leanpath{VOA/Examples/VirasoroSugawara.lean}.
\begin{leancode}
noncomputable def L (k : ℂ) (_hk : k ≠ 0) (n : ℤ) : End ℂ Fock :=
  mode (Y k (ω k)) (n + 1)

theorem virasoro_relation_c1 (k : ℂ) (hk : k ≠ 0) (m n : ℤ) :
    ⁅L k hk m, L k hk n⁆
      = ((m : ℂ) - n) • L k hk (m + n)
        + (if m + n = 0 then ((m : ℂ) ^ 3 - m) / 12 else 0) • (1 : End ℂ Fock) := by
\end{leancode}
File \leanpath{VOA/Examples/VirasoroVOA.lean}.
\begin{leancode}
noncomputable def virasoroVOA1 : VertexOperatorAlgebra Fock :=
  heisenbergVOA 1 one_ne_zero

theorem virasoroVOA1_centralCharge : virasoroVOA1.centralCharge = 1 := rfl
\end{leancode}

\begin{dictionary}
\lean{Y k}, \lean{ω k}, \lean{heisenbergVOA} are the Heisenberg VOA of Chapter~\ref{ch:L05} at level
$k$. \textbf{Caution on what \lean{virasoroVOA1} is}: it is literally the Heisenberg VOA on the full Fock
space (\lean{heisenbergVOA 1 one_}\allowbreak\lean{ne_}\allowbreak\lean{zero}), not the Virasoro sub-VOA generated by $\omega$. It ``is the
$c = 1$ Virasoro VOA'' only in the sense that its conformal vector generates a Virasoro algebra of central
charge $1$. The file's header (written before Section~\ref{l06:sec:voa} existed) still says the general-$c$
Virasoro VOA is ``the documented next step''; that is outdated (Warning~\ref{l06:warn:docstrings}).
\end{dictionary}

\begin{physics}
The stress tensor of a free boson is $T = \frac1{2k}{:}JJ{:}$, the Sugawara construction for the abelian
current algebra, with $c = 1$: one bosonic degree of freedom. The non-abelian version,
$T = \frac{1}{2(k+h^\vee)}\sum{:}J^aJ^a{:}$, is treated for $\widehat{\mathfrak{sl}}_2$ in Chapter~\ref{ch:L08}.
\end{physics}

\begin{proof}[Proof on paper]
The Heisenberg relations $[a_m,a_n] = km\delta_{m+n,0}$ and the explicit formula
$L_n = \frac1{2k}\sum_j{:}a_{n-j}a_j{:}$ give $[L_m,a_n] = -na_{m+n}$ and then \eqref{l06:eq:vir} with
$c = 1$ by the standard normal-ordering computation (done in \leanpath{HeisenbergCentral.lean},
Chapter~\ref{ch:L05}); this file transports it to the VOA modes.
\end{proof}

\begin{status}
\lean{VOA.}\allowbreak\lean{VirasoroVOA.}\allowbreak\lean{virasoro_}\allowbreak\lean{relation_}\allowbreak\lean{c1}, \lean{VOA.virasoroVOA1}: \lean{[propext, Classical.}\allowbreak\lean{choice,
Quot.}\allowbreak\lean{sound]}. Unconditional.
\end{status}

\begin{definition}[Minimal-model central charges]\label{l06:def:cpq}
$c_{p,q} := 1-6(p-q)^2/(pq)\in\Q$ for $p,q\in\N$.
\end{definition}

File \leanpath{VOA/Examples/VirasoroVOA.lean}.
\begin{leancode}
def minimalCentralCharge (p q : ℕ) : ℚ := 1 - 6 * ((p : ℚ) - q) ^ 2 / (p * q)

theorem minimalCentralCharge_ising : minimalCentralCharge 3 4 = 1 / 2 := by
\end{leancode}

\begin{dictionary}
A rational number, cast silently to $\C$ when used as a central charge. Lean defines division by zero as
zero, so $c_{p,0} = 1$; every theorem using it assumes $p,q\ge2$. The formula is symmetric in $(p,q)$.
Examples: $c_{2,3} = 0$, $c_{2,5} = -22/5$, $c_{3,4} = 1/2$, $c_{2,7} = -68/7$, $c_{4,5} = 7/10$.
\end{dictionary}

\begin{physics}
The BPZ minimal-model central charges: $c_{3,4} = \frac12$ is the critical Ising model,
$c_{2,5} = -\frac{22}5$ the Lee--Yang edge singularity, $c_{m,m+1}$ the unitary series of Friedan, Qiu
and Shenker.
\end{physics}

%======================================================================
\section{The simple quotient \texorpdfstring{$L(c,0)$}{L(c,0)}}\label{l06:sec:simple}
%======================================================================

\begin{definition}[Simple Virasoro VOA]\label{l06:def:simple}
$\lsixL := \lsixV/I_{\max}(c)$, where $I_{\max}(c)$ is the sum of all proper graded ideals of $\lsixV$,
with the quotient VOA structure.
\end{definition}

File \leanpath{VOA/Core/QuotientVOA.lean}.
\begin{leancode}
def maxGradedIdeal (W : VertexOperatorAlgebra 𝒱) : Submodule ℂ 𝒱 :=
  sSup (ProperGradedIdeals W)
\end{leancode}
File \leanpath{VOA/Examples/VirasoroSimple.lean}.
\begin{leancode}
noncomputable def maxIdealVir (c : ℂ) : Submodule ℂ (VirUniv c) :=
  maxGradedIdeal (virasoroVOA c)

noncomputable def virasoroSimple (c : ℂ) :
    VertexOperatorAlgebra (VirUniv c ⧸ maxIdealVir c) :=
  maxQuotientVOA (virasoroVOA c)
\end{leancode}

\begin{dictionary}
\lean{ProperGradedIdeals W} is the set of graded ideals not containing the vacuum, and \lean{sSup} their
sum (Chapter~\ref{ch:L04}, ``maximal ideal''). \lean{maxQuotientVOA} is the quotient VOA by it
(Definition~\ref{l04:def:quot}). The central charge of the quotient is still $c$
(\lean{virasoroSimple_}\allowbreak\lean{centralCharge}, by \lean{rfl}); its $n$-th graded piece is the image of $\lsixV_n$
(\lean{virasoroSimple_grading}); the projection intertwines all products, $[u_{(n)}v] = [u]_{(n)}[v]$
(\lean{virasoroSimple_}\allowbreak\lean{mkQ_}\allowbreak\lean{smul}).
\end{dictionary}

\begin{physics}
$\lsixL$ is the identity sector of the CFT with all null states removed: the chiral algebra that is
actually realised. At generic $c$ it equals $\lsixV$; at the minimal-model charges it is the vacuum
module of the minimal model.
\end{physics}

\begin{theorem}[$\lsixL$ is simple and nondegenerate, for every $c$]\label{l06:thm:simple}
For every $c\in\C$: $\lsixL$ is a simple VOA (nonzero, and its only ideals are $0$ and itself); its
vacuum is nonzero; for $c\ne0$ its conformal vector is nonzero; and $\lsixV$ itself is simple if and
only if $I_{\max}(c) = 0$.
\end{theorem}

File \leanpath{VOA/Core/BilinearForm.lean}.
\begin{leancode}
def IsSimple (W : VertexOperatorAlgebra 𝒱) : Prop :=
  (⊤ : Submodule ℂ 𝒱) ≠ ⊥ ∧ ∀ I : Submodule ℂ 𝒱, W.IsIdeal I → I = ⊥ ∨ I = ⊤
\end{leancode}
File \leanpath{VOA/Examples/VirasoroSimple.lean}.
\begin{leancode}
theorem virasoroSimple_isSimple : (virasoroSimple c).IsSimple :=
  isSimple_maxQuotientVOA' (virasoroVOA c) (vacVec_ne_zero c)

theorem virasoroSimple_vacuum_ne_zero : (virasoroSimple c).vacuum ≠ 0 := by

theorem virasoroSimple_conformalVector_ne_zero (hc : c ≠ 0) :
    (virasoroSimple c).conformalVector ≠ 0 := by

theorem virasoroVOA_isSimple_iff : (virasoroVOA c).IsSimple ↔ maxIdealVir c = ⊥ :=
\end{leancode}

\begin{dictionary}
\lean{IsSimple} is the standard definition; \lean{W.IsIdeal I} is a two-sided ideal ($u_{(n)}I\subseteq I$
and $I_{(n)}u\subseteq I$; Chapter~\ref{ch:L04}). The proof term of \lean{virasoroSimple_}\allowbreak\lean{isSimple} is the
general theorem ``a VOA with nonzero vacuum has a simple maximal quotient'' fed with
\lean{vacVec_ne_zero}. \textbf{Graded versus all ideals.} \lean{IsSimple} quantifies over \emph{all}
ideals, but $I_{\max}$ is a sum of \emph{graded} ideals. The gap is closed by
\lean{QuotVOA.}\allowbreak\lean{isGradedIdeal_}\allowbreak\lean{of_}\allowbreak\lean{isIdeal} (every ideal of a VOA is graded, as it is $L_0$-stable and $L_0$
acts semisimply with integer eigenvalues); so the statement is the textbook one.

\textbf{At $c = 0$.} The theorem holds for $c = 0$, where $\lsixL = \C\lsixvac$ is one-dimensional
(Section~\ref{l06:sec:dims}); simplicity is then trivially true, and the conformal vector is $0$ --- which
is why \lean{virasoroSimple_}\allowbreak\lean{conformalVector_}\allowbreak\lean{ne_}\allowbreak\lean{zero} assumes $c\ne0$. The theorem does not say which $c$
have $I_{\max}(c)\ne0$; that is Kac theory (Section~\ref{l06:sec:kac}).
\end{dictionary}

\begin{physics}
Removing all null states gives an irreducible chiral algebra. The conformal vector survives (the stress
tensor is not null) exactly when $\langle T,T\rangle = c/2\ne0$; at $c = 0$ the stress tensor is null and
the irreducible vacuum sector is trivial, as for the trivial minimal model $M(3,2)$.
\end{physics}

\begin{proof}[Proof on paper]
The sum of all proper graded ideals is proper: a proper graded ideal meets $\lsixV_0 = \C\lsixvac$ in
zero, so the sum lies in $\bigoplus_{n\ne0}\lsixV_n$, which does not contain $\lsixvac\ne0$. Ideals of the
quotient pull back to ideals containing $I_{\max}$, which are $I_{\max}$ or everything. For $c\ne0$:
$\omega_{(3)}\omega = L_2\omega = \frac c2\lsixvac$, so an ideal containing $\omega$ contains $\lsixvac$.
The last claim: if $I_{\max} = 0$ the projection is an isomorphism; conversely a simple VOA has no
proper nonzero ideal.
\end{proof}

\begin{status}
\lean{VOA.}\allowbreak\lean{VirSimple.}\allowbreak\lean{virasoroSimple_}\allowbreak\lean{isSimple}, \lean{…virasoroSimple_}\allowbreak\lean{conformalVector_}\allowbreak\lean{ne_}\allowbreak\lean{zero},
\lean{…virasoroVOA_}\allowbreak\lean{isSimple_}\allowbreak\lean{iff}: \lean{[propext, Classical.}\allowbreak\lean{choice, Quot.}\allowbreak\lean{sound]}. Unconditional.
\end{status}

The library also proves that $\lsixL$ is \emph{self-dual} in the sense of Chapter~\ref{ch:L12} (it has a
nondegenerate fully invariant bilinear form), for every $c$:

File \leanpath{VOA/Examples/VirasoroSelfDual.lean}.
\begin{leancode}
theorem virasoroSimple_grading_one : (virasoroSimple c).grading 1 = ⊥ := by

theorem virasoroSimple_selfDual : (virasoroSimple c).SelfDual :=
\end{leancode}
Here \lean{SelfDual W} is ``there is an invariant form $B$, fully invariant, with zero left radical''
(\leanpath{VOA/\allowbreak Modules/\allowbreak StronglyRational.lean}); for a simple VOA it is equivalent to
$L_1\lsixV_1 = 0$ (\lean{selfDual_iff}), which holds vacuously since $\lsixL_1 = 0$. Physically: the
contragredient of the vacuum module is itself, i.e.\ the two-point function of the identity sector is
nondegenerate. Axioms: \lean{[propext, Classical.}\allowbreak\lean{choice, Quot.}\allowbreak\lean{sound]}.

%======================================================================
\section{The invariant form on the vacuum module and its Gram matrices}\label{l06:sec:gram}
%======================================================================

\subsection{The canonical form}

\begin{definition}[Canonical invariant form on $\lsixV$]\label{l06:def:canB}
$\lsixform{\cdot}{\cdot}_c$ is the unique bilinear form on $\lsixV$ that is invariant in the sense of
Frenkel--Huang--Lepowsky (Chapter~\ref{ch:L04}) and normalised by $\lsixform{\lsixvac}{\lsixvac} = 1$. In
particular $\lsixform{L_na}{b} = \lsixform{a}{L_{-n}b}$ for all $n$.
\end{definition}

File \leanpath{VOA/Examples/VirasoroGram4.lean}.
\begin{leancode}
noncomputable def canB : InvariantForm (virasoroVOA c) :=
  canonicalInvariantForm (vacVec_ne_zero c) (virasoroVOA_qp c)

theorem canB_adj (n : ℤ) (a b : VirUniv c) :
    (canB c).form (Lop c n a) b = (canB c).form a (Lop c (-n) b) := by
\end{leancode}
File \leanpath{VOA/Examples/VirasoroFormSymmetry.lean}.
\begin{leancode}
theorem canB_isSymmetric (c : ℂ) : (canB c).IsSymmetric := by
\end{leancode}

\begin{dictionary}
\lean{InvariantForm W} and \lean{canonicalInvariantForm} are the invariant-form structure and Li's
construction of Chapter~\ref{ch:L04} (Theorem~\ref{l04:thm:canonical}); the two arguments discharge its
hypotheses: $\lsixvac\ne0$ and $L_1\lsixV_1 = 0$ (vacuous since $\lsixV_1 = 0$). The form is normalised
(\lean{canB_normalized}). \lean{canB_adj} is Virasoro adjointness. \lean{canB_isSymmetric} is symmetry
of this particular form; the docstring warns that symmetry of a general \lean{InvariantForm} is false
(\lean{not_forall_isSymmetric}, Chapter~\ref{ch:L04}); here it follows because a form with
$\lsixform{\lsixvac}{\lsixvac} = 1$ and adjointness is unique, and its transpose has the same properties.
\end{dictionary}

\begin{physics}
$\lsixform{a}{b}$ is the two-point function of the corresponding fields, normalised by
$\langle0|0\rangle = 1$; $L_n^\dagger = L_{-n}$ is the reality of the stress tensor (BPZ conjugation).
\end{physics}

\subsection{Gram matrices and the dimension of \texorpdfstring{$L(c,0)$}{L(c,0)}}

\begin{theorem}[Graded dimension of $\lsixL$ is a Gram rank]\label{l06:thm:dimrank}
Let $G_n(c)$ be the matrix of $\lsixform\cdot\cdot_c$ on $\lsixV_n$ in the partition basis. Then
$\dim\lsixL_n = \operatorname{rank}G_n(c)$ for every $c$ and $n$. In particular $\dim\lsixL_n =
\lsixpt(n)$ if $\det G_n(c)\ne0$, and $\dim\lsixL_n<\lsixpt(n)$ if $\det G_n(c) = 0$.
\end{theorem}

File \leanpath{VOA/Examples/VirasoroSimpleGram.lean}.
\begin{leancode}
theorem finrank_virasoroSimple_grading_partition (n : ℕ) :
    finrank ℂ ((virasoroSimple c).grading (n : ℤ))
      = (GramMat.canGram (vacVec_ne_zero c) (virasoroVOA_qp c) (n : ℤ)
          (gradedBasis c n)).rank :=

theorem finrank_virasoroSimple_lt_of_det_eq_zero {ι : Type*} [Fintype ι] [DecidableEq ι] (n : ℤ)
    (e : Basis ι ℂ ((virasoroVOA c).grading n))
    (h : (GramMat.canGram (vacVec_ne_zero c) (virasoroVOA_qp c) n e).det = 0) :
    finrank ℂ ((virasoroSimple c).grading n) < finrank ℂ ((virasoroVOA c).grading n) :=
\end{leancode}
File \leanpath{VOA/Core/GradedFormMatrix.lean}.
\begin{leancode}
noncomputable def canGram : Matrix ι ι ℂ := gramMatrix (canonicalInvariantForm hne hqp) n e
\end{leancode}

\begin{dictionary}
\lean{GramMat.}\allowbreak\lean{canGram hne hqp n e} is the matrix $(\lsixform{e_i}{e_j})_{i,j}$ of the canonical form on
the graded piece of weight $n$, in the basis $e$ (Chapter~\ref{ch:L04}, ``Gram-rank spine'').
\lean{Matrix.rank} is the rank. The equation is the precise form of ``$\lsixL_n = \lsixV_n/\mathrm{rad}$''.
\end{dictionary}

\begin{physics}
The number of physical (non-null) states at level $n$ of the identity sector equals the rank of the
matrix of two-point functions: null states are exactly the kernel.
\end{physics}

\begin{proof}[Proof on paper]
The radical of the canonical form on a VOA of this type equals the maximal ideal
(Theorem~\ref{l04:thm:lirad}; Li's theorem, Chapter~\ref{ch:L04}); the radical is graded and its
weight-$n$ part is the kernel of $G_n(c)$; so $\dim\lsixL_n = \dim\lsixV_n - \dim\ker G_n = \operatorname{rank}G_n$.
\end{proof}

\begin{status}
\lean{VOA.}\allowbreak\lean{VirSimpleGram.}\allowbreak\lean{finrank_}\allowbreak\lean{virasoroSimple_}\allowbreak\lean{grading_}\allowbreak\lean{partition}: \lean{[propext, Classical.}\allowbreak\lean{choice,
Quot.}\allowbreak\lean{sound]}. Unconditional.
\end{status}

\subsection{Explicit Kac determinants of the vacuum module at levels 4, 5, 6}

\begin{theorem}[Vacuum Gram matrices at levels $4,5,6$]\label{l06:thm:gram456}
In the bases $(L_{-2}^2\lsixvac,L_{-4}\lsixvac)$, $(L_{-5}\lsixvac,L_{-3}L_{-2}\lsixvac)$ and
$(L_{-6}\lsixvac,L_{-4}L_{-2}\lsixvac,L_{-3}^2\lsixvac,L_{-2}^3\lsixvac)$,
\begin{gather*}
 G_4 = \begin{pmatrix}\frac{c(c+8)}2 & 3c\\ 3c & 5c\end{pmatrix},\qquad
 G_5 = \begin{pmatrix}10c & 4c\\ 4c & c(c+6)\end{pmatrix},\\
 G_6 = \begin{pmatrix}
 \frac{35}2c & 5c & 18c & 24c\\
 5c & \frac{5c^2+16c}2 & 14c & \frac{9c^2+48c}2\\
 18c & 14c & 8c^2+36c & 30c\\
 24c & \frac{9c^2+48c}2 & 30c & \frac{3c(c+8)(c+16)}4
 \end{pmatrix},
\end{gather*}
and
\begin{equation*}
 \det G_4 = \tfrac12c^2(5c+22),\quad \det G_5 = 2c^2(5c+22),\quad
 \det G_6 = \tfrac34c^4(2c-1)(5c+22)^2(7c+68).
\end{equation*}
In particular $\det G_6(c) = 0$ iff $c\in\{0,\frac12,-\frac{22}5,-\frac{68}7\}$.
\end{theorem}

File \leanpath{VOA/Examples/VirasoroGram4.lean}.
\begin{leancode}
noncomputable def canGram4 : Matrix (Fin 2) (Fin 2) ℂ :=
  !![(canB c).form (L2L2vac c) (L2L2vac c), (canB c).form (L2L2vac c) (L4vac c);
     (canB c).form (L4vac c) (L2L2vac c), (canB c).form (L4vac c) (L4vac c)]

theorem canGram4_eq :
    canGram4 c = !![c * (c + 8) / 2, 3 * c; 3 * c, 5 * c] := by

theorem det_canGram4 : (canGram4 c).det = c ^ 2 * (5 * c + 22) / 2 := by
\end{leancode}
File \leanpath{VOA/Examples/VirasoroGram56.lean}.
\begin{leancode}
theorem canGram5_eq : canGram5 c = !![10 * c, 4 * c; 4 * c, c * (c + 6)] := by

theorem det_canGram5 : (canGram5 c).det = 2 * c ^ 2 * (5 * c + 22) := by

theorem canGram6_eq :
    canGram6 c =
      !![35 / 2 * c, 5 * c, 18 * c, 24 * c;
         5 * c, (5 * c ^ 2 + 16 * c) / 2, 14 * c, (9 * c ^ 2 + 48 * c) / 2;
         18 * c, 14 * c, 8 * c ^ 2 + 36 * c, 30 * c;
         24 * c, (9 * c ^ 2 + 48 * c) / 2, 30 * c,
           3 * c * (c + 8) * (c + 16) / 4] := by

theorem det_canGram6 :
    (canGram6 c).det = 3 / 4 * c ^ 4 * (2 * c - 1) * (5 * c + 22) ^ 2 * (7 * c + 68) := by

theorem det_canGram6_eq_zero_iff :
    (canGram6 c).det = 0 ↔ c = 0 ∨ c = isingC ∨ c = leeYangC ∨ c = m27C := by
\end{leancode}

\begin{dictionary}
\lean{!![a, b; c, d]} is Mathlib notation for an explicit matrix, rows separated by \lean{;}.
\lean{L2L2vac c} $= L_{-2}L_{-2}\lsixvac$, \lean{L4vac c} $= L_{-4}\lsixvac$ (defined in
\leanpath{LeeYangSingular.lean}, Chapter~\ref{ch:L19}); the level-5 and level-6 vectors are defined in
\leanpath{VirasoroGram56.lean} in the orders displayed (\lean{canGram5} uses $(L_{-5}\lsixvac,
L_{-3}L_{-2}\lsixvac)$, \lean{canGram6} uses $(L_{-6},L_{-4}L_{-2},L_{-3}L_{-3},L_{-2}^3)\lsixvac$). The
matrices are the Gram matrices of the canonical form (not of some other bilinear form; an earlier
version of the project used a form with the wrong normalisation, which is why the file stresses ``the
right form''). \lean{isingC} $= 1/2$, \lean{leeYangC} $= -22/5$, \lean{m27C} $= -68/7$ are complex
constants; \lean{isingC_}\allowbreak\lean{eq_}\allowbreak\lean{minimalCentralCharge} identifies $1/2$ with $c_{3,4}$, and similarly for $m27C$
and $c_{2,7}$.

\textbf{Consistency with Kac.} The degree of $\det G_n$ in $c$ equals the total number of parts over all
partitions of $n$ into parts $\ge2$: $2+1 = 3$ at level 4, $1+2 = 3$ at level 5, $1+2+2+3 = 8$ at level 6;
the zeros are at the minimal-model charges with $(p-1)(q-1)\le n$: $c_{2,3} = 0$ (level 2), $c_{2,5}$
(level 4), $c_{3,4}$ and $c_{2,7}$ (level 6).
\end{dictionary}

\begin{physics}
These are the first nontrivial vacuum Kac determinants. A zero of $\det G_n$ at $c$ is a central charge
at which the identity family has a null state at level $\le n$; the four values at level 6 are the trivial
model, Ising, Lee--Yang and $M(2,7)$. The factor $(5c+22)^2$ at level 6 is the Lee--Yang null vector of
level 4: its two level-6 descendants $L_{-2}\chi$ and $L_{-1}^2\chi$ are null, and the nullity of $G_6$ at
$c = -\frac{22}5$ is indeed $2$ (Theorem~\ref{l06:thm:dims}).
\end{physics}

\begin{proof}[Proof on paper]
Each entry is computed by moving the left monomial to the right as positive modes and commuting, using
\eqref{l06:eq:vir} and $L_n\lsixvac = 0$ for $n\ge-1$. For example
$\lsixform{L_{-4}\lsixvac}{L_{-2}^2\lsixvac} = \lsixform{\lsixvac}{L_4L_{-2}^2\lsixvac}$ and
$L_4L_{-2}L_{-2}\lsixvac = [L_4,L_{-2}]L_{-2}\lsixvac = 6L_2L_{-2}\lsixvac = 6\cdot\frac c2\lsixvac$, giving
$3c$. The determinants are then polynomial algebra.
\end{proof}

\begin{status}
\lean{VOA.}\allowbreak\lean{VirGram4.}\allowbreak\lean{canGram4_}\allowbreak\lean{eq}, \lean{…det_canGram4}, \lean{VOA.}\allowbreak\lean{VirGram56.}\allowbreak\lean{det_}\allowbreak\lean{canGram5},
\lean{…det_canGram6}, \lean{…det_}\allowbreak\lean{canGram6_}\allowbreak\lean{eq_}\allowbreak\lean{zero_}\allowbreak\lean{iff}: \lean{[propext, Classical.}\allowbreak\lean{choice, Quot.}\allowbreak\lean{sound]}.
Unconditional.
\end{status}

\subsection{The vacuum Gram determinant at every level: degree and leading coefficient}

\begin{definition}[Polynomial vacuum Gram matrix]\label{l06:def:partgram}
$G^{\mathrm{vac}}_n\in\mathrm{Mat}_{\lsixpt(n)}(\C[c])$ is the matrix whose $(\lambda,\mu)$ entry is the
polynomial in $c$ computing $\lsixform{L_\lambda\lsixvac}{L_\mu\lsixvac}_c$ by the commutation algorithm.
\end{definition}

File \leanpath{VOA/Examples/VirasoroPartitionObstruction.lean}.
\begin{leancode}
noncomputable def partitionGramPolynomial (n : ℕ) :
    Matrix (SortedOfEnergy n) (SortedOfEnergy n) (Polynomial ℂ) :=
  fun i j => wordPolynomial ((partitionWord n i).reverse.map Neg.neg ++ partitionWord n j)

theorem partitionGramPolynomial_eval (c : ℂ) (n : ℕ) :
    (partitionGramPolynomial n).map (Polynomial.eval c) =
      GramMat.gramMatrix (VirGram4.canB c) (n:ℤ) (gradedBasis c n) := by
\end{leancode}
File \leanpath{VOA/Examples/VirasoroGramLeadingCoefficient.lean}.
\begin{leancode}
theorem partitionDeterminant_natDegree (n : ℕ) :
    (partitionGramPolynomial n).det.natDegree = ∑ j : SortedOfEnergy n, j.val.length := by

theorem partitionDeterminant_leadingCoeff (n : ℕ) :
    (partitionGramPolynomial n).det.leadingCoeff =
      ∏ j : SortedOfEnergy n, diagonalContraction j.val := by
\end{leancode}
File \leanpath{VOA/Examples/VirasoroContractionDiagonal.lean}.
\begin{leancode}
noncomputable def contractionWeight (a : ℕ) : ℂ :=
  (((a:ℂ)+2)^3-((a:ℂ)+2))/12

noncomputable def diagonalContraction : List ℕ → ℂ
  | [] => 1
  | a::w => (((w.count a+1:ℕ):ℂ)*contractionWeight a)*diagonalContraction w
\end{leancode}

\begin{dictionary}
\lean{partitionWord n i} is the list of mode indices $-(i_j+2)$ of the partition; the word
\lean{(partitionWord n i).}\allowbreak\lean{reverse.}\allowbreak\lean{map Neg.}\allowbreak\lean{neg ++ partitionWord n j} is ``the adjoint of the left
monomial followed by the right monomial'', and \lean{wordPolynomial} is the vacuum coefficient of that
mode word, computed with $c$ as a formal variable (\leanpath{VirasoroUniversalPairing.lean},
\lean{universal_gram_entry}). \lean{partitionGramPolynomial_}\allowbreak\lean{eval} proves that evaluating at $c$ gives the
actual Gram matrix of the canonical form in the partition basis --- so this is not a separately postulated
polynomial. The leading coefficient is $\prod_\lambda\prod_k\bigl(\frac{k^3-k}{12}\bigr)^{m_k}m_k!$, $m_k$ the
multiplicity of the part $k$ in $\lambda$ (each repeated index contributes its running count, giving the
factorial).
\end{dictionary}

\begin{physics}
The highest power of $c$ comes from contracting each $L_{-k}$ on the left with an $L_{-k}$ on the right
through the central term $\frac{c}{12}(k^3-k)$ alone: the ``free-field'' or large-$c$ limit, in which the
modes $L_{-k}$ behave like independent oscillators. That the determinant has exactly the expected degree is
the statement that at large $c$ the vacuum family has no null states.
\end{physics}

\begin{proof}[Proof on paper]
Each entry has $c$-degree at most $\min(\ell(\lambda),\ell(\mu))$ (one power of $c$ per central
contraction); so the determinant has degree at most $\sum_\lambda\ell(\lambda)$, and the coefficient of
that power is the determinant of the ``top'' matrix of equal-length contractions. A content argument shows
this top matrix is diagonal (contractions pair equal modes, so off-diagonal equal-length partitions give
zero), with diagonal entries $\prod_k(\frac{k^3-k}{12})^{m_k}m_k!\ne0$.
\end{proof}

\begin{status}
\lean{VOA.}\allowbreak\lean{VirDeformModes.}\allowbreak\lean{NormalOrdering.}\allowbreak\lean{partitionDeterminant_}\allowbreak\lean{natDegree},
\lean{…partitionDeterminant_}\allowbreak\lean{leadingCoeff}: \lean{[propext, Classical.}\allowbreak\lean{choice, Quot.}\allowbreak\lean{sound]}. Unconditional.
\textbf{Not formalised}: the full factorisation of $\det G^{\mathrm{vac}}_n(c)$ at every level (the vacuum
Kac formula of Gorelik--Kac). What is proved in general is its simple zero at the first singular level
of each minimal model (Section~\ref{l06:sec:vacuumkac}).
\end{status}

%======================================================================
\section{Graded dimensions of the minimal models: Lee--Yang, Ising, \texorpdfstring{$M(2,7)$}{M(2,7)}}\label{l06:sec:dims}
%======================================================================

\begin{theorem}[Low-level dimensions of $L(c_{p,q},0)$]\label{l06:thm:dims}
\begin{enumerate}
\item Lee--Yang, $c = -\frac{22}5$: $\dim\lsixL_4 = 1$, $\dim\lsixL_5 = 1$, $\dim\lsixL_6 = 2$
  (against $2,2,4$ for $\lsixV$).
\item Ising, $c = \frac12$: $\dim\lsixL_6 = 3$; $M(2,7)$, $c = -\frac{68}7$: $\dim\lsixL_6 = 3$.
\item $c = 0$: $\dim\lsixL_6 = 0$.
\item For $c\notin\{0,-\frac{22}5\}$, $\dim\lsixL_4 = 2$; for $c\notin\{0,\frac12,-\frac{22}5,-\frac{68}7\}$,
  $\dim\lsixL_6 = 4$; at each of these four charges $\dim\lsixL_6<4$.
\item No linear isomorphism $L(-\frac{22}5,0)\to L(\frac12,0)$ maps the weight-6 space onto the weight-6
  space.
\item $c = 1$: the maximal ideal of $V(1,0)$ meets $V(1,0)_n$ trivially for $n = 4,5,6$.
\end{enumerate}
\end{theorem}

File \leanpath{VOA/Examples/VirasoroSimpleGram.lean}.
\begin{leancode}
theorem finrank_virasoroSimple_leeYang_four :
    finrank ℂ ((virasoroSimple leeYangC).grading 4) = 1 := by

theorem finrank_virasoroSimple_four_of_generic (h0 : c ≠ 0) (hLY : c ≠ leeYangC) :
    finrank ℂ ((virasoroSimple c).grading 4) = 2 := by

theorem finrank_virasoroSimple_six_of_generic (h0 : c ≠ 0) (hI : c ≠ isingC)
    (hLY : c ≠ leeYangC) (h27 : c ≠ m27C) :
    finrank ℂ ((virasoroSimple c).grading 6) = 4 := by

theorem finrank_virasoroSimple_six_lt_of_degenerate
    (h : c = 0 ∨ c = isingC ∨ c = leeYangC ∨ c = m27C) :
    finrank ℂ ((virasoroSimple c).grading 6) < 4 := by

theorem maxIdealVir_one_inf_grading_six_eq_bot :
    QuotVOA.maxGradedIdeal (virasoroVOA (1 : ℂ)) ⊓ (virasoroVOA (1 : ℂ)).grading 6 = ⊥ :=
\end{leancode}
File \leanpath{VOA/Examples/VirasoroSimpleSix.lean}.
\begin{leancode}
theorem finrank_virasoroSimple_six_zero :
    finrank ℂ ((virasoroSimple (0 : ℂ)).grading 6) = 0 :=

theorem finrank_virasoroSimple_six_ising :
    finrank ℂ ((virasoroSimple isingC).grading 6) = 3 :=

theorem finrank_virasoroSimple_six_m27 :
    finrank ℂ ((virasoroSimple m27C).grading 6) = 3 :=

theorem finrank_virasoroSimple_six_leeYang :
    finrank ℂ ((virasoroSimple leeYangC).grading 6) = 2 :=

theorem finrank_virasoroSimple_five_leeYang :
    finrank ℂ ((virasoroSimple leeYangC).grading 5) = 1 :=

theorem no_gradedEquiv_leeYang_ising :
    ¬ ∃ e : (VirUniv leeYangC ⧸ maxIdealVir leeYangC) ≃ₗ[ℂ] (VirUniv isingC ⧸ maxIdealVir isingC),
        ∀ x, x ∈ (virasoroSimple leeYangC).grading 6 ↔ e x ∈ (virasoroSimple isingC).grading 6 :=
\end{leancode}

\begin{dictionary}
\lean{leeYangC : ℂ := -22 / 5} is defined in \leanpath{VOA/\allowbreak Examples/\allowbreak LeeYangSingular.lean}
(Chapter~\ref{ch:L19}); \lean{isingC} and \lean{m27C} in \leanpath{VirasoroGram56.lean}. All statements are
about the actual simple quotient \lean{virasoroSimple c} at that $c$, through
Theorem~\ref{l06:thm:dimrank}. \lean{⊓} is intersection of subspaces; \lean{≃ₗ[ℂ]} a linear isomorphism.
Item 6 is honest about its scope: it says nothing about $V(1,0)_n$ for $n\ge7$. (Classically $V(1,0)$ is
simple at every level, since $c = 1$ is not of the form $c_{p,q}$ with $p,q\ge2$ coprime; the library does
not state this.) \textbf{What is not proved}: the full
characters. The dimensions given are coefficients of the Rogers--Ramanujan series ($q^4,q^5,q^6$) and of the
Ising vacuum character ($q^6$); they are \emph{evidence} that the definitions are right, not the character
identities.
\end{dictionary}

\begin{physics}
The vacuum characters of the minimal models begin
\begin{equation*}
 \text{Lee--Yang: }1+q^2+q^3+q^4+q^5+2q^6+\cdots,\qquad
 \text{Ising: }1+q^2+q^3+2q^4+2q^5+3q^6+\cdots,
\end{equation*}
against $1+q^2+q^3+2q^4+2q^5+4q^6$ for the universal module. The Lee--Yang null vector
$5L_{-2}^2\lsixvac-3L_{-4}\lsixvac$ at level $4$ removes one state at level 4, its descendant $L_{-1}\chi$
one at level 5, and two at level 6; the Ising null vector at level 6 removes one state. At $c = 0$ the
stress tensor is null, $\lsixL = \C\lsixvac$ and all positive levels vanish. Item 5 says Lee--Yang and Ising
have genuinely different chiral algebras, distinguished already by state counting.
\end{physics}

\begin{proof}[Proof on paper]
By Theorem~\ref{l06:thm:dimrank}, compute the rank of the explicit matrices of
Theorem~\ref{l06:thm:gram456} at the given $c$. Lee--Yang, level 4: $G_4(-\frac{22}5)$ is a nonzero matrix
with zero determinant, so rank $1$. Level 6 at Lee--Yang: the determinant vanishes to second order in
$(5c+22)$; the rank is at most $2$ by an explicit factorisation of $G_6(-\frac{22}5)$
(\lean{canGram6_leeYang_factor}) and at least $2$ by a nonzero $2\times2$ minor. Ising, level 6: the root is simple, and a nonzero $3\times3$ minor gives rank $3$.
At $c = 0$, $G_6(0) = 0$. Item 5: an isomorphism matching the weight-6 spaces would make their dimensions
equal, $2\ne3$.
\end{proof}

\begin{status}
All items: \lean{[propext, Classical.}\allowbreak\lean{choice, Quot.}\allowbreak\lean{sound]} (Section~\ref{l06:sec:verification}).
Unconditional.
\end{status}

%======================================================================
\section{Verma modules}\label{l06:sec:verma}
%======================================================================

\begin{definition}[Verma module $M(c,h)$]\label{l06:def:verma}
For $(c,h)\in\C^2$ let $I(c,h)\subseteq\lsixU$ be the left ideal generated by $L_n$ ($n>0$), $C-c$ and
$L_0-h$. The Verma module is $M(c,h) := \lsixU/I(c,h)$, with highest-weight vector $v_h := [1]$ and
$\lsixVir$ acting by left multiplication. For a list $w = [i_1,\dots,i_k]$ write
$L^{(1)}_wv_h := L_{-i_1-1}\cdots L_{-i_k-1}v_h$ and $|w| := \sum_j(i_j+1)$ (the \emph{level}).
\end{definition}

File \leanpath{VOA/Examples/InverseDSVirasoroVerma.lean}.
\begin{leancode}
def vermaGens (c h : ℂ) : Set UVir :=
  {a | (∃ n : ℤ, 0 < n ∧ a = ι ℂ (L n))
      ∨ a = ι ℂ C - algebraMap ℂ UVir c
      ∨ a = ι ℂ (L 0) - algebraMap ℂ UVir h}

noncomputable def vermaIdeal (c h : ℂ) : Submodule UVir UVir :=
  Submodule.span _ (vermaGens c h)

abbrev Verma (c h : ℂ) : Type := UVir ⧸ vermaIdeal c h

noncomputable def vermaRho (c h : ℂ) : LieHom ℂ Virasoro (End ℂ (Verma c h)) :=
  (Algebra.lsmul ℂ ℂ (Verma c h)).toLieHom.comp (ι ℂ)

noncomputable def highestVector (c h : ℂ) : Verma c h :=
  Submodule.Quotient.mk 1

def vermaMode (c h : ℂ) (m : ℤ) : End ℂ (Verma c h) :=
  vermaRho c h (L m)

def negWord (c h : ℂ) : List ℕ → Verma c h
  | [] => highestVector c h
  | i :: w => vermaMode c h (-(i : ℤ) - 1) (negWord c h w)

def energyHW (w : List ℕ) : ℕ := (w.map fun i => i + 1).sum

noncomputable def sortedSpan (c h : ℂ) (N : ℕ) : Submodule ℂ (Verma c h) :=
  Submodule.span ℂ {x | ∃ w : List ℕ,
    SortedWord w ∧ energyHW w = N ∧ negWord c h w = x}
\end{leancode}

\begin{dictionary}
\textbf{The shift is now one, not two.} In \lean{negWord c h} the entry $i$ means $L_{-i-1}$, so entry
$0$ is $L_{-1}$, which acts nontrivially on $v_h$ when $h\ne0$. Compare \lean{negWordVir}, where entry $i$
means $L_{-i-2}$. \lean{energyHW} is the level. \lean{sortedSpan c h N} is the span of sorted (weakly
decreasing) words of level $N$, i.e.\ the span of the PBW monomials indexed by partitions of $N$ into
positive parts. It is defined as a span; that it is the level-$N$ eigenspace of $L_0-h$ is proved
(\lean{verma_}\allowbreak\lean{mem_}\allowbreak\lean{level_}\allowbreak\lean{iff_}\allowbreak\lean{eigen}, \leanpath{VirasoroVermaEigenspaces.lean}).

\textbf{Same construction as $\lsixV$}, with $L_{-1}$ and $L_0-h$ in place of $L_{-1}$ and $L_0$: indeed
$\lsixV$ is a quotient of $M(c,0)$ (Theorem~\ref{l06:thm:vacexact}).
\end{dictionary}

\begin{physics}
$M(c,h)$ is the conformal family of a primary field of weight $h$ with no null state removed:
the primary state $v_h = |h\rangle$ and all its descendants.
\end{physics}

\begin{definition}[Highest-weight vectors and the universal property]\label{l06:def:hwv}
In a representation $R$ of $\lsixVir$ on $V$, a vector $v$ is a \emph{highest-weight vector of weight
$(c,h)$} if $R$ has central charge $c$, $L_nv = 0$ for $n>0$ and $L_0v = hv$. For such $v$ there is a
unique $\lsixVir$-linear map $M(c,h)\to V$ with $v_h\mapsto v$.
\end{definition}

File \leanpath{VOA/Examples/InverseDSVirasoroVerma.lean}.
\begin{leancode}
structure IsHighestWeightVector (R : VirRep V) (v : V) (c h : ℂ) : Prop where
  centralCharge : R.centralCharge = c
  pos_annihilate : ∀ n : ℤ, 0 < n → virMode R n v = 0
  zero_acts : virMode R 0 v = h • v

noncomputable def universalMap (R : VirRep V) (v : V) {c h : ℂ}
    (hv : IsHighestWeightVector R v c h) : Verma c h →ₗ[ℂ] V :=
  (((vermaIdeal c h).restrictScalars ℂ).liftQ (eval R v)
      (vermaIdeal_le_ker_eval R v hv)).comp
    (Submodule.Quotient.restrictScalarsEquiv ℂ (vermaIdeal c h)).symm.toLinearMap

theorem universalMap_intertwines (R : VirRep V) (v : V) {c h : ℂ}
    (hv : IsHighestWeightVector R v c h) (x : Virasoro) (w : Verma c h) :
    universalMap R v hv (vermaRho c h x w) =
      R.ρ x (universalMap R v hv w) := by
\end{leancode}

\begin{dictionary}
\lean{VirRep V} (Chapter~\ref{ch:L02}) is a Lie representation $\rho$ of $\lsixVir$ on $V$ with a number
\lean{centralCharge} and the proof that $\rho(C)$ is that number times the identity; \lean{virMode R n}
$= \rho(L_n)$. \lean{eval R v} is $a\mapsto a\cdot v$ on $\lsixU$; \lean{liftQ} factors it through the
quotient (the side condition \lean{vermaIdeal_le_ker_eval} is the check that the generators of $I(c,h)$
kill $v$). The remaining composition only changes the scalar ring of the quotient from $\lsixU$ to $\C$.
\lean{universalMap_}\allowbreak\lean{intertwines} says the result is a map of representations.
\end{dictionary}

\begin{physics}
Any primary state $|h\rangle$ in any theory receives a map from the Verma module: one simply lets the
descendants of the abstract primary go to the actual descendants. Null vectors of the actual theory are
the kernel of this map.
\end{physics}

\begin{theorem}[PBW basis of $M(c,h)$, for every $(c,h)$]\label{l06:thm:vermabasis}
For every $(c,h)\in\C^2$: $v_h\ne0$; the sorted monomials $L^{(1)}_wv_h$, $w$ weakly decreasing, form a
basis of $M(c,h)$; at each level $N$ the sorted words of level $N$ form a basis of the level-$N$ subspace,
whose dimension is the number $\lsixpp(N)$ of partitions of $N$.
\end{theorem}

File \leanpath{VOA/Examples/VirasoroVermaNonvanishing.lean}.
\begin{leancode}
theorem highestVector_ne_zero (c h : ℂ) : highestVector c h ≠ 0 := by
\end{leancode}
File \leanpath{VOA/Examples/VirasoroVermaMomentumPolynomial.lean}.
\begin{leancode}
theorem verma_sortedWord_linearIndependent (c h : ℂ) :
    LinearIndependent ℂ (fun w : {w : List ℕ // SortedWord w} => negWord c h w.val) := by
\end{leancode}
File \leanpath{VOA/Examples/VirasoroVermaBasis.lean}.
\begin{leancode}
noncomputable def vermaSortedBasis :
    Module.Basis {w : List ℕ // SortedWord w} ℂ (Verma c h) :=
  Module.Basis.mk (verma_sortedWord_linearIndependent c h) (verma_sortedWord_span_top c h)

def VermaPartitions (n : ℕ) := {w : List ℕ // SortedWord w ∧ energyHW w=n}

noncomputable def vermaLevelBasis (n : ℕ) :
    Module.Basis (VermaPartitions n) ℂ (sortedSpan c h n) :=
  Module.Basis.mk (vermaLevelWord_linearIndependent c h n) (vermaLevelWord_span c h n)

theorem vermaLevel_finrank (n : ℕ) :
    Module.finrank ℂ (sortedSpan c h n)=Fintype.card (VermaPartitions n) :=
\end{leancode}

\begin{dictionary}
\lean{VermaPartitions n} is the finite type of sorted words of level $n$ (entry $i$ = part $i+1$), i.e.\
partitions of $n$ into positive parts; its cardinality is $\lsixpp(n)$. The $c$ and $h$ in
\lean{vermaSortedBasis}, \lean{vermaLevelBasis}, \lean{vermaLevel_finrank} are section variables
(\lean{variable (c h : ℂ)} in the file), so they are universally quantified. Independence is proved
without \lean{VOA.pbw}; the affine field of \lean{VOA.pbw} would not apply here, although its relative field
(\lean{PBWTheorem.bottom_injective}, Chapter~\ref{ch:L02}) would give the nonvanishing of Verma highest-weight
vectors.
\end{dictionary}

\begin{physics}
The Verma character $q^{h-c/24}\prod_{n\ge1}(1-q^n)^{-1}$, for every central charge and weight.
\end{physics}

\begin{proof}[Proof on paper]
\emph{Nonvanishing}: as for the vacuum, find a concrete representation with a nonzero highest-weight vector
of weight $(c,h)$: the charged free boson of Section~\ref{l06:sec:screening} with background $\ell$ and
momentum $p$ has central charge $1-24\ell^2$ and vacuum weight $p^2-2\ell p$; choose $\ell^2 = (1-c)/24$
and $(p-\ell)^2 = h+\ell^2$ (\lean{exists_}\allowbreak\lean{chargedFock_}\allowbreak\lean{parameters}). \emph{Spanning}: straightening, as in
Theorem~\ref{l06:thm:pbw}. \emph{Independence}: map $M(c,h)$ to a family of Fock modules depending
polynomially on a momentum parameter; the coefficient of the top power of the momentum in the image of a
sorted monomial is a monomial that determines the word (\lean{vermaWordPolynomial_}\allowbreak\lean{top}), and a triangularity
argument on these leading terms gives independence (\leanpath{VirasoroVermaMomentumPolynomial.lean}).
\end{proof}

\begin{status}
\lean{VOA.}\allowbreak\lean{InverseDS.}\allowbreak\lean{VirasoroVerma.}\allowbreak\lean{highestVector_}\allowbreak\lean{ne_}\allowbreak\lean{zero},
\lean{…verma_}\allowbreak\lean{sortedWord_}\allowbreak\lean{linearIndependent}, \lean{…vermaLevel_finrank}:
\lean{[propext, Classical.}\allowbreak\lean{choice, Quot.}\allowbreak\lean{sound]}. Unconditional, \lean{VOA.pbw}-free.
\end{status}

\begin{theorem}[The vacuum module as a Verma quotient]\label{l06:thm:vacexact}
For every $c$, $L_{-1}v_0$ is a highest-weight vector of weight $(c,1)$ in $M(c,0)$, the induced map
$\iota:M(c,1)\to M(c,0)$ is injective, the map $\pi:M(c,0)\to\lsixV$, $v_0\mapsto\lsixvac$, is surjective,
and $\ker\pi = \operatorname{im}\iota$. Hence $M(c,0)/M(c,1)\cong\lsixV$.
\end{theorem}

File \leanpath{VOA/Examples/VirasoroVacuumExactSequence.lean}.
\begin{leancode}
noncomputable def translationEmbedding (c : ℂ) : Verma c 1 →ₗ[ℂ] Verma c 0 :=
  universalMap (vermaRep c 0) _ (translation_highest c)

theorem translationEmbedding_injective (c : ℂ) :
    Function.Injective (translationEmbedding c) := by

theorem vacuumProjection_surjective (c : ℂ) : Function.Surjective (vacuumProjection c) := by

theorem vacuumProjection_ker (c : ℂ) :
    LinearMap.ker (vacuumProjection c) = LinearMap.range (translationEmbedding c) := by

noncomputable def vacuumQuotientEquiv (c : ℂ) :
    (Verma c 0 ⧸ LinearMap.range (translationEmbedding c)) ≃ₗ[ℂ] VirUniv c :=
  (Submodule.quotEquivOfEq _ _ (vacuumProjection_ker c).symm).trans
    ((vacuumProjection c).quotKerEquivOfSurjective (vacuumProjection_surjective c))
\end{leancode}
File \leanpath{VOA/Examples/VirasoroVacuumProjection.lean}.
\begin{leancode}
noncomputable def vacuumProjection (c : ℂ) : Verma c 0 →ₗ[ℂ] VirUniv c :=
  universalMap (virVacRep c) (vacVec c) (vacuum_highest c)
\end{leancode}

\begin{dictionary}
\lean{translation_highest c} is the proof that $L_{-1}v_0$ is a highest-weight vector of weight $(c,1)$;
\lean{vacuum_highest c} that $\lsixvac$ is one of weight $(c,0)$ in the representation \lean{virVacRep c}
on $\lsixV$. \lean{LinearMap.ker}/\lean{LinearMap.range} are kernel and image; \lean{≃ₗ[ℂ]} a linear
isomorphism. The statement is an exact sequence $0\to M(c,1)\to M(c,0)\to\lsixV\to0$, proved in the actual
carriers, at every $c$ (including the minimal-model charges, where $M(c,0)$ has further singular vectors
that survive in $\lsixV$).
\end{dictionary}

\begin{physics}
The vacuum module is the conformal family of the identity with exactly one null vector removed:
$L_{-1}|0\rangle$, the statement that the identity operator is constant. Everything else about the identity
family at special $c$ is in the further quotient to $\lsixL$.
\end{physics}

\begin{proof}[Proof on paper]
$L_1L_{-1}v_0 = 2L_0v_0 = 0$, $L_2L_{-1}v_0 = 3L_1v_0 = 0$, and $L_0L_{-1}v_0 = L_{-1}v_0$; so $L_{-1}v_0$ is
highest of weight $1$. The images under $\iota$ of the PBW basis of $M(c,1)$ are the PBW monomials of
$M(c,0)$ ending in $L_{-1}$ (sorted words with a final entry $0$), a subfamily of a basis --- hence
injectivity. The PBW monomials of $M(c,0)$ \emph{not} ending in $L_{-1}$ map bijectively to the partition
basis of $\lsixV$, giving surjectivity and the kernel.
\end{proof}

\begin{status}
\lean{VOA.}\allowbreak\lean{InverseDS.}\allowbreak\lean{VirasoroVerma.}\allowbreak\lean{vacuumProjection_}\allowbreak\lean{ker}: \lean{[propext, Classical.}\allowbreak\lean{choice, Quot.}\allowbreak\lean{sound]}.
Unconditional.
\end{status}

%======================================================================
\section{The Shapovalov form, singular vectors and the irreducible quotient}\label{l06:sec:shap}
%======================================================================

\begin{definition}[Contravariant form]\label{l06:def:shap}
The \emph{contravariant} (Shapovalov) form on $M(c,h)$ is the bilinear form $\lsixform\cdot\cdot$ with
$\lsixform{v_h}{v_h} = 1$ and $\lsixform{L_mu}{v} = \lsixform{u}{L_{-m}v}$ for all $m\in\Z$.
\end{definition}

File \leanpath{VOA/Examples/VirasoroVermaContravariant.lean}.
\begin{leancode}
noncomputable def vermaHighestCoordinate (c h : ℂ) : Module.Dual ℂ (Verma c h) :=
  (Finsupp.lapply ⟨[], sortedWord_nil⟩).comp (vermaSortedBasis c h).repr.toLinearMap

noncomputable def vermaContravariant (c h : ℂ) :
    Verma c h →ₗ[ℂ] Module.Dual ℂ (Verma c h) :=
  universalMap (vermaDualRep c h) (vermaHighestCoordinate c h)
    (vermaHighestCoordinate_isHighest c h)

theorem vermaContravariant_adj (c h : ℂ) (m : ℤ) (u v : Verma c h) :
    vermaContravariant c h (vermaMode c h m u) v =
      vermaContravariant c h u (vermaMode c h (-m) v) := by

theorem vermaContravariant_symm (c h : ℂ) (u v : Verma c h) :
    vermaContravariant c h u v = vermaContravariant c h v u := by

theorem vermaContravariant_unique (c h : ℂ)
    (B : Verma c h →ₗ[ℂ] Module.Dual ℂ (Verma c h))
    (hnorm : B (highestVector c h) (highestVector c h) = 1)
    (hadj : ∀ m u v, B (vermaMode c h m u) v = B u (vermaMode c h (-m) v)) :
    B = vermaContravariant c h := by

theorem vermaContravariant_translation_norm (c h : ℂ) :
    vermaContravariant c h (negWord c h [0]) (negWord c h [0]) = 2*h := by
\end{leancode}

\begin{dictionary}
\textbf{The construction.} \lean{Module.Dual ℂ M} is the algebraic dual $M^* = \mathrm{Hom}(M,\C)$; a
bilinear form is a linear map $M\to M^*$. \lean{vermaHighestCoordinate} is the functional ``coefficient of
$v_h$ in the PBW basis''. \lean{vermaDualRep} makes $M^*$ a representation by $(L_m f)(v) = f(L_{-m}v)$
(\lean{vermaDualMode}; reversing indices and composition preserves \eqref{l06:eq:vir}). The highest
coordinate is a highest-weight vector of weight $(c,h)$ in $M^*$, so the universal property gives a
$\lsixVir$-map $M\to M^*$; that is the form. There is \textbf{no complex conjugation}: this is the
bilinear Shapovalov form, not a Hermitian form, so no positivity is involved.

\textbf{Pinning it down.} \lean{vermaContravariant_}\allowbreak\lean{unique}: any linear $B:M\to M^*$ with
$B(v_h)(v_h) = 1$ and adjointness equals this one. \lean{vermaContravariant_adj}: adjointness.
\lean{vermaContravariant_}\allowbreak\lean{symm}: symmetry. \lean{vermaContravariant_}\allowbreak\lean{translation_}\allowbreak\lean{norm}: the smoke test
$\lsixform{L_{-1}v_h}{L_{-1}v_h} = 2h$. So this is exactly Definition~\ref{l06:def:shap}.
\end{dictionary}

\begin{physics}
$\lsixform{u}{v}$ is the two-point function $\langle u|v\rangle$ of descendants of a primary with
$L_n^\dagger = L_{-n}$. The norm $2h$ of $L_{-1}|h\rangle$ is the statement that $\partial\phi$ has two-point
function $\propto h$: for $h = 0$ the field is constant and its derivative is null.
\end{physics}

\begin{definition}[Gram matrix, radical, irreducible quotient]\label{l06:def:gramverma}
$G_n(c,h)$ is the matrix of $\lsixform\cdot\cdot$ on the level-$n$ subspace in the partition basis. The
\emph{radical} $\mathrm{Rad}(c,h)$ is $\{u : \lsixform{u}{\cdot} = 0\}$, and $L(c,h) := M(c,h)/\mathrm{Rad}(c,h)$.
\end{definition}

File \leanpath{VOA/Examples/VirasoroVermaGramPolynomial.lean}.
\begin{leancode}
noncomputable def vermaGramMatrix (c h : ℂ) (n : ℕ) :
    Matrix (VermaPartitions n) (VermaPartitions n) ℂ := fun i j =>
  vermaContravariant c h (negWord c h i.val) (negWord c h j.val)

noncomputable def vermaGramPolynomial (n : ℕ) :
    Matrix (VermaPartitions n) (VermaPartitions n) (MvPolynomial (Fin 2) ℂ) :=
  vermaAlgebraGramMatrix (MvPolynomial.X 0) (MvPolynomial.X 1) n

theorem vermaGramPolynomial_det_eval (c h : ℂ) (n : ℕ) [DecidableEq (VermaPartitions n)] :
    MvPolynomial.aeval (fun i : Fin 2 => if i=0 then c else h)
      (vermaGramPolynomial n).det = (vermaGramMatrix c h n).det := by
\end{leancode}
File \leanpath{VOA/Examples/VirasoroVermaMaximalRadical.lean}.
\begin{leancode}
noncomputable def vermaRadical (c h : ℂ) : Submodule ℂ (Verma c h) :=
  LinearMap.ker (vermaContravariant c h)

abbrev VermaSimple (c h : ℂ) := Verma c h ⧸ vermaRadical c h
\end{leancode}

\begin{dictionary}
\lean{vermaGramMatrix c h n} is literally the matrix of the form on the partition words. 
\lean{vermaGramPolynomial n} is a matrix over $\C[X_0,X_1]$ ($X_0$ = central charge, $X_1$ = highest
weight) computed by the finite commutation algorithm; \lean{vermaGramPolynomial_}\allowbreak\lean{det_}\allowbreak\lean{eval} proves that
its determinant, evaluated at $(c,h)$, is the actual determinant --- so $\det G_n$ is a polynomial in
$(c,h)$, proved, not postulated. \lean{vermaRadical} is the kernel of $u\mapsto\lsixform{u}{\cdot}$.
\end{dictionary}

\begin{physics}
The Gram matrix collects all two-point functions of level-$n$ descendants. The radical is the space of
null states; $L(c,h)$ is the physical conformal family.
\end{physics}

\begin{theorem}[Radical, irreducibility, dimensions]\label{l06:thm:radical}
For every $(c,h)$: $\mathrm{Rad}(c,h)$ is the unique maximal proper submodule of $M(c,h)$; $L(c,h)$ is
irreducible and nonzero; $\dim L(c,h)_n = \operatorname{rank}G_n(c,h)$; and $L(c,0)$ (Verma version) is
linearly isomorphic to the simple VOA $\lsixL$ of Definition~\ref{l06:def:simple}.
\end{theorem}

File \leanpath{VOA/Examples/VirasoroVermaMaximalRadical.lean}.
\begin{leancode}
theorem vermaRadical_maximal (c h : ℂ) (S : Submodule ℂ (Verma c h))
    (hS : VermaModeStable c h S) (hr : vermaRadical c h≤S) :
    S=vermaRadical c h ∨ S=⊤ := by

theorem vermaSimple_irreducible (c h : ℂ) (S : Submodule ℂ (VermaSimple c h))
    (hS : ∀ m : ℤ, ∀ v∈S, vermaSimpleMode c h m v∈S) : S=⊥ ∨ S=⊤ := by

theorem vermaSimple_highest_ne_zero (c h : ℂ) :
    (vermaRadical c h).mkQ (highestVector c h)≠0 := by
\end{leancode}
File \leanpath{VOA/Examples/VirasoroVermaQuotientDimensions.lean}.
\begin{leancode}
noncomputable def vermaSimpleLevel (c h : ℂ) (n : ℕ) : Submodule ℂ (VermaSimple c h) :=
  (sortedSpan c h n).map (vermaRadical c h).mkQ

theorem vermaSimpleLevel_finrank_eq_gramRank (c h : ℂ) (n : ℕ) :
    Module.finrank ℂ (vermaSimpleLevel c h n)=(vermaGramMatrix c h n).rank := by
\end{leancode}
File \leanpath{VOA/Examples/VirasoroVermaSimpleVacuum.lean}.
\begin{leancode}
noncomputable def vermaSimpleVacuumEquiv (c : ℂ) :
    VermaSimple c 0 ≃ₗ[ℂ] (VirUniv c ⧸ maxIdealVir c) :=
  (Submodule.quotEquivOfEq _ _ (simpleVacuumProjection_ker c).symm).trans
    ((simpleVacuumProjection c).quotKerEquivOfSurjective (simpleVacuumProjection_surjective c))
\end{leancode}

\begin{dictionary}
\lean{VermaModeStable c h S} says $S$ is stable under every $L_m$ (and $\lsixVir$-submodule);
\lean{vermaSimpleMode} is the induced action on the quotient. \lean{vermaRadical_maximal} is maximality,
\lean{vermaSimple_}\allowbreak\lean{irreducible} irreducibility of the quotient (no nonzero proper invariant subspace), and
\lean{vermaSimple_}\allowbreak\lean{highest_}\allowbreak\lean{ne_}\allowbreak\lean{zero} says the quotient is nonzero. (\lean{verma_}\allowbreak\lean{proper_}\allowbreak\lean{submodule_}\allowbreak\lean{le_}\allowbreak\lean{radical},
in the same file, gives ``every proper submodule lies in the radical''.) The isomorphism
\lean{vermaSimpleVacuumEquiv} is only a \emph{linear} isomorphism, built from the kernel identity
\lean{simpleVacuumProjection_}\allowbreak\lean{ker}: the composite $M(c,0)\to\lsixV\to\lsixL$ has kernel exactly
$\mathrm{Rad}(c,0)$; \lean{vermaSimpleVacuumEquiv_}\allowbreak\lean{mode} records that it intertwines the modes.
\end{dictionary}

\begin{physics}
$L(c,h)$ is the irreducible conformal family; its character is $\sum_n\operatorname{rank}G_n(c,h)\,q^{h+n-c/24}$.
For $h = 0$ it is the vacuum module of the chiral algebra $\lsixL$.
\end{physics}

\begin{proof}[Proof on paper]
A proper submodule $S$ is graded and does not contain $v_h$, so it lies in $\bigoplus_{n>0}M_n$; for
$u\in S$ and any $v$, $\lsixform{u}{L_{-m_1}\cdots v_h} = \lsixform{L_{m_k}\cdots L_{m_1}u}{v_h}$, which is the
$v_h$-coefficient of an element of $S$, hence $0$. So $S\subseteq\mathrm{Rad}$; the radical is itself
proper ($\lsixform{v_h}{v_h} = 1$) and a submodule (adjointness). Levels are orthogonal
(\lean{vermaContravariant_}\allowbreak\lean{orthogonal}), so the radical at level $n$ is the kernel of $G_n$, giving the rank
formula. For $h = 0$, both $\mathrm{Rad}(c,0)$ and the preimage of $I_{\max}(c)$ are the maximal proper
submodule.
\end{proof}

\begin{status}
\lean{VOA.}\allowbreak\lean{InverseDS.}\allowbreak\lean{VirasoroVerma.}\allowbreak\lean{vermaSimple_}\allowbreak\lean{irreducible}, \lean{…vermaSimpleLevel_}\allowbreak\lean{finrank_}\allowbreak\lean{eq_}\allowbreak\lean{gramRank},
\lean{…vermaSimpleVacuumEquiv}, \lean{…vermaContravariant_}\allowbreak\lean{unique}, \lean{…vermaContravariant_}\allowbreak\lean{symm}:
\lean{[propext, Classical.}\allowbreak\lean{choice, Quot.}\allowbreak\lean{sound]}. Unconditional.
\end{status}

\subsection{Singular vectors}

\begin{definition}[Singular vector]\label{l06:def:singular}
A \emph{singular vector} at level $n>0$ in $M(c,h)$ is a nonzero $v$ in the level-$n$ subspace with
$L_1v = L_2v = 0$ (equivalently $L_mv = 0$ for all $m>0$, since $L_1,L_2$ generate the positive part).
\end{definition}

File \leanpath{VOA/Examples/VirasoroKacCollisionStates.lean}.
\begin{leancode}
def HasNonzeroSingular (c h : ℂ) (n : ℕ) : Prop :=
  ∃ v : Verma c h, v≠0 ∧ v∈sortedSpan c h n ∧
    vermaMode c h 1 v=0 ∧ vermaMode c h 2 v=0
\end{leancode}
File \leanpath{VOA/Examples/VirasoroVermaWardSpecialization.lean}.
\begin{leancode}
theorem verma_positive_zero_of_one_two (c h : ℂ) (v : Verma c h)
    (h1 : vermaMode c h 1 v=0) (h2 : vermaMode c h 2 v=0) :
    ∀ n : ℤ, 0<n → vermaMode c h n v=0 := by
\end{leancode}

\begin{dictionary}
Many statements of the library write the singular condition inline as ``$v\ne0$, $v\in$
\lean{sortedSpan c h n}, \lean{vermaMode c h 1 v = 0}, \lean{vermaMode c h 2 v = 0}'' (sometimes with $v$
an element of the subtype \lean{sortedSpan c h n}). \lean{verma_}\allowbreak\lean{positive_}\allowbreak\lean{zero_}\allowbreak\lean{of_}\allowbreak\lean{one_}\allowbreak\lean{two} justifies using
only $L_1,L_2$: $[L_2,L_1] = L_3$, $[L_{m},L_1] = (m-1)L_{m+1}$ generate all $L_m$, $m>0$.
\end{dictionary}

\begin{physics}
A singular vector is a descendant that is itself primary: a null field. Its existence is the condition
under which BPZ differential equations hold.
\end{physics}

\begin{theorem}[Singular vectors versus the Gram determinant]\label{l06:thm:singgram}
Let $n>0$.
\begin{enumerate}
\item If $M(c,h)$ has a singular vector at level $n$, then $\det G_n(c,h) = 0$; conversely if
  $\det G_n(c,h)\ne0$ there is none.
\item If $\det G_n(c,h) = 0$ and $\det G_k(c,h)\ne0$ for all $k<n$, then $M(c,h)$ has a singular vector at
  level $n$.
\item A singular vector $v$ at level $n$ generates an injective image of $M(c,h+n)$; consequently
  $(X-h)^{\lsixpp(m)}$ divides $\det G_{n+m}(c,X)$ as a polynomial in the highest weight $X$, for every
  $m\ge0$.
\item Singular vectors at a given level are unique up to scalar.
\end{enumerate}
\end{theorem}

File \leanpath{VOA/Examples/VirasoroVermaGramPolynomial.lean}.
\begin{leancode}
theorem vermaGramMatrix_nonsingular_no_singular (c h : ℂ) (n : ℕ)
    [DecidableEq (VermaPartitions n)] (hn : 0 < n)
    (hd : (vermaGramMatrix c h n).det ≠ 0)
    (v : Verma c h) (hv : v ∈ sortedSpan c h n)
    (h1 : vermaMode c h 1 v = 0) (h2 : vermaMode c h 2 v = 0) : v = 0 := by

theorem verma_singular_of_first_gram_root (c h : ℂ) (n : ℕ)
    [∀ k : ℕ, DecidableEq (VermaPartitions k)]
    (hroot : (vermaGramMatrix c h n).det = 0)
    (hlower : ∀ k < n, (vermaGramMatrix c h k).det ≠ 0) :
    ∃ v : Verma c h, v ≠ 0 ∧ v ∈ sortedSpan c h n ∧
      ∀ m : ℤ, 0 < m → vermaMode c h m v = 0 := by
\end{leancode}
File \leanpath{VOA/Examples/VirasoroVermaSingularEmbedding.lean}.
\begin{leancode}
noncomputable def vermaSingularMap (c h : ℂ) (n : ℕ) (v : Verma c h)
    (hv : v ∈ sortedSpan c h n)
    (h1 : vermaMode c h 1 v=0) (h2 : vermaMode c h 2 v=0) :
    Verma c (h+(n:ℂ)) →ₗ[ℂ] Verma c h :=
  universalMap (vermaRep c h) v (verma_singular_isHighest c h n v hv h1 h2)

theorem vermaSingularMap_injective (c h : ℂ) (n : ℕ) (v : Verma c h)
    (hv : v ∈ sortedSpan c h n) (hne : v≠0)
    (h1 : vermaMode c h 1 v=0) (h2 : vermaMode c h 2 v=0) :
    Function.Injective (vermaSingularMap c h n v hv h1 h2) :=

noncomputable def vermaWeightGramPolynomial (c : ℂ) (n : ℕ) :
    Matrix (VermaPartitions n) (VermaPartitions n) (Polynomial ℂ) :=
  vermaAlgebraGramMatrix (Polynomial.C c) Polynomial.X n

theorem vermaWeightDeterminant_dvd_of_singular (c h : ℂ) (n m : ℕ)
    [DecidableEq (VermaPartitions (n+m))] (hn : 0<n)
    (v : sortedSpan c h n) (hv : v≠0)
    (h1 : vermaMode c h 1 (v : Verma c h)=0)
    (h2 : vermaMode c h 2 (v : Verma c h)=0) :
    (Polynomial.X-Polynomial.C h)^(Fintype.card (VermaPartitions m)) ∣
      (vermaWeightGramPolynomial c (n+m)).det := by
\end{leancode}
File \leanpath{VOA/Examples/VirasoroSingularUniqueness.lean}.
\begin{leancode}
theorem verma_singular_eq_smul (c h : ℂ) (n : ℕ) (v w : Verma c h)
    (hv : v ∈ sortedSpan c h n) (hw : w ∈ sortedSpan c h n) (hwne : w ≠ 0)
    (hv1 : vermaMode c h 1 v = 0) (hv2 : vermaMode c h 2 v = 0)
    (hw1 : vermaMode c h 1 w = 0) (hw2 : vermaMode c h 2 w = 0) :
    v = (vermaTopCoordinate c h n v / vermaTopCoordinate c h n w) • w := by
\end{leancode}

\begin{dictionary}
\lean{vermaWeightGramPolynomial c n} is the Gram matrix at fixed $c$ as a matrix of polynomials in the
highest weight $X$ (\lean{Polynomial.C c} is the constant $c$); \lean{vermaWeightGramPolynomial_}\allowbreak\lean{eval}
shows that evaluating at $h$ gives \lean{vermaGramMatrix c h n}. \lean{∣} is divisibility of polynomials.
\lean{[DecidableEq …]} arguments are bookkeeping required to form a determinant and carry no
mathematical content. \lean{vermaTopCoordinate c h n v} is the coefficient of $L_{-1}^nv_h$ in $v$;
uniqueness is proved by showing it is nonzero on every singular vector (\lean{verma_}\allowbreak\lean{singular_}\allowbreak\lean{top_}\allowbreak\lean{ne_}\allowbreak\lean{zero}).
In item 2 the hypothesis \lean{hlower} is needed: without it a zero of $\det G_n$ may be caused only by
descendants of lower singular vectors.
\end{dictionary}

\begin{physics}
Item 3 is the source of the exponents $\lsixpp(N-rs)$ in Kac's formula: a null vector at level $rs$ drags
along its whole conformal family of null descendants, $\lsixpp(m)$ of them at level $rs+m$. Item 4 is the
classical uniqueness of the null vector at given level (used by BPZ to write explicit null-vector
equations).
\end{physics}

\begin{proof}[Proof on paper]
(1) A singular vector lies in the radical (it is orthogonal to everything, as moving creation modes across
turns them into annihilation modes); its coefficient vector is a nonzero kernel vector of $G_n$. If
$\det G_n\ne0$ the kernel is zero. (2) A kernel vector $u$ at level $n$ is in the radical; then $L_mu$,
$m>0$, is in the radical at level $n-m<n$, where the form is nondegenerate, so $L_mu = 0$. (3) The universal
map is injective because a nonzero vector of a Verma module has linearly independent PBW descendants
(\lean{vermaDescendant_}\allowbreak\lean{linearIndependent}: $\lsixU^-$ acts freely). The $\lsixpp(m)$ independent
descendants at level $n+m$ lie in the radical of $M(c,h)$. Now use an elementary fact about polynomial
matrices: if $A(X)$ is a square polynomial matrix and $k$ linearly independent vectors lie in the kernel
of $A(h)$, then $(X-h)^k$ divides $\det A(X)$ (complete them to a basis; the corresponding columns of
$A(X)P$ vanish at $h$, so each is divisible by $X-h$). (4) Two singular vectors at level $n$ with the same
top coordinate differ by a singular vector with zero top coordinate, which must be zero.
\end{proof}

\begin{status}
All items: \lean{[propext, Classical.}\allowbreak\lean{choice, Quot.}\allowbreak\lean{sound]}. Unconditional.
\end{status}

%======================================================================
\section{The Kac determinant}\label{l06:sec:kac}
%======================================================================

\subsection{The Kac parametrisation and the Kac product}

\begin{definition}[Kac parametrisation]\label{l06:def:kacparam}
For $t\in\C\setminus\{0\}$ and $r,s\in\C$ put $c(t) := 13-6t-6/t$ and
$h_{r,s}(t) := \bigl((rt-s)^2-(t-1)^2\bigr)/(4t)$.
\end{definition}

File \leanpath{VOA/Examples/VirasoroVermaBSA.lean}.
\begin{leancode}
noncomputable def vermaBSACharge (t : ℂ) : ℂ := 13-6*t-6/t
\end{leancode}
File \leanpath{VOA/Examples/VirasoroFusionLabelRoots.lean}.
\begin{leancode}
noncomputable def fusionKacWeight (u v t : ℂ) : ℂ :=
  ((u*t-v)^2-(t-1)^2)/(4*t)
\end{leancode}

\begin{dictionary}
The names are historical (\lean{vermaBSACharge}: the charge on the Benoit--Saint-Aubin curve of
Chapter~\ref{ch:L07}; \lean{fusionKacWeight}: introduced for the fusion construction) but the formulas are
exactly \eqref{l06:eq:kacparam}. Lean's total division gives $6/0 = 0$, so at $t = 0$ the value $c(0) = 13$
is meaningless; every theorem below assumes $t\ne0$. \textbf{Surjectivity}: every $c\in\C$ is $c(t)$ for
some $t\ne0$ (the quadratic $6t^2-(13-c)t+6 = 0$ has roots with product $1$), so statements ``for all
$t\ne0$'' cover all central charges. \textbf{Orientation}: at $t = p/q$, $c(p/q) = c_{p,q}$ and
$h_{r,s}(p/q) = ((pr-qs)^2-(p-q)^2)/(4pq)$ (\lean{fusionKacWeight_}\allowbreak\lean{minimal},
\lean{minimalKacWeight}); with $p>q$ the vacuum label at the first singular level is $(r,s) = (q-1,p-1)$.
Under $t\mapsto1/t$, $c$ is invariant and $h_{r,s}\leftrightarrow h_{s,r}$ (\lean{vermaBSACharge_inv},
\lean{fusionKacWeight_}\allowbreak\lean{swap_}\allowbreak\lean{inv}).
\end{dictionary}

\begin{physics}
In the Coulomb gas $t = \alpha_+^2$ (up to normalisation) and $h_{r,s}$ is the weight of the vertex
operator with charge $\alpha_{r,s} = \frac{1-r}2\alpha_++\frac{1-s}2\alpha_-$. The symmetry $t\leftrightarrow1/t$
exchanges the two screening charges.
\end{physics}

\begin{definition}[The Kac product]\label{l06:def:kacproduct}
For $t\ne0$ and $n\in\N$,
\begin{equation*}
 \mathcal K_n(t;X) := \prod_{\substack{r,s\ge1\\ rs\le n}}\bigl(X-h_{r,s}(t)\bigr)^{\lsixpp(n-rs)}\in\C[X],
 \qquad
 K_n := \prod_{\lambda\vdash n}\ \prod_{k\ge1}(2k)^{m_k(\lambda)}m_k(\lambda)!,
\end{equation*}
where $m_k(\lambda)$ is the multiplicity of the part $k$ in the partition $\lambda$.
\end{definition}

File \leanpath{VOA/Examples/VirasoroKacProduct.lean}.
\begin{leancode}
noncomputable def kacLabels (n : ℕ) : Finset (Fin n × Fin n) :=
  Finset.univ.filter (fun i => (i.1.val+1)*(i.2.val+1) ≤ n)

noncomputable def kacWeightProduct (t : ℂ) (n : ℕ) : Polynomial ℂ :=
  ∏ i ∈ kacLabels n,
    (Polynomial.X-Polynomial.C (fusionKacWeight (i.1.val+1) (i.2.val+1) t)) ^
      Fintype.card (VermaPartitions (n-(i.1.val+1)*(i.2.val+1)))
\end{leancode}
File \leanpath{VOA/Examples/VirasoroVermaGramLeading.lean}.
\begin{leancode}
noncomputable def vermaDiagonalContraction : List ℕ → ℂ
  | [] => 1
  | a::w => (((w.count a+1:ℕ):ℂ)*(2*((a:ℂ)+1)))*vermaDiagonalContraction w
\end{leancode}

\begin{dictionary}
\textbf{Labels.} \lean{Fin n × Fin n} indexes pairs $(i,j)$ with $0\le i,j<n$; the label is
$(r,s) = (i+1,j+1)$, and the filter keeps $rs\le n$. Since $rs\le n$ forces $r,s\le n$, no label is lost by
the bound $<n$ on the indices. \textbf{Exponents.} \lean{Fintype.}\allowbreak\lean{card (VermaPartitions m)} is $\lsixpp(m)$
(partitions into positive parts; \lean{n-…} is natural-number subtraction, harmless since $rs\le n$).
\textbf{Constant.} For a sorted word $w$ (entry $a$ = part $a+1$) \lean{vermaDiagonalContraction w}
multiplies, for each entry, $2(a+1)$ (i.e.\ $2k$ for the part $k$) by the number of later equal entries plus
one; over a block of $m$ equal parts this gives $(2k)^mm!$. So
$\prod_{w}$\lean{vermaDiagonalContraction w} $= K_n$. Example: $K_2 = 4\cdot(2^2\cdot2!) = 32$, and indeed
$\det G_2 = 32h\,(h-h_{1,2})(h-h_{2,1})$.
\end{dictionary}

\begin{physics}
$\mathcal K_n$ lists the expected null vectors at level $\le n$, each with its family of descendants;
$K_n$ is the large-$h$ normalisation, in which each $L_{-k}$ contracted with $L_k$ contributes $2kh$.
\end{physics}

\subsection{What the library proves about the determinant in general}

\begin{theorem}[Degree, leading coefficient and the degree count]\label{l06:thm:degree}
For every $c\in\C$ and $n\in\N$, $\det G_n(c,X)\in\C[X]$ has degree exactly $\sum_{\lambda\vdash n}\ell(\lambda)$
and leading coefficient $K_n$ (independent of $c$). Moreover
$\sum_{\lambda\vdash n}\ell(\lambda) = \sum_{rs\le n}\lsixpp(n-rs)$, so $\mathcal K_n(t;X)$ is monic of the same
degree.
\end{theorem}

File \leanpath{VOA/Examples/VirasoroVermaGramLeading.lean}.
\begin{leancode}
theorem vermaWeightDeterminant_natDegree (c : ℂ) (n : ℕ) :
    (vermaWeightGramPolynomial c n).det.natDegree=∑ j : VermaPartitions n, j.val.length := by

theorem vermaWeightDeterminant_leadingCoeff (c : ℂ) (n : ℕ) :
    (vermaWeightGramPolynomial c n).det.leadingCoeff=
      ∏ j : VermaPartitions n, vermaDiagonalContraction j.val := by
\end{leancode}
File \leanpath{VOA/Examples/VirasoroKacDegreeCount.lean}.
\begin{leancode}
theorem verma_partition_lengths_eq_kac_count (n : ℕ) :
    (∑ w : VermaPartitions n, w.val.length)=
      ∑ a ∈ Finset.range n, ∑ k ∈ Finset.range n,
        if (a+1)*(k+1)≤n then
          Fintype.card (VermaPartitions (n-(a+1)*(k+1))) else 0 := by
\end{leancode}
File \leanpath{VOA/Examples/VirasoroKacProduct.lean}.
\begin{leancode}
theorem kacWeightProduct_monic (t : ℂ) (n : ℕ) : (kacWeightProduct t n).Monic := by

theorem kacWeightProduct_natDegree (t : ℂ) (n : ℕ) :
    (kacWeightProduct t n).natDegree = ∑ w : VermaPartitions n, w.val.length := by
\end{leancode}

\begin{dictionary}
\lean{natDegree}, \lean{leadingCoeff}, \lean{Monic} are Mathlib's degree, leading coefficient and
``leading coefficient $1$''. \lean{w.val.length} is the number of parts $\ell(\lambda)$.
\lean{Finset.range n} is $\{0,\dots,n-1\}$. The combinatorial identity is the classical one: every
partition with a marked part of size $a+1$ occurring at least $k+1$ times corresponds to a partition of
$n-(a+1)(k+1)$, so both sides count pairs (partition, marked part-occurrence).
\end{dictionary}

\begin{physics}
At large $h$ the descendants of a primary behave like free oscillators and the Gram matrix becomes
diagonal: $\lsixform{L_\lambda v_h}{L_\lambda v_h}\approx\prod_k(2kh)^{m_k}m_k!$. The degree identity
says the null vectors predicted by Kac account for exactly the total number of ``oscillator
excitations'' --- no room is left for unexpected zeros.
\end{physics}

\begin{proof}[Proof on paper]
Each entry $\lsixform{L_\lambda v_h}{L_\mu v_h}$ has $h$-degree $\le\min(\ell(\lambda),\ell(\mu))$ (each power of
$h$ comes from one contraction $[L_k,L_{-k}]\ni2kL_0$), so $\deg\det\le\sum\ell(\lambda)$, with coefficient
the determinant of the matrix of top coefficients. That matrix vanishes between partitions of different
content (contractions must pair equal modes), is diagonal on equal content, and its diagonal entries are
$\prod_k(2k)^{m_k}m_k!\ne0$; it is independent of $c$ (\lean{vermaWeightTopMatrix_}\allowbreak\lean{charge_}\allowbreak\lean{independent}).
\end{proof}

\begin{status}
\lean{VOA.}\allowbreak\lean{InverseDS.}\allowbreak\lean{VirasoroVerma.}\allowbreak\lean{vermaWeightDeterminant_}\allowbreak\lean{natDegree}, \lean{…leadingCoeff},
\lean{…verma_}\allowbreak\lean{partition_}\allowbreak\lean{lengths_}\allowbreak\lean{eq_}\allowbreak\lean{kac_}\allowbreak\lean{count}, \lean{…kacWeightProduct_}\allowbreak\lean{natDegree}:
\lean{[propext, Classical.}\allowbreak\lean{choice, Quot.}\allowbreak\lean{sound]}. Unconditional.
\end{status}

\begin{theorem}[Kac formula from singular vectors (conditional form)]\label{l06:thm:kaccond}
Let $c\in\C$, $n\in\N$ and $t\in\C$ with $\operatorname{Im}t\ne0$. Suppose that for every label $(r,s)$ with
$rs\le n$ the Verma module $M(c,h_{r,s}(t))$ has a singular vector at level $rs$. Then
\begin{equation*}
 \det G_n(c,X) = K_n\cdot\mathcal K_n(t;X).
\end{equation*}
\end{theorem}

File \leanpath{VOA/Examples/VirasoroKacProduct.lean}.
\begin{leancode}
theorem fusionKacWeight_injective_positive (t : ℂ) (ht : t.im ≠ 0)
    (r s r' s' : ℕ) (hr : 0 < r) (hr' : 0 < r')
    (he : fusionKacWeight r s t = fusionKacWeight r' s' t) : r = r' ∧ s = s' := by

theorem vermaWeightDeterminant_eq_kac_of_singular (c t : ℂ) (n : ℕ) (ht : t.im ≠ 0)
    (hs : ∀ i ∈ kacLabels n,
      ∃ v : sortedSpan c (fusionKacWeight (i.1.val+1) (i.2.val+1) t)
        ((i.1.val+1)*(i.2.val+1)), v ≠ 0 ∧
        vermaMode c (fusionKacWeight (i.1.val+1) (i.2.val+1) t) 1 v = 0 ∧
        vermaMode c (fusionKacWeight (i.1.val+1) (i.2.val+1) t) 2 v = 0) :
    (vermaWeightGramPolynomial c n).det =
      Polynomial.C (∏ w : VermaPartitions n, vermaDiagonalContraction w.val) *
        kacWeightProduct t n := by
\end{leancode}

\begin{dictionary}
\textbf{Hypotheses.} \lean{ht : t.im ≠ 0}: $t$ is not real. This guarantees that the weights $h_{r,s}(t)$ of
distinct positive labels are distinct (\lean{fusionKacWeight_}\allowbreak\lean{injective_}\allowbreak\lean{positive}), so that the factors
$X-h_{r,s}$ are pairwise coprime. \lean{hs}: for every label in \lean{kacLabels n}, a nonzero
$v\in$\lean{sortedSpan c h_{r,s} (rs)} with $L_1v = L_2v = 0$: a singular vector. \textbf{Note} that $c$
and $t$ are independent variables; the hypothesis \lean{hs} can only hold for all labels when $c$ is
compatible with $t$ (for $n\ge2$ it forces $c = c(t)$), but the theorem does not need to say so.
\textbf{Conclusion}: an equality of polynomials in $X$ (\lean{Polynomial.C} is the constant polynomial).

\textbf{Status of the hypothesis.} The docstring says ``The displayed singular existence premise remains
open in general, so this is not an unconditional Kac or rigidity theorem.'' When this chapter was first written, this was already outdated (the docstring has since been corrected): the premise is discharged, for $c = c(t)$ and every nonreal $t$, by
\lean{screening_}\allowbreak\lean{nonreal_}\allowbreak\lean{kac_}\allowbreak\lean{singular} (Theorem~\ref{l06:thm:screenkac}); see
Theorem~\ref{l06:thm:kacfull} and Warning~\ref{l06:warn:docstrings}. The theorem itself is correct as stated.
\end{dictionary}

\begin{physics}
Each primary of weight $h_{r,s}$ with a null vector at level $rs$ forces the factor
$(h-h_{r,s})^{\lsixpp(n-rs)}$ in the level-$n$ determinant, by Theorem~\ref{l06:thm:singgram}(3). Kac's
insight is that these are all the zeros: the counting of Theorem~\ref{l06:thm:degree} leaves no room for more.
\end{physics}

\begin{proof}[Proof on paper]
By Theorem~\ref{l06:thm:singgram}(3), each $(X-h_{r,s}(t))^{\lsixpp(n-rs)}$ divides $D(X) := \det G_n(c,X)$. For
nonreal $t$ the $h_{r,s}(t)$ are distinct: if $h_{r,s} = h_{r',s'}$ then $(rt-s)^2 = (r't-s')^2$, so
$(r-r')t = s-s'$ or $(r+r')t = s+s'$; taking imaginary parts, $r = r'$ (then $s = s'$) or $r+r' = 0$,
impossible. Pairwise coprime factors each dividing $D$ give $\mathcal K_n\mid D$. Both have the same degree
(Theorem~\ref{l06:thm:degree}), $\mathcal K_n$ is monic and $D$ has leading coefficient $K_n$; hence
$D = K_n\mathcal K_n$.
\end{proof}

\begin{status}
\lean{VOA.}\allowbreak\lean{InverseDS.}\allowbreak\lean{VirasoroVerma.}\allowbreak\lean{vermaWeightDeterminant_}\allowbreak\lean{eq_}\allowbreak\lean{kac_}\allowbreak\lean{of_}\allowbreak\lean{singular}:
\lean{[propext, Classical.}\allowbreak\lean{choice, Quot.}\allowbreak\lean{sound]}. Conditional on the displayed hypothesis \lean{hs}, which is
dischargeable (Theorem~\ref{l06:thm:kacfull}).
\end{status}

\begin{theorem}[Extension from nonreal to all nonzero $t$]\label{l06:thm:kacextend}
Fix $N$. If for every nonreal $z$ and every label with $rs\le N$, $M(c(z),h_{r,s}(z))$ has a singular vector
at level $rs$, then for every $t\ne0$ and every $k\le N$, $\det G_k(c(t),X) = K_k\,\mathcal K_k(t;X)$.
\end{theorem}

File \leanpath{VOA/Examples/VirasoroKacVacuumCoefficient.lean}.
\begin{leancode}
theorem vermaWeightDeterminant_kac_extend (n : ℕ)
    (hs : ∀ z : ℂ, z.im ≠ 0 →
      (vermaWeightGramPolynomial (vermaBSACharge z) n).det =
        Polynomial.C (∏ w : VermaPartitions n, vermaDiagonalContraction w.val) *
          kacWeightProduct z n)
    (t : ℂ) (ht : t ≠ 0) :
    (vermaWeightGramPolynomial (vermaBSACharge t) n).det =
      Polynomial.C (∏ w : VermaPartitions n, vermaDiagonalContraction w.val) *
        kacWeightProduct t n := by
\end{leancode}
File \leanpath{VOA/Examples/VirasoroKacSingularTransfer.lean}.
\begin{leancode}
theorem verma_kac_all_of_nonreal_singular (N : ℕ)
    (hs : ∀ z : ℂ, z.im ≠ 0 → ∀ r s : ℕ, 0 < r → 0 < s → r*s ≤ N →
      ∃ v : sortedSpan (vermaBSACharge z) (fusionKacWeight r s z) (r*s),
        v ≠ 0 ∧ vermaMode (vermaBSACharge z) (fusionKacWeight r s z) 1 v = 0 ∧
          vermaMode (vermaBSACharge z) (fusionKacWeight r s z) 2 v = 0)
    (t : ℂ) (ht : t ≠ 0) (k : ℕ) (hk : k ≤ N) :
    (vermaWeightGramPolynomial (vermaBSACharge t) k).det =
      Polynomial.C (∏ w : VermaPartitions k,vermaDiagonalContraction w.val) *
        kacWeightProduct t k := by
\end{leancode}

\begin{dictionary}
In both statements the central charge is now tied to the parameter: \lean{vermaBSACharge z} $= c(z)$.
The conclusion holds at \emph{every} nonzero $t$, including real $t$ where labels collide (e.g.\ $t = p/q$,
the minimal models): the identity of polynomials is extended by continuity in $t$ of both sides (the Gram
entries are polynomials in $c(t)$ and $X$, continuous in $t\ne0$), not by re-running the coprimality
argument, which would fail at collisions. The hypothesis of \lean{verma_}\allowbreak\lean{kac_}\allowbreak\lean{all_}\allowbreak\lean{of_}\allowbreak\lean{nonreal_}\allowbreak\lean{singular} is
stated with natural-number labels $r,s>0$ instead of \lean{kacLabels}; the proof converts.
\end{dictionary}

\begin{physics}
The Kac formula holds uniformly, including at the rational values of $t$ where the Kac table
degenerates and weights coincide (e.g.\ $h_{r,s} = h_{q-r,p-s}$ in the minimal models) --- there the factors
simply merge, and the exponents add.
\end{physics}

\begin{proof}[Proof on paper]
For nonreal $z$ Theorem~\ref{l06:thm:kaccond} applies at every level $k\le N$. The coefficients of both
sides, as polynomials in $X$, are continuous functions of $z$ on $\C\setminus\{0\}$; nonreal numbers are
dense; so the coefficients agree at every $t\ne0$.
\end{proof}

\begin{status}
\lean{…vermaWeightDeterminant_}\allowbreak\lean{kac_}\allowbreak\lean{extend}, \lean{…verma_}\allowbreak\lean{kac_}\allowbreak\lean{all_}\allowbreak\lean{of_}\allowbreak\lean{nonreal_}\allowbreak\lean{singular}:
\lean{[propext, Classical.}\allowbreak\lean{choice, Quot.}\allowbreak\lean{sound]}. Conditional on \lean{hs}, discharged in
Theorem~\ref{l06:thm:kacfull}.
\end{status}

\subsection{Roots imply singular vectors; Laurent interpolation}

\begin{theorem}[Bootstrap from roots]\label{l06:thm:bootstrap}
Let $\operatorname{Im}t\ne0$ and $N\in\N$. If $\det G_{rs}(c(t),h_{r,s}(t)) = 0$ for every label with $rs\le N$,
then for every such label $M(c(t),h_{r,s}(t))$ has a singular vector at level $rs$.
\end{theorem}

File \leanpath{VOA/Examples/VirasoroKacRootBootstrap.lean}.
\begin{leancode}
theorem verma_singular_of_nonreal_kac_roots (t : ℂ) (ht : t.im ≠ 0) (N : ℕ)
    (hroots : ∀ r s : ℕ, 0 < r → 0 < s → r*s ≤ N →
      (vermaGramMatrix (vermaBSACharge t) (fusionKacWeight r s t) (r*s)).det = 0) :
    ∀ r s : ℕ, 0 < r → 0 < s → r*s ≤ N →
      ∃ v : sortedSpan (vermaBSACharge t) (fusionKacWeight r s t) (r*s),
        v ≠ 0 ∧ vermaMode (vermaBSACharge t) (fusionKacWeight r s t) 1 v = 0 ∧
          vermaMode (vermaBSACharge t) (fusionKacWeight r s t) 2 v = 0 := by
\end{leancode}

\begin{dictionary}
The hypothesis is only that $\det G_{rs}$ \emph{vanishes} at the first expected level; nothing about
lower levels is assumed. The conclusion is the hypothesis \lean{hs} of Theorem~\ref{l06:thm:kaccond} for
$c = c(t)$.
\end{dictionary}

\begin{physics}
A zero of the determinant at the first level where the label $(r,s)$ can contribute is a genuine null
\emph{primary} (not a descendant of an earlier null vector), because at nonreal $t$ no other null vector
lives below.
\end{physics}

\begin{proof}[Proof on paper]
Strong induction on $rs$. For a label $(r,s)$ and $k<rs$, the labels $(r',s')$ with $r's'\le k$ have
singular vectors by induction, so Theorem~\ref{l06:thm:kaccond} gives $\det G_k(c(t),h) = K_k\mathcal K_k(t;h)$
at $h = h_{r,s}(t)$; this is nonzero because $h_{r,s}(t)\ne h_{r',s'}(t)$ for the (different) labels with
$r's'\le k<rs$. Hence $\det G_{rs} = 0$ is the first degenerate level, and
Theorem~\ref{l06:thm:singgram}(2) produces a singular vector.
\end{proof}

\begin{status}
\lean{…verma_}\allowbreak\lean{singular_}\allowbreak\lean{of_}\allowbreak\lean{nonreal_}\allowbreak\lean{kac_}\allowbreak\lean{roots}: \lean{[propext, Classical.}\allowbreak\lean{choice, Quot.}\allowbreak\lean{sound]}. Conditional on
\lean{hroots}, discharged in Theorem~\ref{l06:thm:screenkac}.
\end{status}

\begin{theorem}[Infinitely many roots on a Kac curve suffice]\label{l06:thm:laurent}
Fix $r,s,n$. The function $t\mapsto\det G_n(c(t),h_{r,s}(t))$ is a Laurent polynomial in $t$. If it vanishes on
an infinite set of nonzero $t$, it vanishes at every $t\ne0$.
\end{theorem}

File \leanpath{VOA/Examples/VirasoroKacLaurent.lean}.
\begin{leancode}
theorem kacGram_root_of_infinite (r s n : ℕ) (S : Set ℂ) (hS : S.Infinite)
    (hne : ∀ t ∈ S, t ≠ 0)
    (hz : ∀ t ∈ S, (vermaGramMatrix (vermaBSACharge t)
      (fusionKacWeight r s t) n).det = 0) (t : ℂ) (ht : t ≠ 0) :
    (vermaGramMatrix (vermaBSACharge t) (fusionKacWeight r s t) n).det = 0 := by
\end{leancode}

\begin{dictionary}
\lean{S.Infinite}: the set $S$ is infinite. No density or accumulation point is required; that is the
point of using Laurent polynomials (\lean{ℂ[T;T⁻¹]}, Mathlib's \lean{LaurentPolynomial}): the
determinant is the polynomial \lean{vermaGramPolynomial n} evaluated at the Laurent polynomials
$13-6T-6T^{-1}$ and $\frac14T^{-1}((rT-s)^2-(T-1)^2)$ (\lean{kacGramLaurent}, \lean{kacGramLaurent_eval}).
\end{dictionary}

\begin{physics}
Algebraic continuation in the Coulomb-gas parameter: a statement verified at the special values where the
screening integrals are finite sums propagates to all central charges.
\end{physics}

\begin{proof}[Proof on paper]
Multiply the Laurent polynomial by $T^m$ to clear negative powers; an ordinary polynomial with infinitely
many roots is zero, and $T^m$ is a unit.
\end{proof}

\begin{status}
\lean{…kacGram_}\allowbreak\lean{root_}\allowbreak\lean{of_}\allowbreak\lean{infinite}: \lean{[propext, Classical.}\allowbreak\lean{choice, Quot.}\allowbreak\lean{sound]}. Unconditional.
\end{status}

%======================================================================
\section{Screening operators and the Kac roots}\label{l06:sec:screening}
%======================================================================

\subsection{The charged free boson}

\begin{definition}[Charged Fock representation]\label{l06:def:chargedfock}
Let $\mathcal F = \C[x_1,x_2,\dots]$ be the rank-one Fock space with oscillators $\alpha_n$,
$[\alpha_m,\alpha_n] = 2m\,\delta_{m+n,0}$ (a lattice vector of norm $2$). For a background $\ell\in\C$ and a
momentum $p\in\C$ define
\begin{equation*}
 L^{\ell,p}_n := \tfrac14\sum_j{:}\alpha_{n-j}\alpha_j{:}+p\,\alpha_n+p^2\delta_{n,0}
 -(n+1)\ell\bigl(\alpha_n+2p\,\delta_{n,0}\bigr).
\end{equation*}
These satisfy \eqref{l06:eq:vir} with $c = 1-24\ell^2$, and the Fock vacuum $1$ is a highest-weight vector
of weight $h = p^2-2\ell p$.
\end{definition}

File \leanpath{VOA/Examples/InverseDSLeeYangFockVirasoro.lean}.
\begin{leancode}
def alphaMode (n : ℤ) : End ℂ LeeYangFock := HeisenbergN.oscN 1 2 0 n

def sectorCurrent (p : ℂ) (n : ℤ) : End ℂ LeeYangFock :=
  alphaMode n + (if n = 0 then 2 * p else 0) • (1 : End ℂ LeeYangFock)

def sectorSugawara (p : ℂ) (n : ℤ) : End ℂ LeeYangFock :=
  HeisenbergN.LvirN 1 2 n + p • alphaMode n
    + (if n = 0 then p ^ 2 else 0) • (1 : End ℂ LeeYangFock)

def sectorBackgroundCurrent (ell p : ℂ) (n : ℤ) : End ℂ LeeYangFock :=
  ell • sectorCurrent p n

def sectorVirasoro (ell p : ℂ) (n : ℤ) : End ℂ LeeYangFock :=
  sectorSugawara p n - ((n : ℂ) + 1) • sectorBackgroundCurrent ell p n

theorem sectorVirasoro_virasoro (ell p : ℂ) (m n : ℤ) :
    sectorVirasoro ell p m * sectorVirasoro ell p n
        - sectorVirasoro ell p n * sectorVirasoro ell p m =
      ((m : ℂ) - n) • sectorVirasoro ell p (m + n)
        + (if m + n = 0 then
            ((m : ℂ) ^ 3 - m) / 12 * (1 - 24 * ell ^ 2) else 0) •
          (1 : End ℂ LeeYangFock) := by
\end{leancode}
File \leanpath{VOA/Examples/VirasoroVermaNonvanishing.lean}.
\begin{leancode}
noncomputable def chargedFockRep (ell p : ℂ) : VirRep LeeYangFock where
  ρ := fixedRho (chargedFockData ell p) rfl
  centralCharge := 1-24*ell^2
  acts_central := fixedRho_C (chargedFockData ell p) rfl

theorem exists_chargedFock_parameters (c h : ℂ) :
    ∃ ell p : ℂ, 1-24*ell^2=c ∧ p^2-2*ell*p=h := by
\end{leancode}

\begin{dictionary}
\lean{LeeYangFock} is (an abbreviation for) \lean{FockN 1} $= \C[x_{(k,0)}]$, the rank-one Fock space of
Chapter~\ref{ch:L05}; the name comes from its first use in the Lee--Yang inverse Drinfeld--Sokolov work
of Chapter~\ref{ch:L19}. \lean{HeisenbergN.}\allowbreak\lean{oscN 1 2 0 n} is the oscillator at level $2$, i.e.\
$[\alpha_m,\alpha_n] = 2m\delta$; \lean{LvirN 1 2 n} $= \frac14\sum{:}\alpha\alpha{:}$ is its Sugawara mode.
\lean{sectorCurrent p} shifts the zero mode by the momentum: $\alpha_0 = 2p$ on the sector of momentum
$p\alpha$. \lean{chargedFockRep} packages the modes as a \lean{VirRep} via \lean{fixedRho}, the
extension from mode relations to the Lie algebra. \lean{exists_}\allowbreak\lean{chargedFock_}\allowbreak\lean{parameters}: every $(c,h)$ is
realised.
\end{dictionary}

\begin{physics}
This is the Coulomb gas: a free boson $\varphi$ with background charge $\ell$ (stress tensor
$T = \frac14(\partial\varphi)^2-\ell\,\partial^2\varphi$ in this normalisation) and a vertex operator of
momentum $p$, whose conformal weight $p^2-2\ell p$ is the familiar $\alpha(\alpha-2\alpha_0)$.
\end{physics}

\begin{definition}[Integer screening state]\label{l06:def:screenstate}
Let $\beta\in\Q\setminus\{0\}$ with $t := \beta^2\in\N$, and $r,s\ge1$. Put
$\ell := \frac12(\beta-\beta^{-1})$ and $p := \ell+\frac12(r\beta-s\beta^{-1})$. The \emph{screening state}
$S_{r,s}\in\mathcal F$ is the homogeneous polynomial of degree $rs$ (energy $rs$) whose coefficient on the
oscillator monomial $\prod_k x_k^{d_k}$ is
\begin{equation*}
 \frac{\beta^{\sum d_k}}{\prod_k k^{d_k}d_k!}\cdot
 \operatorname{coeff}_{\,z_1^{(r-1)t+s}\cdots z_r^{(r-1)t+s}}\Bigl[\,\prod_{i<j}(z_i-z_j)^{2t}\cdot
 \prod_k p_k(z)^{d_k}\Bigr],
\end{equation*}
where $p_k(z) = z_1^k+\dots+z_r^k$ are the power sums.
\end{definition}

File \leanpath{VOA/Examples/VirasoroScreeningSingular.lean}.
\begin{leancode}
def screeningBackground (β : ℚ) : ℚ := (β - 1/β)/2

def screeningMomentum (r s : ℕ) (β : ℚ) : ℚ :=
  screeningBackground β + ((r:ℚ)*β - (s:ℚ)/β)/2
\end{leancode}
File \leanpath{VOA/Examples/VirasoroScreeningMoment.lean}.
\begin{leancode}
noncomputable def screeningMoment (r s t : ℕ) (P : MvPolynomial (Fin r) ℚ) : ℚ :=
  (vandermondeLaurent r^(2*t) * auxiliaryLaurent r P)
    (fun _ : Fin r => (((r-1)*t+s:ℕ):ℤ))
\end{leancode}
File \leanpath{VOA/Examples/VirasoroScreeningFock.lean}.
\begin{leancode}
noncomputable def screeningPrefactor (β : ℚ) (d : ℕ →₀ ℕ) : ℚ :=
  β ^ (d.sum fun _ n => n) /
    ∏ i ∈ d.support, ((i+1 : ℕ) : ℚ) ^ d i * (Nat.factorial (d i) : ℚ)

noncomputable def screeningFockCoefficient (r s t : ℕ) (β : ℚ) (d : ℕ →₀ ℕ) : ℚ :=
  screeningPrefactor β d * screeningMoment r s t
    (powerSumSubstitution (Fin r) (MvPolynomial.monomial d 1))

noncomputable def screeningFockState (r s t : ℕ) (β : ℚ) : FockN 1 :=
  ∑ d ∈ (weightDegSetN_finite 1 (r * s)).toFinset,
    MvPolynomial.monomial d (screeningFockCoefficient r s t β (plainOccupation d) : ℂ)
\end{leancode}

\begin{dictionary}
\textbf{Decoding.} \lean{d : ℕ →₀ ℕ} is an occupation: $d(i)$ = number of oscillators of mode $i+1$ (so
$x_k^{d_k}$ with $k = i+1$; the shift by one is why \lean{(i+1)} appears in \lean{screeningPrefactor}).
\lean{powerSumSubstitution} sends the formal variable $i$ to the power sum $p_{i+1}$ in $r$ variables;
\lean{auxiliaryLaurent} embeds polynomials into Laurent polynomials; \lean{vandermondeLaurent r} is
$\prod_{i<j}(z_i-z_j)$, raised to $2t$; evaluating a Laurent polynomial (a finitely supported function on
exponent vectors) at the constant vector $((r-1)t+s,\dots)$ extracts that coefficient. The sum in
\lean{screeningFockState} runs over the finite set of occupations of energy $rs$; \lean{plainOccupation}
forgets the trivial colour index of \lean{FockN 1}. So the definition is the displayed formula; in
particular it is a finite sum with rational coefficients, with no infinite series and no contour integral.

\textbf{Parameters.} $c = 1-24\ell^2 = 13-6\beta^2-6\beta^{-2} = c(\beta^2)$ and
$h = p^2-2\ell p = h_{r,s}(\beta^2)$ (\lean{screening_}\allowbreak\lean{central_}\allowbreak\lean{charge}, \lean{screening_}\allowbreak\lean{highest_}\allowbreak\lean{weight}): the
Kac parametrisation with $t = \beta^2$. The library uses the screening only at $\beta = m\in\N_{>0}$, so
$t = m^2$; the variable \lean{t} of \lean{screeningFockState} is the exponent and is set to $m^2$.
\end{dictionary}

\begin{physics}
$S_{r,s}$ is the Dotsenko--Fateev/Felder state $Q^r|p_0\rangle$: $r$ screening currents $e^{\beta\varphi}$
integrated around the origin. The Vandermonde power is the OPE $\prod(z_i-z_j)^{2\beta^2}$ of the screening
currents, the exponential $\exp(\beta\sum\alpha_{-k}p_k(z)/k)$ is their oscillator part (whence the
prefactor $\beta^{\sum d}/\prod k^{d_k}d_k!$), and the contour integral picks the residue --- here, since
$2\beta^2$ is an even integer, a single coefficient.
\end{physics}

\begin{theorem}[Integer screening states are nonzero singular vectors]\label{l06:thm:screensing}
For all $r,s,m\ge1$ (with $\beta = m$, $t = m^2$): $S_{r,s}\ne0$, $S_{r,s}$ has energy $rs$, and
$L^{\ell,p}_1S_{r,s} = L^{\ell,p}_2S_{r,s} = 0$.
\end{theorem}

File \leanpath{VOA/Examples/VirasoroScreeningSingular.lean}.
\begin{leancode}
theorem integer_screening_nonzero_singular (r s m : ℕ) (hr : 0 < r) (hm : 0 < m) :
    screeningFockState r s (m^2) (m:ℚ) ≠ 0 ∧
    screeningFockState r s (m^2) (m:ℚ) ∈ gradingWN 1 (r*s) ∧
    sectorVirasoro (screeningBackground (m:ℚ) : ℂ)
      (screeningMomentum r s (m:ℚ) : ℂ) 1 (screeningFockState r s (m^2) (m:ℚ)) = 0 ∧
    sectorVirasoro (screeningBackground (m:ℚ) : ℂ)
      (screeningMomentum r s (m:ℚ) : ℂ) 2 (screeningFockState r s (m^2) (m:ℚ)) = 0 := by
\end{leancode}

\begin{dictionary}
\lean{gradingWN 1 (r*s)} is the energy-$rs$ subspace of the Fock space. No hypothesis on $s$ is needed
($s = 0$ is allowed by the statement and harmless). The two operators are the \emph{actual} modes
\lean{sectorVirasoro} of Definition~\ref{l06:def:chargedfock}, not a model of them.
\end{dictionary}

\begin{physics}
The screening charge commutes with the Virasoro algebra, so it carries the highest-weight vacuum of one
Fock module to a highest-weight vector of another: a null vector of the Coulomb gas at level $rs$ and
weight $h_{r,s}$.
\end{physics}

\begin{proof}[Proof on paper]
\emph{Ward identities.} On the coefficients of $S_{r,s}$, the modes $L_1,L_2$ act as the differential
operators $\sum_i z_i^{n+1}\partial_{z_i}$ ($n = 1,2$) plus multiplications by power sums, applied to the
integrand $\Delta^{2t}\prod p_k^{d_k}$; at the chosen momentum the result is a total derivative
$\sum_i\partial_{z_i}(z_i^{n+1}\cdots)$ minus terms that cancel. The finite substitute for ``the contour
integral of a total derivative vanishes'' is a coefficient identity: at the common exponent $D$, the
coefficient of
$\sum_iz_i^{n+1}\partial_iF$ equals $(D-n)$ times the coefficient of $p_n\cdot F$
(\lean{screening_}\allowbreak\lean{coefficient_}\allowbreak\lean{integration}). Combined with the product rule for $\Delta^{2t}$
(\lean{screeningVectorField_}\allowbreak\lean{one_}\allowbreak\lean{vandermonde}), these give two exact linear relations among the residues,
which are exactly the vanishing of the $L_1$- and $L_2$-coefficients (\lean{screeningFockState_}\allowbreak\lean{L1_}\allowbreak\lean{zero},
\lean{screeningFockState_}\allowbreak\lean{L2_}\allowbreak\lean{zero}).

\emph{Nonvanishing.} By Newton's identities, $(z_1\cdots z_r)^s$ is a polynomial in power sums of weighted
degree $rs$; its residue is the coefficient of $\prod z_i^{(r-1)t}$ in $\Delta^{2t}$. That central
coefficient is nonzero: with $P := \Delta^t$, reversing all exponents gives
$P(1/z) = \pm P(z)/\prod_iz_i^{(r-1)t}$ (each factor $z_i^{-1}-z_j^{-1} = (z_j-z_i)/(z_iz_j)$), so the
coefficient in question is $\pm$ the constant term of $P(z)P(1/z)$, which is the sum of the squares of the
rational coefficients of $P$, hence nonzero (\lean{vandermonde_}\allowbreak\lean{central_}\allowbreak\lean{coefficient_}\allowbreak\lean{ne_}\allowbreak\lean{zero}).
By linearity some single monomial has nonzero residue (\lean{screening_}\allowbreak\lean{nonzero_}\allowbreak\lean{monomial}); the prefactor is
nonzero for $\beta\ne0$.
\end{proof}

\begin{status}
\lean{VOA.}\allowbreak\lean{VirasoroScreening.}\allowbreak\lean{integer_}\allowbreak\lean{screening_}\allowbreak\lean{nonzero_}\allowbreak\lean{singular},
\lean{VOA.}\allowbreak\lean{VirasoroScreening.}\allowbreak\lean{screeningFockState_}\allowbreak\lean{ne_}\allowbreak\lean{zero}: \lean{[propext, Classical.}\allowbreak\lean{choice, Quot.}\allowbreak\lean{sound]}.
Unconditional. No Dyson/Selberg product formula is used.
\end{status}

\subsection{From Fock singular vectors to Verma Gram roots}

\begin{theorem}[Transfer]\label{l06:thm:transfer}
Let $N>0$ and $P\in\mathcal F$ be a nonzero vector of energy $N$ with $L^{\ell,p}_1P = L^{\ell,p}_2P = 0$.
Then $\det G_N(1-24\ell^2,\,p^2-2\ell p) = 0$.
\end{theorem}

File \leanpath{VOA/Examples/VirasoroFockGramTransfer.lean}.
\begin{leancode}
theorem chargedFock_singular_gram_root (ell p : ℂ) (N : ℕ) (hN : 0 < N)
    (P : FockN 1) (hP : P ∈ gradingWN 1 N) (hne : P ≠ 0)
    (h1 : sectorVirasoro ell p 1 P = 0) (h2 : sectorVirasoro ell p 2 P = 0) :
    (vermaGramMatrix (1-24*ell^2) (p^2-2*ell*p) N).det = 0 := by
\end{leancode}
File \leanpath{VOA/Examples/VirasoroFockLevelDimension.lean}.
\begin{leancode}
theorem fock_verma_level_finrank (c h : ℂ) (N : ℕ) :
    Module.finrank ℂ (gradingWN 1 N) = Module.finrank ℂ (sortedSpan c h N) := by
\end{leancode}

\begin{dictionary}
The conclusion is about the actual Verma Gram matrix, at the central charge and weight of the Fock
module. Faithfulness of the free-field realisation is neither assumed nor claimed (the universal map
$M\to\mathcal F$ is in general neither injective nor surjective).
\end{dictionary}

\begin{physics}
A null vector in the Coulomb gas forces a zero of the Kac determinant at the same $(c,h)$: if the
abstract conformal family had no null vector at level $N$, the free-field realisation would be
isomorphic to it up to level $N$, and could not contain one either.
\end{physics}

\begin{proof}[Proof on paper]
Suppose $\det G_N\ne0$. Then the radical of $M(c,h)$ vanishes at every level $k\le N$ (a nonzero radical
vector at level $k$ would have a nonzero descendant at level $N$ in the radical, by
\lean{verma_}\allowbreak\lean{radical_}\allowbreak\lean{zero_}\allowbreak\lean{below}). The kernel of the universal map $M(c,h)\to\mathcal F$ is a proper submodule,
hence in the radical, hence zero at levels $\le N$. Levels of $M$ and $\mathcal F$ have the same dimension
$\lsixpp(k)$ (\lean{fock_}\allowbreak\lean{verma_}\allowbreak\lean{level_}\allowbreak\lean{finrank}), so the map is bijective at levels $\le N$. Lift $P$ to
$u\ne0$ at level $N$; then $L_1u$, $L_2u$ map to $0$ at levels $<N$, where the map is injective, so $u$ is
singular --- contradicting $\det G_N\ne0$ (Theorem~\ref{l06:thm:singgram}(1)).
\end{proof}

\begin{status}
\lean{VOA.}\allowbreak\lean{VirasoroScreening.}\allowbreak\lean{chargedFock_}\allowbreak\lean{singular_}\allowbreak\lean{gram_}\allowbreak\lean{root}: \lean{[propext, Classical.}\allowbreak\lean{choice, Quot.}\allowbreak\lean{sound]}.
Unconditional.
\end{status}

\begin{theorem}[Kac roots and singular vectors at all labels]\label{l06:thm:screenkac}
For all $r,s\ge1$:
\begin{enumerate}
\item for every integer $m\ge1$, $\det G_{rs}(c(m^2),h_{r,s}(m^2)) = 0$;
\item hence for every $t\ne0$, $\det G_{rs}(c(t),h_{r,s}(t)) = 0$;
\item for every nonreal $t$, $M(c(t),h_{r,s}(t))$ has a singular vector at level $rs$.
\end{enumerate}
\end{theorem}

File \leanpath{VOA/Examples/VirasoroScreeningKacRoots.lean}.
\begin{leancode}
theorem integer_screening_kac_root (r s m : ℕ) (hr : 0 < r) (hs : 0 < s)
    (hm : 0 < m) :
    (vermaGramMatrix (vermaBSACharge ((m:ℂ)^2))
      (fusionKacWeight r s ((m:ℂ)^2)) (r*s)).det = 0 := by

theorem infinite_integer_screening_roots (r s : ℕ) (hr : 0 < r) (hs : 0 < s) :
    ∃ S : Set ℂ, S.Infinite ∧ (∀ t ∈ S, t ≠ 0) ∧
      ∀ t ∈ S, (vermaGramMatrix (vermaBSACharge t)
        (fusionKacWeight r s t) (r*s)).det = 0 := by

theorem screening_kac_root (r s : ℕ) (hr : 0 < r) (hs : 0 < s)
    (t : ℂ) (ht : t ≠ 0) :
    (vermaGramMatrix (vermaBSACharge t) (fusionKacWeight r s t) (r*s)).det = 0 := by

theorem screening_nonreal_kac_singular (t : ℂ) (ht : t.im ≠ 0)
    (r s : ℕ) (hr : 0 < r) (hs : 0 < s) :
    ∃ v : sortedSpan (vermaBSACharge t) (fusionKacWeight r s t) (r*s),
      v ≠ 0 ∧ vermaMode (vermaBSACharge t) (fusionKacWeight r s t) 1 v = 0 ∧
        vermaMode (vermaBSACharge t) (fusionKacWeight r s t) 2 v = 0 := by
\end{leancode}

\begin{dictionary}
The set $S$ of \lean{infinite_}\allowbreak\lean{integer_}\allowbreak\lean{screening_}\allowbreak\lean{roots} is $\{1,4,9,\dots\}$. Item 3 is \emph{exactly} the
hypothesis \lean{hs} of \lean{verma_}\allowbreak\lean{kac_}\allowbreak\lean{all_}\allowbreak\lean{of_}\allowbreak\lean{nonreal_}\allowbreak\lean{singular} (for any $N$), and its specialisation is
the hypothesis of \lean{vermaWeightDeterminant_}\allowbreak\lean{eq_}\allowbreak\lean{kac_}\allowbreak\lean{of_}\allowbreak\lean{singular} with $c = c(t)$.
\end{dictionary}

\begin{physics}
Every entry of the (infinite) Kac table gives a zero of $\det G_{rs}$ at every central charge, and a
singular vector at level $rs$ at every nonreal $t$ --- proved by building the null vectors in the Coulomb gas at
the special values $t = 1,4,9,\dots$ and continuing algebraically. At real $t$, which includes all minimal
models, the theorem gives only the vanishing of the determinant.
\end{physics}

\begin{proof}[Proof on paper]
(1) Theorems~\ref{l06:thm:screensing} and~\ref{l06:thm:transfer} with $\beta = m$. (2) Theorem~\ref{l06:thm:laurent}
with $S = \{m^2\}$. (3) Theorem~\ref{l06:thm:bootstrap}, whose hypothesis is (2).
\end{proof}

\begin{status}
All four: \lean{[propext, Classical.}\allowbreak\lean{choice, Quot.}\allowbreak\lean{sound]}. Unconditional.
\end{status}

\subsection{The Kac determinant formula, unconditionally}\label{l06:sec:kacunconditional}

\begin{theorem}[Kac determinant formula]\label{l06:thm:kacfull}
For every $t\in\C\setminus\{0\}$, every $n\in\N$ and every $h\in\C$,
\begin{equation*}
 \det G_n\bigl(c(t),h\bigr) = K_n\prod_{\substack{r,s\ge1\\ rs\le n}}\bigl(h-h_{r,s}(t)\bigr)^{\lsixpp(n-rs)} .
\end{equation*}
Since every $c\in\C$ is some $c(t)$, this is the Kac determinant for all Virasoro Verma modules.
\end{theorem}

This theorem is \textbf{not} a declaration of the library. It is the composition of
Theorem~\ref{l06:thm:kacextend} and Theorem~\ref{l06:thm:screenkac}(3): applying
\lean{verma_kac_all_of_nonreal_singular} to the singular vectors \lean{screening_nonreal_kac_singular} gives the
identity of polynomials in $h$, and evaluating it (\lean{vermaWeightDeterminant_eval}) gives the displayed formula.
We performed this composition in Lean against the library, as two theorems \lean{kac_determinant} and
\lean{kac_determinant_eval}; the check is not part of this release.

\begin{dictionary}
The identity is stated with symbols defined above: \lean{vermaWeightGramPolynomial} (Theorem~\ref{l06:thm:singgram}),
\lean{vermaBSACharge} (Definition~\ref{l06:def:kacparam}), \lean{vermaDiagonalContraction} and
\lean{kacWeightProduct} (Definition~\ref{l06:def:kacproduct}). The only hypothesis is $t\ne0$. The
check was compiled against an earlier commit, when this chapter was first written (not re-run for this edition), and \texttt{\#print axioms} reports
\lean{[propext, Classical.}\allowbreak\lean{choice, Quot.}\allowbreak\lean{sound]} for both theorems. Why it is not trivially true: the left side
is the determinant of the actual Shapovalov form, computed by an independent algorithm, and the right side
is an explicit product; the degree and leading coefficient match by Theorem~\ref{l06:thm:degree}, so any
error would show up as a wrong root.
\end{dictionary}

\begin{physics}
This is Kac's theorem (1979), proved by Feigin--Fuchs (1982--84) via the Coulomb gas: it identifies all the
central charges and weights at which conformal families contain null states, and is the basis of the
classification of minimal models.
\end{physics}

\begin{proof}[Proof on paper]
Combine: singular vectors at nonreal $t$ (Theorem~\ref{l06:thm:screenkac}) give the formula at nonreal $t$
(Theorem~\ref{l06:thm:kaccond}), which extends to all $t\ne0$ by continuity (Theorem~\ref{l06:thm:kacextend}).
\end{proof}

\begin{status}
Checked in Lean against the library (not a library declaration; the check is not part of this release), axioms
\lean{[propext, Classical.}\allowbreak\lean{choice, Quot.}\allowbreak\lean{sound]}. \textbf{Unconditional}. The library additionally contains
older partial results that this subsumes: the formula at levels $<9$ (\lean{vermaWeightDeterminant_}\allowbreak\lean{kac_}\allowbreak\lean{below_}\allowbreak\lean{nine_}\allowbreak\lean{all},
via Benoit--Saint-Aubin vectors for labels with a part $\le2$), singular vectors at $t = \pm1$ for all
labels (\lean{kacSingular_at_one}, \lean{kacSingular_at_neg_one}), and the Benoit--Saint-Aubin factors
$(X-h_{1,N})^{\lsixpp(m)}$ for all $t\ne0$ (\lean{vermaWeightDeterminant_}\allowbreak\lean{bsa_}\allowbreak\lean{factor_}\allowbreak\lean{all}).
\end{status}

%======================================================================
\section{From Verma modules back to the vacuum: the simple zero at \texorpdfstring{$c_{p,q}$}{c(p,q)}}\label{l06:sec:vacuumkac}
%======================================================================

The graded (first-order) rigidity theorem of Chapter~\ref{ch:L07} needs one fact about the \emph{vacuum} Gram determinant
$D_n(c) := \det G^{\mathrm{vac}}_n(c)$ (Definition~\ref{l06:def:partgram}): at $c = c_{p,q}$ and the first
singular level $n = (p-1)(q-1)$, $D_n$ has a \emph{simple} zero. This section explains how the library
derives it from the Verma determinant.

\begin{theorem}[Vacuum block of the Verma determinant]\label{l06:thm:vacblock}
For every $c$ and $n$:
\begin{enumerate}
\item the vacuum Gram matrix $G^{\mathrm{vac}}_n(c)$ is the submatrix of the Verma Gram matrix $G_n(c,0)$ on
  the partitions without parts equal to $1$;
\item the coefficient of $X^{\lsixpp(n)}$ in $\det G_{n+1}(c,X)$ equals
  $D_{n+1}(c)\cdot\det T_n(c)$, where $T_n(c)$ is the first-order (in $X$) block of the Gram matrix on the
  words ending with $L_{-1}$.
\end{enumerate}
\end{theorem}

File \leanpath{VOA/Examples/VirasoroVermaVacuumForm.lean}.
\begin{leancode}
theorem vermaGramMatrix_vacuum_submatrix (c : ℂ) (n : ℕ) :
    (vermaGramMatrix c 0 n).submatrix (vacuumVermaPartition n) (vacuumVermaPartition n) =
      (partitionGramPolynomial n).map (Polynomial.eval c) := by
\end{leancode}
File \leanpath{VOA/Examples/VirasoroVacuumDeterminantBlock.lean}.
\begin{leancode}
noncomputable def translationFirstGram (c : ℂ) (n : ℕ) :
    Matrix (VermaPartitions n) (VermaPartitions n) ℂ := fun i j =>
  (vermaWeightGramPolynomial c (n+1)
    (translationVermaPartition n i) (translationVermaPartition n j)).coeff 1

theorem vermaWeightDeterminant_vacuum_first (c : ℂ) (n : ℕ) :
    (vermaWeightGramPolynomial c (n+1)).det.coeff (Fintype.card (VermaPartitions n)) =
      ((partitionGramPolynomial (n+1)).map (Polynomial.eval c)).det *
        (translationFirstGram c n).det := by
\end{leancode}

\begin{dictionary}
\lean{vacuumVermaPartition n} sends a vacuum partition (entries $i$ meaning $L_{-i-2}$) to the same
modes in the Verma convention (entries $i+1$ meaning $L_{-i-2}$); \lean{submatrix} restricts rows and
columns. \lean{translationVermaPartition n} appends an $L_{-1}$ to a level-$n$ word. \lean{.coeff k} is the
coefficient of $X^k$. Since $\det G_{n+1}(c,X)$ is divisible by $X^{\lsixpp(n)}$ (the factor of the label
$(1,1)$), item 2 identifies its first possibly nonzero coefficient.
\end{dictionary}

\begin{physics}
At $h\to0$ the descendants $L_{-1}(\dots)$ become null (they are descendants of the null vector
$L_{-1}|0\rangle$); the leading behaviour of the determinant as $h\to0$ factors into the norm of the
genuine vacuum states times the norm of these ``translation'' states, which behaves like $h$ times the
Gram matrix of $M(c,1)$.
\end{physics}

\begin{proof}[Proof on paper]
Order the partition basis of level $n+1$ as (words without $L_{-1}$) $\sqcup$ (words ending in $L_{-1}$).
Entries pairing a word ending in $L_{-1}$ with anything are divisible by $X$ (move the $L_{-1}$ across as
$L_1$, which acts on $v_h$-descendants producing a factor that vanishes with $h$); so the matrix has the
block form $\begin{pmatrix}A(X) & XB(X)\\ XB'(X) & XT(X)\end{pmatrix}$ and the coefficient of $X^{\lsixpp(n)}$
in its determinant is $\det A(0)\det T(0)$, with $A(0) = G^{\mathrm{vac}}_{n+1}(c)$ by item 1.
\end{proof}

\begin{status}
Both items: \lean{[propext, Classical.}\allowbreak\lean{choice, Quot.}\allowbreak\lean{sound]}. Unconditional.
\end{status}

\begin{theorem}[Simple zero of the vacuum determinant at the minimal models]\label{l06:thm:vacsimple}
Let $p>q\ge2$ be coprime and $N = (p-1)(q-1)$. Then $D_N(c_{p,q}) = 0$ and $D_N'(c_{p,q})\ne0$.
\end{theorem}

File \leanpath{VOA/Examples/VirasoroKacVacuumTransfer.lean}.
\begin{leancode}
theorem vacuum_simple_of_verma_kac (p q n : ℕ) (hq : 2 ≤ q)
    (hpq : q < p) (hcop : p.Coprime q) (hn : n+1 = (p-1)*(q-1))
    (hk : ∀ z : ℂ, z ≠ 0 → ∀ k ≤ n+1,
      (vermaWeightGramPolynomial (vermaBSACharge z) k).det =
        Polynomial.C (∏ w : VermaPartitions k, vermaDiagonalContraction w.val) *
          kacWeightProduct z k) :
    let P := (VirDeformModes.NormalOrdering.partitionGramPolynomial (n+1)).det
    P.eval (vermaBSACharge ((p:ℂ)/q)) = 0 ∧
      P.derivative.eval (vermaBSACharge ((p:ℂ)/q)) ≠ 0 := by
\end{leancode}
File \leanpath{VOA/Examples/VirasoroKacScalarRoot.lean}.
\begin{leancode}
theorem kacVacuumScalar_minimal_simple (p q n : ℕ) (hq : 2 ≤ q)
    (hpq : q < p) (hcop : p.Coprime q) (hn : n+1 = (p-1)*(q-1)) :
    kacVacuumScalar ((p:ℂ)/q) n = 0 ∧
      _root_.deriv (fun z => kacVacuumScalar z n) ((p:ℂ)/q) ≠ 0 := by
\end{leancode}
File \leanpath{VOA/Examples/VirasoroKacVacuumCoefficient.lean}.
\begin{leancode}
noncomputable def kacVacuumScalar (t : ℂ) (n : ℕ) : ℂ :=
  ∏ i ∈ (kacLabels (n+1)).erase (0,0),
    (-fusionKacWeight (i.1.val+1) (i.2.val+1) t) ^
      Fintype.card (VermaPartitions (n+1-(i.1.val+1)*(i.2.val+1)))
\end{leancode}

\begin{dictionary}
\textbf{Hypotheses.} \lean{hq}, \lean{hpq}, \lean{hcop}: $2\le q<p$, coprime --- the minimal-model labels with
the order fixed (the symmetric case $p<q$ follows by $c_{p,q} = c_{q,p}$). \lean{hn}: $n+1 = N$.
\lean{hk}: the Kac formula at every $z\ne0$ and every level $\le N$ --- which is
Theorem~\ref{l06:thm:kacfull}, hence dischargeable. \textbf{Conclusion}: $P = D_N$ as a polynomial in $c$;
\lean{P.derivative} is its derivative; evaluated at \lean{vermaBSACharge (p/q)} $= c_{p,q}$.
\lean{kacVacuumScalar t n} is $\mathcal K_N(t;0)$ with the factor of the label $(1,1)$ removed (the
\lean{.erase (0,0)}); \lean{_root_.deriv} is the complex derivative.

\textbf{Docstrings.} Both docstrings say the Kac identity ``remains the missing ordinary
representation-theory input'' / ``still requires the missing general ordinary Kac representation input''.
Outdated (Warning~\ref{l06:warn:docstrings}).

\textbf{Unconditional form.} Composing \lean{vacuum_simple_of_verma_kac} with the Kac formula
(Theorem~\ref{l06:thm:kacfull}, which supplies \lean{hk}) gives the theorem without \lean{hk}; this composition,
\lean{vacuum_simple_zero}, was checked in Lean against the library (the check is not part of this release),
with axioms \lean{[propext, Classical.}\allowbreak\lean{choice, Quot.}\allowbreak\lean{sound]}. The library's own unconditional consumer of
this chain is \lean{minimal_}\allowbreak\lean{graded_}\allowbreak\lean{rigidity} (Chapter~\ref{ch:L07}), which takes the equivalent route
through \lean{infinite_}\allowbreak\lean{integer_}\allowbreak\lean{screening_}\allowbreak\lean{roots}.
\end{dictionary}

\begin{physics}
At a minimal-model central charge the vacuum family acquires exactly one new null vector at level
$(p-1)(q-1)$ --- the field whose decoupling defines the minimal model --- and its norm, as a function of
$c$, crosses zero linearly. That is why the central charge cannot be moved while keeping the null
relation (Chapter~\ref{ch:L07}).
\end{physics}

\begin{proof}[Proof on paper]
Set $t$ near $p/q$ and $n+1 = N$. By Theorem~\ref{l06:thm:vacblock}(2) and the Kac formula, comparing
coefficients of $X^{\lsixpp(n)}$:
$2^{\lsixpp(n)}D_N(c(t))\det G_n(c(t),1) = K_N\cdot\mathcal K'_N(t)$, where $\mathcal K'_N(t)$ is
\lean{kacVacuumScalar} (\lean{vacuumDeterminant_}\allowbreak\lean{kac_}\allowbreak\lean{coefficient}). At $t = p/q$: the label $(q-1,p-1)$ has
$h_{q-1,p-1}(p/q) = 0$ with nonzero $t$-derivative (\lean{fusionKacWeight_}\allowbreak\lean{minimal_}\allowbreak\lean{derivative}) and
exponent $\lsixpp(0) = 1$; every other label $(r,s)\ne(1,1)$ with $rs\le N$ has $h_{r,s}(p/q)\ne0$
(\lean{minimalKacWeight_}\allowbreak\lean{zero_}\allowbreak\lean{before_}\allowbreak\lean{first}, a Diophantine argument using coprimality). Hence
$\mathcal K'_N$ has a simple zero at $p/q$. Also $\det G_n(c(p/q),1)\ne0$, because $1 = h_{r,s}(p/q)$ does
not happen for $rs\le n<N$ (\lean{kacWeightProduct_}\allowbreak\lean{one_}\allowbreak\lean{minimal_}\allowbreak\lean{ne}). The map $t\mapsto c(t)$ has
$c'(p/q) = -6+6q^2/p^2\ne0$. By the chain rule (\lean{polynomial_}\allowbreak\lean{simple_}\allowbreak\lean{of_}\allowbreak\lean{curve_}\allowbreak\lean{product}),
$D_N(c_{p,q}) = 0$ and $D_N'(c_{p,q})\ne0$.
\end{proof}

\begin{status}
\lean{VOA.}\allowbreak\lean{InverseDS.}\allowbreak\lean{VirasoroVerma.}\allowbreak\lean{vacuum_}\allowbreak\lean{simple_}\allowbreak\lean{of_}\allowbreak\lean{verma_}\allowbreak\lean{kac},
\lean{…kacVacuumScalar_}\allowbreak\lean{minimal_}\allowbreak\lean{simple}: \lean{[propext, Classical.}\allowbreak\lean{choice, Quot.}\allowbreak\lean{sound]}; the first is
conditional on \lean{hk}, discharged by composition with Theorem~\ref{l06:thm:kacfull} (checked in Lean, not part of
this release; unconditional,
\lean{[propext, Classical.}\allowbreak\lean{choice, Quot.}\allowbreak\lean{sound]}).
\end{status}

%======================================================================
\section{Other files of this chapter}\label{l06:sec:others}
%======================================================================

\begingroup\raggedright
The remaining files of the chapter are infrastructure for the theorems above or for the rigidity
argument of Chapter~\ref{ch:L07}. None introduces a hypothesis or axiom; all are \lean{VOA.pbw}-free. We
group them.
\begin{itemize}
\item \emph{Normal ordering and universal pairings} (namespace \lean{VirDeformModes.}\allowbreak\lean{NormalOrdering}):
  \leanpath{VirasoroNormalOrdering}, \leanpath{VirasoroUniversalPairing}, \leanpath{VirasoroCentralContractions},
  \leanpath{VirasoroContractionDiagonal}, \leanpath{VirasoroContractionSupport}, \leanpath{VirasoroGramDegree},
  \leanpath{VirasoroGramLeadingCoefficient}, \leanpath{VirasoroGramDivisibility},
  \leanpath{VirasoroPartitionObstruction}, \leanpath{VirasoroDeterminantObstruction},
  \leanpath{VirasoroDeterminantFirstBlock}, \leanpath{VirasoroSingularGram}, \leanpath{VirasoroSingularDescendants},
  \leanpath{VirasoroSingularSpecialization}, \leanpath{VirasoroWordPairing}, \leanpath{VirasoroFixedReconstruction},
  \leanpath{VirasoroVacuumEvaluation}. A formal normal-ordering algorithm for mode words with the central charge
  as a variable; the vacuum Gram polynomial (Definition~\ref{l06:def:partgram}) and its degree and leading
  coefficient; divisibility of determinants by $(X-c)^{\text{nullity}}$; and the deformation-theoretic results
  \lean{isTrivial_}\allowbreak\lean{of_}\allowbreak\lean{zero_}\allowbreak\lean{central_}\allowbreak\lean{variation}, \lean{leeYang_}\allowbreak\lean{graded_}\allowbreak\lean{rigidity},
  \lean{graded_}\allowbreak\lean{rigidity_}\allowbreak\lean{of_}\allowbreak\lean{simple_}\allowbreak\lean{determinant}, \lean{virWord_}\allowbreak\lean{pairing_}\allowbreak\lean{obstruction}, which are stated and
  explained in Chapter~\ref{ch:L07}.
\item \emph{Verma infrastructure} (namespace \lean{InverseDS.}\allowbreak\lean{VirasoroVerma}): \leanpath{VirasoroVermaEigenspaces},
  \leanpath{VirasoroVermaUniformStraightening}, \leanpath{VirasoroVermaMomentumPolynomial},
  \leanpath{VirasoroVermaPositiveSymbols}, \leanpath{VirasoroVermaDescendantSymbols},
  \leanpath{VirasoroVermaDescendantIndependence}, \leanpath{VirasoroVermaHighestProjection},
  \leanpath{VirasoroVermaPolynomialWard}, \leanpath{VirasoroVermaWeightContractions},
  \leanpath{VirasoroVermaGramDegree}, \leanpath{VirasoroVermaTensorFock}, \leanpath{VirasoroLeadingInvariants},
  \leanpath{VirasoroTopCoordinate}, \leanpath{VirasoroVermaWardSpecialization}, \leanpath{VirasoroWardClosedKernel}
  (a singular family at nonreal parameters extends by continuity: \lean{verma_}\allowbreak\lean{highest_}\allowbreak\lean{of_}\allowbreak\lean{nonreal}),
  \leanpath{VirasoroVermaChargeTransport}, \leanpath{VirasoroVermaChargeResponse} (Ward-type identities used in the
  fusion construction of Chapter~\ref{ch:L07}).
\item \emph{Ward operators for descendant fields}: \leanpath{VirasoroDescendantWard}, \leanpath{VirasoroAllDescendantWard},
  \leanpath{VirasoroDescendantTriangular}, \leanpath{VirasoroDescendantTower}. Coefficient sequences of a
  would-be intertwining (vertex) operator between Verma modules and the defects of the $L_1$, $L_2$ covariance
  identities (the conformal Ward identities $[L_m,\phi(z)] = z^m(z\partial+(m+1)\Delta)\phi(z)$ written
  coefficientwise). They support the fusion construction of singular vectors (Chapter~\ref{ch:L07}); no fusion rule
  is assumed.
\item \emph{Screening auxiliaries}: \leanpath{VirasoroScreening\{Laurent, Vandermonde, Moment, Newton, Coefficient,
  Derivation, PowerSumDerivative, WardWeight, Integration, ResidueFunctional, Differentiation, Factorial,
  Migration, ShiftCoefficient, ModeCoefficients, PairSums\}}, \leanpath{VirasoroChargedFockGrading},
  \leanpath{VirasoroFockPositiveOperators}, \leanpath{VirasoroFockShiftSum}, \leanpath{VirasoroFockLevelDimension}:
  the finite residue calculus of Theorem~\ref{l06:thm:screensing}.
\item \emph{Vacuum-module auxiliaries}: \leanpath{VirasoroGradedField} (graded fields and the grading axiom),
  \leanpath{VirasoroGoodFiltration}, \leanpath{VirasoroFiltrationTranslation} (Li-type filtrations, used in
  Chapter~\ref{ch:L11}), \leanpath{VirasoroFockConjugation} (complex conjugation on Fock space, used for real
  structures in Chapter~\ref{ch:L19}), \leanpath{VirasoroQuotientForm} (the canonical form \lean{Bsimp} on
  $\lsixL$ and \lean{gram_simple_negWord}: its Gram entries equal those of $\lsixV$),
  \leanpath{VirasoroFormSymmetry}, \leanpath{VirasoroDescendantIndependence}, \leanpath{VirasoroVacuumFourShift},
  \leanpath{VirasoroVacuumFourContinuant} (a finite continuant recurrence at the parameters $t = 4/p$ and a
  polynomial certificate relating $p$ and $p+4$; the docstring states that the needed nonvanishing is neither
  assumed nor proved there; no headline theorem of this chapter depends on these files).
\item \emph{Kac arithmetic}: \leanpath{VirasoroKacLabelSeparation} (\lean{minimalKacWeight} and which labels have
  weight $0$ or $1$ at $t = p/q$), \leanpath{VirasoroKacSmallLabels}, \leanpath{VirasoroKacCollisionStates}
  (the partial results superseded by Theorem~\ref{l06:thm:kacfull}), \leanpath{VirasoroKacSingularTransfer}
  (also containing the conditional graded-rigidity theorems, with conclusion \lean{IsRigidGraded}, \lean{minimal_}\allowbreak\lean{rigidity_}\allowbreak\lean{of_}\allowbreak\lean{nonreal_}\allowbreak\lean{kac_}\allowbreak\lean{singular} and
  \lean{minimal_}\allowbreak\lean{graded_}\allowbreak\lean{rigidity_}\allowbreak\lean{of_}\allowbreak\lean{kac_}\allowbreak\lean{singular}), \leanpath{VirasoroKacRootBootstrap}
  (\lean{minimal_}\allowbreak\lean{rigidity_}\allowbreak\lean{of_}\allowbreak\lean{infinite_}\allowbreak\lean{kac_}\allowbreak\lean{roots}, the form consumed by Chapter~\ref{ch:L07}).
\item \emph{Combinatorics}: \leanpath{VirasoroAdmissibleCombinatorics} --- Salazar's $n$-admissible partitions
  (\lean{IsNAdmissible}: $w_j\ge w_{j+n-1}+2$) and their counts in lengths $\le2$, the combinatorial side of
  the conjectural basis of the boundary models $L(c_{2,2n+1},0)$. The file asserts no basis theorem.
\end{itemize}
\par\endgroup

%======================================================================
\section{Warnings, recorded defects and dead ends}\label{l06:sec:warnings}
%======================================================================

\begin{warning}[title={Outdated docstrings}]\label{l06:warn:docstrings}
Several docstrings, when this chapter was first written, described the state of the
library before later files landed. The Lean \emph{statements} are correct; the prose was wrong. At a pre-release build commit
most of the wording quoted below has been corrected in the sources (the ``Phase-4/5 target'' remark in the header
of \leanpath{VirasoroVOA.lean} remains); the list is kept as a record.
\begin{enumerate}
\item \emph{``The singular-existence premise remains open''.} Module or declaration docstrings of
  \leanpath{VirasoroKacProduct.lean} (\lean{vermaWeightDeterminant_}\allowbreak\lean{eq_}\allowbreak\lean{kac_}\allowbreak\lean{of_}\allowbreak\lean{singular}),
  \leanpath{VirasoroKacSingularTransfer.lean} (module header; \lean{minimal_}\allowbreak\lean{rigidity_}\allowbreak\lean{of_}\allowbreak\lean{nonreal_}\allowbreak\lean{kac_}\allowbreak\lean{singular}),
  \leanpath{VirasoroKacVacuumTransfer.lean} (\lean{vacuum_}\allowbreak\lean{simple_}\allowbreak\lean{of_}\allowbreak\lean{verma_}\allowbreak\lean{kac}, \lean{graded_}\allowbreak\lean{rigidity_}\allowbreak\lean{of_}\allowbreak\lean{verma_}\allowbreak\lean{kac}),
  \leanpath{VirasoroKacVacuumCoefficient.lean} (module header), \leanpath{VirasoroKacScalarRoot.lean}
  (\lean{kacVacuumScalar_}\allowbreak\lean{minimal_}\allowbreak\lean{simple}), \leanpath{VirasoroKacRootBootstrap.lean} and
  \leanpath{VirasoroKacLaurent.lean} (module headers), and \leanpath{VirasoroKacLabelSeparation.lean} (line 191)
  say the singular-vector or root input is missing. It is supplied by \lean{screening_}\allowbreak\lean{nonreal_}\allowbreak\lean{kac_}\allowbreak\lean{singular} and
  \lean{infinite_}\allowbreak\lean{integer_}\allowbreak\lean{screening_}\allowbreak\lean{roots} (Theorem~\ref{l06:thm:screenkac}); the composition is
  Theorem~\ref{l06:thm:kacfull}.
\item \emph{``General $c$ is the next step''.} The headers of \leanpath{VirasoroVOA.lean} and
  \leanpath{VirasoroSugawara.lean} say that a standalone Virasoro VOA for general $c$ is ``deferred''/``the
  documented next step''. It exists: \lean{VirModule.virasoroVOA} (Definition~\ref{l06:def:virVOA}) and
  \lean{VirSimple.}\allowbreak\lean{virasoroSimple} (Definition~\ref{l06:def:simple}). The same header calls the rationality of
  minimal models a ``Phase-4/5 target''; it is not proved in this chapter.
\item \emph{``$L(c,0)$'' for the universal module.} The docstrings of \lean{virasoroVOA} (``$\mathrm{Vir}(c) =
  L(c,0)$, the universal Virasoro vacuum VOA''), \leanpath{VirasoroField.lean}, \leanpath{VirasoroSpanning.lean},
  \leanpath{VirasoroWordFields.lean}, \leanpath{VirasoroGrading.lean}, \leanpath{VirasoroNonvanishing.lean} and
  \leanpath{VirasoroFock.lean} call \lean{VirUniv c} ``$L(c,0)$''. It is the universal $V(c,0)$; the simple
  $L(c,0)$ is \lean{virasoroSimple c}. \leanpath{VirasoroSimple.lean} records the intended renaming.
\end{enumerate}
Severity: documentation only.
\end{warning}

\begin{warning}[title={Redundant hypothesis $c\ne0$ (since removed from \texttt{omegaVir\_ne\_zero})}]\label{l06:warn:omega}
\textbf{Update.} Following this finding the library removed the hypothesis: \lean{omegaVir_ne_zero} and
\lean{virasoroVOA_conformalVector_ne_zero} are now stated for every $c$ (quoted in Section~\ref{l06:sec:voa}'s dictionary). The text below records the situation when the chapter
was written.
\lean{VirModule.}\allowbreak\lean{omegaVir_}\allowbreak\lean{ne_}\allowbreak\lean{zero} and \lean{VirModule.}\allowbreak\lean{virasoroVOA_}\allowbreak\lean{conformalVector_}\allowbreak\lean{ne_}\allowbreak\lean{zero}
(\leanpath{VirasoroInstance.lean}) assume \lean{hc : c ≠ 0}. In the universal module the conclusion
$\omega\ne0$ holds for every $c$, because $\omega = L_{-2}\lsixvac$ is a partition-basis vector
(\lean{negWordVir_ne_zero}). A Lean check made for this book (not part of this release) proves
\lean{omegaVir_}\allowbreak\lean{ne_}\allowbreak\lean{zero_}\allowbreak\lean{all (c : ℂ) : omegaVir c ≠ 0} in one line, axioms \lean{[propext, Classical.}\allowbreak\lean{choice,
Quot.}\allowbreak\lean{sound]}. (The hypothesis is \emph{not} redundant for the simple quotient:
\lean{virasoroSimple_}\allowbreak\lean{conformalVector_}\allowbreak\lean{ne_}\allowbreak\lean{zero} genuinely needs $c\ne0$.) Severity: weaker than possible;
harmless.
\end{warning}

\begin{warning}[title={The Kac determinant is not a library declaration}]\label{l06:warn:kacnotlib}
A referee searching the library for ``the Kac determinant formula'' finds only conditional statements
(\lean{vermaWeightDeterminant_}\allowbreak\lean{eq_}\allowbreak\lean{kac_}\allowbreak\lean{of_}\allowbreak\lean{singular}, \lean{verma_}\allowbreak\lean{kac_}\allowbreak\lean{all_}\allowbreak\lean{of_}\allowbreak\lean{nonreal_}\allowbreak\lean{singular},
\lean{vermaWeightDeterminant_}\allowbreak\lean{kac_}\allowbreak\lean{extend}) and the finite-level \lean{vermaWeightDeterminant_}\allowbreak\lean{kac_}\allowbreak\lean{below_}\allowbreak\lean{nine_}\allowbreak\lean{all}.
The unconditional formula at all levels is a two-line composition that the library never states; likewise the
unconditional vacuum simple-zero theorem. Both were checked in Lean for this book (Theorems~\ref{l06:thm:kacfull} and~\ref{l06:thm:vacsimple}); the checks are
not part of this release. Severity: conditional gap
that is closed by composition.
\end{warning}

\begin{warning}[title={Two word conventions, two free-field normalisations}]\label{l06:warn:conventions}
(1) In the vacuum module a word entry $i$ means $L_{-i-2}$ (\lean{negWordVir}, \lean{SortedOfEnergy}); in Verma
modules it means $L_{-i-1}$ (\lean{negWord}, \lean{VermaPartitions}). The map \lean{vacuumVermaPartition} adds one
to each entry. (2) The linear-dilaton realisation of Theorem~\ref{l06:thm:nonvanishing} uses
$L_n = L^{\mathrm{Sug}}_n+\lambda(n+1)a_n$ at Heisenberg level $1$ with $c = 1-12\lambda^2$; the charged
realisation of Definition~\ref{l06:def:chargedfock} uses level $2$, a minus sign and $c = 1-24\ell^2$. Both are
correct; the formulas must not be mixed.
\end{warning}

\begin{warning}[title={Dead ends and expectations that turned out false}]\label{l06:warn:deadends}
\begin{enumerate}
\item \emph{Isolated collision parameters do not suffice.} Singular vectors constructed only at special
  parameters (e.g.\ $t = \pm1$, \lean{kacSingular_at_one}) cannot feed the Laurent interpolation of
  Theorem~\ref{l06:thm:laurent}, which needs infinitely many roots; the docstring of
  \lean{minimal_}\allowbreak\lean{rigidity_}\allowbreak\lean{of_}\allowbreak\lean{infinite_}\allowbreak\lean{kac_}\allowbreak\lean{roots} records this. The integer screenings supply the infinite
  family $t\in\{1,4,9,\dots\}$.
\item \emph{Benoit--Saint-Aubin vectors give only labels with a part $\le2$.} The explicit BSA singular
  vectors (Chapter~\ref{ch:L07}) cover labels $(1,N)$, $(N,1)$ and (by composition) $(2,N)$; this gives the
  Kac formula only below level $9$ (\lean{vermaWeightDeterminant_}\allowbreak\lean{kac_}\allowbreak\lean{below_}\allowbreak\lean{nine_}\allowbreak\lean{all}), since $(3,3)$ is the
  first label with both parts $\ge3$. The general case needed the screening construction.
\item \emph{The free-field map of the vacuum module is injective at every $c$}
  (Theorem~\ref{l06:thm:ffembed}). The opposite expectation (that it ``sees'' the minimal-model null vectors as
  a kernel) is false; the kernels in the Felder complex belong to screening maps between Fock modules.
\item \emph{Determinant versus rank.} A vanishing vacuum determinant only says ``some null state''; the
  dimension of $\lsixL_n$ needs the rank (Theorem~\ref{l06:thm:dimrank}). The multiplicity of a root of
  $\det G_n$ bounds the nullity only from above (a root of order $k$ allows nullity anywhere from $1$ to $k$), so
  the library computes ranks directly, by explicit factorisations and nonzero minors, rather than reading them
  off root multiplicities.
\end{enumerate}
\end{warning}

%======================================================================
\section{Verification record}\label{l06:sec:verification}
%======================================================================

When this chapter was first written we ran \texttt{\#print axioms} on the already-built tree at an earlier commit (not re-run for this edition; no new
aggregate \texttt{lake build} was run), with a scratch file importing the relevant modules and
\begin{plaincode}
export PATH="$HOME/.elan/bin:$PATH"
lake env lean -DmaxSynthPendingDepth=3 Axioms.lean
\end{plaincode}
(the flag avoids the audit trap of Section~\ref{ch15:sec:trap}: a bare \texttt{lake env lean} drops
\texttt{maxSynthPendingDepth}). All 69 declarations checked
reported exactly
\begin{plaincode}
depends on axioms: [propext, Classical.choice, Quot.sound]
\end{plaincode}
None uses \lean{VOA.pbw} or \lean{sorryAx}. The list:
\begin{quote}\small\raggedright
\lean{VirModule.virasoroVOA}, \lean{virYA_exists}, \lean{YomegaVir_self_local}, \lean{YomegaVir_isField},
\lean{gradingVir_isInternal}, \lean{omegaVir_ne_zero}, \lean{VOA.vacVec_ne_zero}, \lean{gradedBasis},
\lean{finrank_gradingVir}, \lean{fockMap_injective}, \lean{finrank_}\allowbreak\lean{gradingVir_}\allowbreak\lean{sixteen};
\lean{VirSimple.}\allowbreak\lean{virasoroSimple_}\allowbreak\lean{isSimple}, \lean{virasoroSimple_}\allowbreak\lean{conformalVector_}\allowbreak\lean{ne_}\allowbreak\lean{zero},
\lean{virasoroVOA_}\allowbreak\lean{isSimple_}\allowbreak\lean{iff}, \lean{virasoroSimple_}\allowbreak\lean{selfDual}; \lean{VirGram4.canGram4_eq},
\lean{det_canGram4}, \lean{VirGram56.det_canGram5}, \lean{det_canGram6}, \lean{det_}\allowbreak\lean{canGram6_}\allowbreak\lean{eq_}\allowbreak\lean{zero_}\allowbreak\lean{iff},
\lean{VirFormSym.}\allowbreak\lean{canB_}\allowbreak\lean{isSymmetric}; \lean{VirSimpleGram.}\allowbreak\lean{finrank_}\allowbreak\lean{virasoroSimple_}\allowbreak\lean{grading_}\allowbreak\lean{partition},
\lean{finrank_}\allowbreak\lean{virasoroSimple_}\allowbreak\lean{leeYang_}\allowbreak\lean{four}, \lean{maxIdealVir_}\allowbreak\lean{one_}\allowbreak\lean{inf_}\allowbreak\lean{grading_}\allowbreak\lean{six_}\allowbreak\lean{eq_}\allowbreak\lean{bot};
\lean{VirSimpleSix.}\allowbreak\lean{finrank_}\allowbreak\lean{virasoroSimple_}\allowbreak\lean{six_}\allowbreak\lean{ising}, \lean{…six_leeYang}, \lean{…five_leeYang},
\lean{…six_m27}, \lean{no_}\allowbreak\lean{gradedEquiv_}\allowbreak\lean{leeYang_}\allowbreak\lean{ising}; in \lean{InverseDS.}\allowbreak\lean{VirasoroVerma}:
\lean{verma_}\allowbreak\lean{sortedWord_}\allowbreak\lean{linearIndependent}, \lean{vermaLevel_finrank}, \lean{highestVector_ne_zero},
\lean{vermaContravariant_}\allowbreak\lean{unique}, \lean{vermaContravariant_}\allowbreak\lean{symm}, \lean{vermaContravariant_}\allowbreak\lean{singular_}\allowbreak\lean{radical},
\lean{vermaSimple_}\allowbreak\lean{irreducible}, \lean{vermaSimpleLevel_}\allowbreak\lean{finrank_}\allowbreak\lean{eq_}\allowbreak\lean{gramRank}, \lean{vermaSimpleVacuumEquiv},
\lean{vacuumProjection_ker}, \lean{vermaGramPolynomial_}\allowbreak\lean{det_}\allowbreak\lean{eval}, \lean{vermaWeightDeterminant_}\allowbreak\lean{natDegree},
\lean{vermaWeightDeterminant_}\allowbreak\lean{leadingCoeff}, \lean{vermaWeightDeterminant_}\allowbreak\lean{dvd_}\allowbreak\lean{of_}\allowbreak\lean{singular},
\lean{verma_}\allowbreak\lean{partition_}\allowbreak\lean{lengths_}\allowbreak\lean{eq_}\allowbreak\lean{kac_}\allowbreak\lean{count}, \lean{kacWeightProduct_}\allowbreak\lean{natDegree},
\lean{vermaWeightDeterminant_}\allowbreak\lean{eq_}\allowbreak\lean{kac_}\allowbreak\lean{of_}\allowbreak\lean{singular}, \lean{vermaWeightDeterminant_}\allowbreak\lean{kac_}\allowbreak\lean{extend},
\lean{verma_}\allowbreak\lean{kac_}\allowbreak\lean{all_}\allowbreak\lean{of_}\allowbreak\lean{nonreal_}\allowbreak\lean{singular}, \lean{verma_}\allowbreak\lean{singular_}\allowbreak\lean{of_}\allowbreak\lean{first_}\allowbreak\lean{gram_}\allowbreak\lean{root},
\lean{verma_}\allowbreak\lean{singular_}\allowbreak\lean{of_}\allowbreak\lean{nonreal_}\allowbreak\lean{kac_}\allowbreak\lean{roots}, \lean{kacGram_}\allowbreak\lean{root_}\allowbreak\lean{of_}\allowbreak\lean{infinite},
\lean{vermaWeightDeterminant_}\allowbreak\lean{kac_}\allowbreak\lean{below_}\allowbreak\lean{nine_}\allowbreak\lean{all}, \lean{kacSingular_at_one}, \lean{verma_singular_eq_smul},
\lean{vacuum_}\allowbreak\lean{simple_}\allowbreak\lean{of_}\allowbreak\lean{verma_}\allowbreak\lean{kac}, \lean{kacVacuumScalar_}\allowbreak\lean{minimal_}\allowbreak\lean{simple}, \lean{vermaWeightDeterminant_}\allowbreak\lean{vacuum_}\allowbreak\lean{first};
in \lean{VirasoroScreening}: \lean{screeningFockState_}\allowbreak\lean{ne_}\allowbreak\lean{zero}, \lean{integer_}\allowbreak\lean{screening_}\allowbreak\lean{nonzero_}\allowbreak\lean{singular},
\lean{chargedFock_}\allowbreak\lean{singular_}\allowbreak\lean{gram_}\allowbreak\lean{root}, \lean{fock_}\allowbreak\lean{verma_}\allowbreak\lean{level_}\allowbreak\lean{finrank}, \lean{integer_}\allowbreak\lean{screening_}\allowbreak\lean{kac_}\allowbreak\lean{root},
\lean{infinite_}\allowbreak\lean{integer_}\allowbreak\lean{screening_}\allowbreak\lean{roots}, \lean{screening_kac_root}, \lean{screening_}\allowbreak\lean{nonreal_}\allowbreak\lean{kac_}\allowbreak\lean{singular};
\lean{VirDeformModes.}\allowbreak\lean{NormalOrdering.}\allowbreak\lean{partitionDeterminant_}\allowbreak\lean{natDegree}, \lean{…partitionDeterminant_}\allowbreak\lean{leadingCoeff};
\lean{VirasoroVOA.}\allowbreak\lean{virasoro_}\allowbreak\lean{relation_}\allowbreak\lean{c1}, \lean{VOA.virasoroVOA1}.
\end{quote}
The Lean checks of Theorem~\ref{l06:thm:kacfull}, of the unconditional form of Theorem~\ref{l06:thm:vacsimple} and of
Warning~\ref{l06:warn:omega} (made outside the library; not part of this release) were compiled the same way and
reported the same three axioms.

%======================================================================
\section{Referee's checklist}\label{l06:sec:checklist}
%======================================================================

{\small\begin{longtable}{>{\raggedright}p{4.1cm}>{\raggedright}p{5.2cm}>{\raggedright}p{2.6cm}>{\raggedright\arraybackslash}p{2.2cm}}
\toprule
Lean declaration & Intended meaning & Hypotheses / status & Fidelity check\\
\midrule
\endhead
\lean{VirModule.virasoroVOA} & $V(c,0)$ is a VOA of central charge $c$, for every $c$ & none; unconditional & Def.~\ref{l06:def:virVOA}\\
\lean{vacVec_ne_zero} & the vacuum of $V(c,0)$ is nonzero & none; uncond. & Thm.~\ref{l06:thm:nonvanishing}\\
\lean{gradedBasis}, \lean{finrank_gradingVir} & partition basis; $\dim V(c,0)_n = \mathrm{p}_2(n)$ & none; uncond. & Thm.~\ref{l06:thm:pbw}\\
\lean{fockMap_injective} & $V(c,0)\hookrightarrow$ Fock, all $c$ & $1-12\lambda^2 = c$; uncond. & Thm.~\ref{l06:thm:ffembed}\\
\lean{VirSimple.}\allowbreak\lean{virasoroSimple_}\allowbreak\lean{isSimple} & $L(c,0)$ is simple, all $c$ & none; uncond. & Thm.~\ref{l06:thm:simple}\\
\lean{finrank_}\allowbreak\lean{virasoroSimple_}\allowbreak\lean{grading_}\allowbreak\lean{partition} & $\dim L(c,0)_n = \operatorname{rank}G_n(c)$ & none; uncond. & Thm.~\ref{l06:thm:dimrank}\\
\lean{det_canGram4/5/6} & vacuum Kac determinants at levels 4,5,6 & none; uncond. & Thm.~\ref{l06:thm:gram456}\\
\lean{finrank_}\allowbreak\lean{virasoroSimple_}\allowbreak\lean{six_}\allowbreak\lean{ising} etc. & Lee--Yang/Ising/$M(2,7)$ dimensions & none; uncond. & Thm.~\ref{l06:thm:dims}\\
\lean{verma_}\allowbreak\lean{sortedWord_}\allowbreak\lean{linearIndependent}, \lean{vermaLevel_finrank} & PBW basis of $M(c,h)$ & none; uncond. & Thm.~\ref{l06:thm:vermabasis}\\
\lean{vacuumProjection_ker} & $V(c,0) = M(c,0)/M(c,1)$ & none; uncond. & Thm.~\ref{l06:thm:vacexact}\\
\lean{vermaContravariant_}\allowbreak\lean{unique} & Shapovalov form exists, unique & none; uncond. & Def.~\ref{l06:def:shap}\\
\lean{vermaSimple_}\allowbreak\lean{irreducible}, \lean{vermaSimpleLevel_}\allowbreak\lean{finrank_}\allowbreak\lean{eq_}\allowbreak\lean{gramRank} & $L(c,h)$ irreducible, $\dim = $ rank & none; uncond. & Thm.~\ref{l06:thm:radical}\\
\lean{vermaWeightDeterminant_}\allowbreak\lean{dvd_}\allowbreak\lean{of_}\allowbreak\lean{singular} & singular vector $\Rightarrow(X-h)^{\mathrm p(m)}\mid\det$ & singular vector; uncond. as implication & Thm.~\ref{l06:thm:singgram}\\
\lean{vermaWeightDeterminant_}\allowbreak\lean{natDegree/leadingCoeff} & degree and leading coeff.\ of $\det G_n$ & none; uncond. & Thm.~\ref{l06:thm:degree}\\
\lean{vermaWeightDeterminant_}\allowbreak\lean{eq_}\allowbreak\lean{kac_}\allowbreak\lean{of_}\allowbreak\lean{singular} & Kac formula given singular vectors & $\operatorname{Im}t\ne0$, \lean{hs}; conditional (discharged) & Thm.~\ref{l06:thm:kaccond}\\
\lean{verma_}\allowbreak\lean{kac_}\allowbreak\lean{all_}\allowbreak\lean{of_}\allowbreak\lean{nonreal_}\allowbreak\lean{singular} & Kac formula at all $t\ne0$ given singular vectors & \lean{hs}; conditional (discharged) & Thm.~\ref{l06:thm:kacextend}\\
\lean{integer_}\allowbreak\lean{screening_}\allowbreak\lean{nonzero_}\allowbreak\lean{singular} & screening states are singular & $r,m>0$; uncond. & Thm.~\ref{l06:thm:screensing}\\
\lean{chargedFock_}\allowbreak\lean{singular_}\allowbreak\lean{gram_}\allowbreak\lean{root} & Fock singular $\Rightarrow$ Gram root & Fock singular vector; uncond.\ implication & Thm.~\ref{l06:thm:transfer}\\
\lean{screening_}\allowbreak\lean{nonreal_}\allowbreak\lean{kac_}\allowbreak\lean{singular} & singular vectors at every label, nonreal $t$ & $r,s>0$; uncond. & Thm.~\ref{l06:thm:screenkac}\\
\lean{kac_determinant} (not in the library) & full Kac formula & $t\ne0$; uncond.\ (check not in this release) & Thm.~\ref{l06:thm:kacfull}\\
\lean{vacuum_}\allowbreak\lean{simple_}\allowbreak\lean{of_}\allowbreak\lean{verma_}\allowbreak\lean{kac} & simple zero of $D_N$ at $c_{p,q}$ & labels, \lean{hk}; conditional (discharged by composition) & Thm.~\ref{l06:thm:vacsimple}\\
\lean{VirasoroVOA.}\allowbreak\lean{virasoro_}\allowbreak\lean{relation_}\allowbreak\lean{c1} & abelian Sugawara at $c = 1$ & $k\ne0$; uncond. & Thm.~\ref{l06:thm:sugawara}\\
\bottomrule
\end{longtable}}

%======================================================================
\section{Notes and references}\label{l06:sec:notes}
%======================================================================

\begin{itemize}
\item A.~A.~Belavin, A.~M.~Polyakov, A.~B.~Zamolodchikov, \emph{Infinite conformal symmetry in two-dimensional
  quantum field theory}, Nucl.\ Phys.\ B \textbf{241} (1984) 333--380. Null vectors, their decoupling, minimal
  models.
\item V.~G.~Kac, \emph{Contravariant form for infinite-dimensional Lie algebras and superalgebras}, in Lecture
  Notes in Physics \textbf{94}, Springer (1979) 441--445. The determinant formula.
\item B.~L.~Feigin, D.~B.~Fuchs, \emph{Verma modules over the Virasoro algebra}, Funct.\ Anal.\ Appl.\
  \textbf{17} (1983) 241--242, and \emph{Representations of the Virasoro algebra}, in Representations of Lie
  groups and related topics, Gordon and Breach (1990). Proof of Kac's formula, structure of Verma modules,
  free-field realisations.
\item Vl.~S.~Dotsenko, V.~A.~Fateev, \emph{Conformal algebra and multipoint correlation functions in 2D
  statistical models}, Nucl.\ Phys.\ B \textbf{240} (1984) 312--348. Coulomb gas and screening charges;
  G.~Felder, \emph{BRST approach to minimal models}, Nucl.\ Phys.\ B \textbf{317} (1989) 215--236.
\item A.~Tsuchiya, Y.~Kanie, \emph{Fock space representations of the Virasoro algebra --- intertwining
  operators}, Publ.\ RIMS \textbf{22} (1986) 259--327. Singular vectors as screening integrals.
\item I.~B.~Frenkel, Y.~Zhu, \emph{Vertex operator algebras associated to representations of affine and Virasoro
  algebras}, Duke Math.\ J.\ \textbf{66} (1992) 123--168. The VOA structure on $V(c,0)$ and $L(c,0)$.
\item W.~Wang, \emph{Rationality of Virasoro vertex operator algebras}, Int.\ Math.\ Res.\ Notices (1993) no.~7,
  197--211. Rationality of $L(c_{p,q},0)$ (not formalised here).
\item M.~Gorelik, V.~G.~Kac, \emph{On simplicity of vacuum modules}, Adv.\ Math.\ \textbf{211} (2007) 621--677,
  arXiv:math-ph/0606002. The vacuum determinant and simplicity of $V(c,0)$.
\item L.~Benoit, Y.~Saint-Aubin, \emph{Degenerate conformal field theories and explicit expressions for some
  null vectors}, Phys.\ Lett.\ B \textbf{215} (1988) 517--522.
\item A.~Rocha-Caridi, \emph{Vacuum vector representations of the Virasoro algebra}, in Vertex Operators in
  Mathematics and Physics, MSRI Publ.\ \textbf{3} (1985) 451--473. Characters of minimal models.
\item P.~Di Francesco, P.~Mathieu, D.~S\'en\'echal, \emph{Conformal Field Theory}, Springer (1997), Chapters 6--9:
  the textbook account of everything in Section~\ref{l06:sec:physics}.
\item D.~Friedan, Z.~Qiu, S.~Shenker, \emph{Conformal invariance, unitarity and critical exponents in two
  dimensions}, Phys.\ Rev.\ Lett.\ \textbf{52} (1984) 1575--1578.
\item C.~Salazar, \emph{A basis for the boundary Virasoro models} (admissible monomials), the source of
  \leanpath{VirasoroAdmissibleCombinatorics}; T.~Gannon, \emph{Moonshine beyond the Monster}, Cambridge (2006),
  \S3.1.2 and \S5.2--5.3, cited in the file headers.
\end{itemize}

\endgroup
